% concatenated from main.tex, comands.tex
\documentclass[11pt,a4paper,reqno]{amsart}

\usepackage[british]{babel}
\usepackage{csquotes}
\usepackage{epsfig}
\usepackage{amsmath}
\usepackage{bbm}
\numberwithin{equation}{section}
\usepackage[hidelinks]{hyperref}

\usepackage[left=2.5cm,right=2.5cm,top=2.5cm,bottom=2.5cm,footskip=1.5cm,headsep=1cm]{geometry}
\usepackage{appendix}



\newcommand{\HRule}{\rule{\linewidth}{0.5mm}}
\newcommand{\Mod}[1]{\ (\mathrm{mod}\ #1)}
\newcommand{\interior}[1]{%
  {\kern0pt#1}^{\mathrm{o}}%
}
\DeclareMathOperator*{\esssup}{ess\,sup}



%%%%%%%%%%%%%%%%%%%%%%%%%%%%%%%%%%%%%%%%%%%%
% Global
%%%%%%%%%%%%%%%%%%%%%%%%%%%%%%%%%%%%%%%%%%%%
\usepackage[T1]{fontenc} % To switch to the T1 encoding
\usepackage{lmodern} % To switch to Latin Modern
\rmfamily % To load Latin Modern Roman and enable the following NFSS declarations.
% Declare that Latin Modern Roman (lmr) should take
% its bold (b) and bold extended (bx) weight, and small capital (sc) shape, 
% from the corresponding Computer Modern Roman (cmr) font, for the T1 font encoding.
\DeclareFontShape{T1}{lmr}{b}{sc}{<->ssub*cmr/bx/sc}{}
\DeclareFontShape{T1}{lmr}{bx}{sc}{<->ssub*cmr/bx/sc}{}
\usepackage{lipsum}
\usepackage{mathtools}
\usepackage{amsmath}
\usepackage{amssymb}
\usepackage{tikz}
\usetikzlibrary{matrix,positioning,decorations.pathreplacing,calc}
\usetikzlibrary{
decorations.text,%
decorations.markings,%
shadows}
% For tikz with pagewidth
\usepackage{adjustbox}
\usepackage{graphicx} 

\usepackage{caption}
\usepackage{subcaption}

\usepackage{dsfont} % for 1 bbmath
\usepackage[format=hang]{caption}

\usepackage[normalem]{ulem}


%    Absolute value notation
\newcommand{\abs}[1]{\lvert#1\rvert}
\usepackage{bm}

%%%%%%%%%%%%%%%%%%%%%%%%%%%%%%%%%%%%%%%%%%%%%%%%%%%%%%%%%%%%
% Biblio
\usepackage[
sorting=nty,
maxnames=99,
maxalphanames=5,
natbib=true,
sortcites]{biblatex}

\addbibresource{bibliography.bib}
\renewcommand{\mkbibnamefamily}[1]{\textsc{#1}}
\DeclareNameAlias{default}{family-given}
\renewcommand*{\newunitpunct}{\addcomma\space}
\renewcommand*{\revsdnamepunct}{}

\graphicspath{{figures/}}

\AtEveryBibitem{% Clean up the bibtex rather than editing it
 \clearfield{url}
 \clearfield{issn}
 \clearfield{isbn}
 \clearfield{urldate}
 
 \ifentrytype{book}{
  \clearfield{pages}}{% 
 }
}


%%%%%%%%%%%%%%%%%%%%%%%%%%%%%%%%%%%%%%%%%%%%%%%%%%%%%%%%%%%%
% TODOs
\usepackage{xargs}                      % Use more than one optional parameter in a new commands
%\usepackage[pdftex,dvipsnames]{xcolor}  % Coloured text etc.
% 
\usepackage[prependcaption,textsize=tiny,textwidth=3cm]{todonotes}
\newcommandx{\unsure}[2][1=]{\todo[linecolor=red,backgroundcolor=red!25,bordercolor=red,#1]{#2}}
\newcommandx{\change}[2][1=]{\todo[linecolor=blue,backgroundcolor=blue!25,bordercolor=blue,#1]{#2}}
\newcommandx{\info}[2][1=]{\todo[linecolor=OliveGreen,backgroundcolor=OliveGreen!25,bordercolor=OliveGreen,#1]{#2}}
\newcommandx{\improvement}[2][1=]{\todo[linecolor=black,backgroundcolor=black!25,bordercolor=black,#1]{#2}}
\newcommandx{\thiswillnotshow}[2][1=]{\todo[disable,#1]{#2}}
%%%%%%%%%%%%%%%%%%%%%%%%%%%%%%%%%%%%%%%%%%%%%%%%%%%%%%%%%%%%
\usepackage{amsthm}
\usepackage[noabbrev,capitalize]{cleveref}
\crefname{proposition}{Proposition}{Propositions}
\crefname{equation}{}{}


\newtheorem{theorem}{Theorem}[section]
\newtheorem{lemma}[theorem]{Lemma}
\newtheorem{proposition}[theorem]{Proposition}
\newtheorem{corollary}[theorem]{Corollary}
\newtheorem{conjecture}{Conjecture}
\theoremstyle{definition}
\newtheorem{definition}[theorem]{Definition}
\newtheorem{example}[theorem]{Example}
\newtheorem{assumption}[theorem]{Assumption}
\newtheorem{remark}[theorem]{Remark}


\crefname{assumption}{Assumption}{Assumptions}
\crefname{definition}{Definition}{Definitions}
\crefname{corollary}{Corollary}{Corollaries}
\crefname{enumi}{item}{items}
%%%%%%%%%%%%%%%%%%%%%%%%%%%%%%%%%%%%%%%%%%%%%%%%%%%%%%%%%%%%%%%%%%%%%%



\usepackage{fancyhdr}

\pagestyle{fancy}
\fancyhf{}
\fancyhead[CE]{{\footnotesize\textsc{\nouppercase\rightmark}}}
\fancyhead[CO]{{\footnotesize\textsc{\nouppercase\leftmark}}}
\fancyfoot[CE,CO]{\footnotesize\thepage}
\fancyfoot[RO,LE]{}
\renewcommand{\headrulewidth}{0pt}
\renewcommand{\footrulewidth}{0pt}
%
\fancypagestyle{plain}{%
\fancyhead[C]{}
\fancyfoot[C]{}
\renewcommand{\headrulewidth}{0pt}
\renewcommand{\footrulewidth}{0pt}}
%
\fancypagestyle{nosection}{%
\fancyhead[CE]{}
\fancyhead[CO]{}
\fancyfoot[CE,CO]{\thepage}
%
}
%%%%%%%%%%%%%%%%%%%%%%%%%%%%%%%%%%%%%%%%%%%%%%%%%%%%%%%%%%%%%%%%%%%%%%%%%
%%%%%%%%%%%%%%%%%%%%%%%%%%%%%%%%%%%%%%%%%%%%
% Local
%%%%%%%%%%%%%%%%%%%%%%%%%%%%%%%%%%%%%%%%%%%%
\newcommand{\nm}{\noalign{\smallskip}}
\newcommand{\ds}{\displaystyle}

\DeclareMathOperator{\N}{\mathbb{N}}
\DeclareMathOperator{\Z}{\mathbb{Z}}
\DeclareMathOperator{\Q}{\mathbb{Q}}
\DeclareMathOperator{\R}{\mathbb{R}}
\DeclareMathOperator{\C}{\mathbb{C}}
\DeclareMathOperator{\sgn}{sign}
\DeclareMathOperator{\tr}{tr}
\renewcommand{\i}{\mathbf{i}}
\renewcommand{\tilde}{\widetilde}
\renewcommand{\hat}{\widehat}
\renewcommand{\bar}[1]{\overline{#1}}
\newcommand{\inv}{^{-1}}
\newcommand{\<}{\ensuremath \leq}
\renewcommand{\>}{\ensuremath \geq}
\newcommand{\If}{\textrm{ if }}
\renewcommand{\qedsymbol}{$\blacksquare$}


\usepackage{fancyhdr}

\pagestyle{fancy}
\fancyhf{}
\fancyhead[CE]{\nouppercase{\footnotesize\textsc{\rightmark}}}
\fancyhead[CO]{\footnotesize\textsc{\nouppercase\leftmark}}
\fancyfoot[CE,CO]{\footnotesize\thepage}
\fancyfoot[RO,LE]{}
\renewcommand{\headrulewidth}{0pt}
\renewcommand{\footrulewidth}{0pt}
%
\fancypagestyle{plain}{%
\fancyhead[C]{}
\fancyfoot[C]{}
\renewcommand{\headrulewidth}{0pt}
\renewcommand{\footrulewidth}{0pt}}
%
\fancypagestyle{nosection}{%
\fancyhead[CE]{}
\fancyhead[CO]{}
\fancyfoot[CE,CO]{\thepage}
%
}
%%%%%%%%%%%%%%%%%%%%%%%%%%%%%%%%%%%%%%%%%%%%%%%%%%%%%%%%%%%%%%%%%%%%%%%%%

\DeclareMathOperator{\slp}{\mathcal{S}}
\DeclareMathOperator{\dtn}{\mathcal{T}}
\DeclareMathOperator{\diag}{diag}
\DeclareMathOperator{\BO}{\mathcal{O}}
\DeclareMathOperator{\exactcapmat}{\mathbf{C}}
\DeclareMathOperator{\capmat}{\mathcal{C}}
\DeclareMathOperator{\qpcapmat}{\mathcal{C}^\alpha}
\DeclareMathOperator{\capmatg}{\mathcal{C}^\gamma}
\DeclareMathOperator{\qpcapmatg}{\mathcal{C}^{\gamma,\alpha}}
\newcommand{\bs}{\backslash}
\newcommand{\pri}{^\prime}
\newcommand{\prii}{^{\prime\prime}}
\renewcommand{\epsilon}{\varepsilon}
\DeclareMathOperator{\dd}{d\!}
\let\oldemptyset\emptyset
\let\emptyset\varnothing
\newcommand{\zp}[1]{\phi^{\text{zak}}_{#1}}
\renewcommand{\i}{\mathbf{i}}
\renewcommand{\tilde}{\widetilde}
\renewcommand{\hat}{\widehat}
\newcommand{\zakn}{\varphi^\text{zak}}
\newcommand{\zak}[1]{\varphi^\text{zak}_{#1}}
\DeclareMathOperator{\Ind}{Ind}
\renewcommand{\sp}{\mathrm{sp}}



\DeclareMathOperator{\iL}{{\mathsf{L}}}
\DeclareMathOperator{\iR}{{\mathsf{R}}}
\DeclareMathOperator{\iLR}{{\mathsf{L},\mathsf{R}}}

\usepackage{tcolorbox}
\usetikzlibrary{tikzmark,calc}


%Trigonometric functions
\DeclareMathOperator{\sech}{sech}
\DeclareMathOperator{\csch}{csch}
\DeclareMathOperator{\arcsec}{arcsec}
\DeclareMathOperator{\arccot}{arcCot}
\DeclareMathOperator{\arccsc}{arcCsc}
\DeclareMathOperator{\arccosh}{arcCosh}
\DeclareMathOperator{\arcsinh}{arcsinh}
\DeclareMathOperator{\arctanh}{arctanh}
\DeclareMathOperator{\arcsech}{arcsech}
\DeclareMathOperator{\arccsch}{arcCsch}
\DeclareMathOperator{\arccoth}{arcCoth} 
\DeclareMathOperator*{\argmin}{argmin}


\usepackage{enumerate}
\usepackage{upgreek}

\begin{document}




\title[]{Regularized Brascamp-Lieb inequalities via Optimal Transport and Study of Equality Cases}
\author[B. Ammari]{Bader Ammari}
\address[]{Department of Mathematics, ETH Z\"{u}rich, R\"{a}mistrasse 101, CH-8092 Z\"{u}rich, Switzerland}
\email{bader.ammari@math.ethz.ch}

\large
\maketitle



\begin{abstract}
    We consider regularized Brascamp-Lieb inequalities using the theory of optimal transportation, more precisely an anisotropic version of Caffarelli's contraction theorem. Furthermore, we provide a full picture concerning the issues of finiteness of the Brascamp-Lieb constant and of the existence of Gaussian extremizers. We also find all optimizers for these regularized Brascamp-Lieb inequalities by employing heat flow methods that were already used to settle this question for the non-regularized Brascamp-Lieb inequality and introducing new ideas to deal with several difficulties, which do not appear for the non-regularized Brascamp-Lieb datum. Finally, we give some interesting applications.
\end{abstract}

\hfill \\


\vspace{0.2cm}
{\small
\noindent{\textbf{2020 Mathematics Subject Classification:} 26D15, 52A40, 49Q20, 35K05.} }


\vspace{0.2cm}
{\small
\noindent{\textbf{Keywords:} Brascamp-Lieb inequality, Log-convexity/concavity, Gaussian extremizer, Caffarelli's contraction, Heat flow.}}
\vspace{0.5cm}



\section{Introduction}



Forward Brascamp-Lieb inequalities provide a natural generalization to many important inequalities as the inequalities of Hölder, Young convolution or Loomis-Whitney and are nowadays very well-understood. In addition, applications given by Brascamp-Lieb (BL) inequalities may be found in harmonic analysis (see \cite{BCT} in the context of Fourier restriction theory), analytic number theory (see \cite{Guo_Zhang}), convex geometry (see \cite{KBall,KBall2} where the author studies estimates for the volume of sections of the unit ball of $(\mathbb{R}^m,\|\cdot\|_p)$ as well as the volume ratio), probability, and much more. 
\hfill \\ \\ The first crucial result that made the study of Brascamp-Lieb inequalities flourish is the following observation of Gaussian saturation by Lieb in \cite{Lieb}: 

\begin{theorem}[Lieb \cite{Lieb}]
    Let $\mathbb{R}^{n \times n} \ni \mathcal{Q} \geq 0$, $C_i: \mathbb{R}^n \rightarrow \mathbb{R}^{n_i}$ surjective linear transformations for $i=1,\cdots,m$ and $p_1,\cdots,p_m>0.$ Then, the largest constant $C$ such that $$\int_{\mathbb{R}^n}e^{-\pi \langle \mathcal{Q}x,x\rangle}\prod_{i=1}^mf_i(C_ix)^{p_i}\, dx \leq C\prod_{i=1}^m \Bigl( \int_{\mathbb{R}^{n_i}}f_i \Bigl)^{p_i}$$ holds for all non-negative functions $f_1,\cdots,f_m$ is also the largest constant such that the inequality holds for centered Gaussian functions and is usually called the Brascamp-Lieb constant.
\end{theorem}



From this work, two main approaches have emerged: \hfill \\

The first is based on heat flow methods and was first introduced in \cite{Carlen}. Using a monotonicity formula for positive solutions to heat equations in multilinear settings, which was inspired by \cite{Carlen}, \cite{BCCT} established the necessary conditions for the Brascamp-Lieb constant to be finite or for its value to be attained. In addition, Valdimarsson in \cite{Vald2} settled the question of the form of functions that give equality when it is known that the Brascamp-Lieb constant is attained.

\hfill \\

The second method has its roots in optimal transport. The pioneering papers of Brenier and McCann (see \cite{Brenier,MCann}) led to the well-known Brenier's theorem which enabled identifying the optimal transport map between absolutely continuous measures in $\mathbb{R}^n$ as the gradient of a convex function. These works also hinted at the connection between optimal transportation theory and functional/geometric inequalities. In Euclidean spaces, many classical inequalities can then be derived or redemonstrated using Brenier maps combined with the change of variables formula for convex gradients. The forward Brascamp-Lieb inequality was one of the first inequalities to which this strategy was applied in \cite{Barthe} and, surprisingly, one even gets a dual inequality which is now known as the reverse Brascamp-Lieb inequality. 

\hfill \\ Our goal in this work is to give a complete study of a regularized version of the forward Brascamp-Lieb inequalities in the sense that a log-convexity/log-concavity condition will be imposed on the inputs. Regularizing functional inequalities by restricting the inputs to a smaller subclass of functions is a natural procedure whose main advantage consists in obtaining an improved form of the inequality. This viewpoint is often adopted in areas such as convex geometry, differential geometry, or information theory. Among them, let us mention three examples (for more references, see \cite{Nakamura}). In \cite{Fathi}, the constant of the log-Sobolev inequality is improved thanks to a regularization in terms of the Poincaré constant. In \cite{Nakamura_Tsuji}, deficit estimates for Nelson's hypercontractivity, the logarithmic Sobolev inequality and Talagrand's cost inequality are given under the assumptions of log-subharmonicity, log-convexity and log-concavity. Notice that the hypercontractivity is proven via Fokker-Planck flows which behave well with respect to log-subharmonicity while Talagrand's cost inequality is proven by an optimal transport argument inspired by \cite{Cordero}. 
Finally in the works of \cite{Shohei1,Shohei2} (see also \cite{Cordero_Fradelizi,Milman,Shohei3}), the authors manage to first prove an inverse Brascamp-Lieb inequality under centering, log-concavity and log-convexity assumptions. At the limit, it then retrieves the inverse Brascamp-Lieb inequality of \cite{Barthe_Wolff} without non-degeneracy conditions, for arbitrary centered log-concave inputs, which allows proving the Gaussian correlation inequality for centered convex sets. In addition, we refer the reader to \cite{Strong_Gaussian_Correlation} for the beautiful resolution of the strong Gaussian correlation conjecture by the Forward-Reverse Brascamp-Lieb inequality for centered log-concave functions. 

\hfill \\


Let us explain the main results in this work. We first manage to prove regularized Brascamp-Lieb inequalities and their dual (see \cref{Theorem_REG_BL}) using an anisotropic extension of Caffarelli's contraction theorem and a strategy similar to the proof in \cite{Barthe} of the classical forward Brascamp-Lieb inequality. 
This generalizes the works of \cite{Vald1} and \cite{BCCT} where type $G$ functions, \textit{i.e.} convolutions with $N(0,G^{-1})$, are considered. 
Then, we find necessary and sufficient conditions for the finiteness of a regularized Brascamp-Lieb constant (see \cref{Necessary_suff_cond_finite_BL}) and show that extremizable Brascamp-Lieb data are also Gaussian extremizable (see \cref{Ext_implies_Gaussian_ext}). Finally, using the heat flow approach employed in \cite{BCCT}, we define a notion of generalized geometric Brascamp-Lieb data to which any extremizable Brascamp-Lieb datum is equivalent (see \cref{Prop_Equiv_BL_constant}) and prove that extremizers of such data have a very specific form (see \cref{Main_Thm_Eq_Geometric_BL,Main_Thm_Eq_BL}).


\hfill \\
To the best of our knowledge, this is one of the first works dealing with the structure and equality cases of a functional inequality where the inputs are supposed to satisfy a log-convexity/log-concavity condition. In particular, new ideas are needed, especially for the equality cases of generalized geometric Brascamp-Lieb datum where several difficulties, which do not appear for the usual geometric Brascamp-Lieb datum, arise. 



\hfill \\ \\ This work is organized as follows. In the first section, we prove regularized versions of Brascamp-Lieb inequalities along with their dual reverse inequalities using arguments from the optimal transportation theory. In the second section, we focus on establishing finiteness conditions on the Brascamp-Lieb constant and a satisfactory description of the Brascamp-Lieb datum for which extremizers or Gaussian extremizers exist, thus generalizing the main results of \cite{BCCT}. The third section is devoted to the study of equality cases for extremizable Brascamp-Lieb data and generalizes the work of \cite{Vald2}. Finally, some applications are provided in the last section and a general Caffarelli contraction is proven in the appendix, where an extension of the results of \cite{Guido} is also included, as it is an important tool for the proof of the regularized Brascamp-Lieb inequalities. 




\bigskip


\section{Proof of Regularized Brascamp-Lieb Inequalities}

We will consider regularized Brascamp-Lieb inequalities on Euclidean spaces, \textit{i.e.} a finite-dimensional real Hilbert space $H$ endowed with the usual Lebesgue measure $dx.$ 
All functions below will be assumed to be non-negative. 


\hfill \\


Given a positive definite transformation $A: H \rightarrow H$, we denote $$g_A(x)=\exp(-\pi \langle Ax,x \rangle)$$ for $x \in H.$ We say $0<A \leq B$ for $A,B: H \rightarrow H$ two positive definite transformations iff $B-A$ is positive semi-definite. 

Moreover, let $C_i: H \rightarrow H_i$ surjective linear transformations for $i=1,\cdots,m$ and $p_1,\cdots,p_m>0.$ 



 

\begin{definition}
A function $f: H \rightarrow [0;+\infty)$ is log-concave (resp. log-convex) if $\log(f): H \rightarrow \mathbb{R} \cup \{-\infty\}$ is concave (resp. convex). Given a positive semi-definite $Q: H \rightarrow H$, we say that $f$ is more log-concave (resp. more log-convex) than $g_Q$ if $f=g_Qh$ for some log-concave (resp. log-convex) function $h.$ 
\end{definition}


\begin{remark}
For a smooth function $f$, $f$ is more log-concave (resp. more log-convex) than $g_Q$ iff $$\mathrm{Hess}(-\log f) \geq 2\pi Q \; (\text{resp. } \leq).$$
\end{remark}

\hfill \\



\begin{theorem} \label{Reformulation_REG_BL} 
Let $\mathcal{Q}: H \rightarrow H$ be a positive semi-definite transformation, $I \subset \{1,\cdots,m\}, Q_i: H_i \rightarrow H_i$ be positive definite transformations for any $1 \leq i \leq m.$ If $f_i \in L^1(H_i)$ is more log-convex than $g_{Q_i}$ for any $i \in I$ and $f_j \in L^1(H_j)$ is more log-concave than $g_{Q_j}$ for any $j \in I^c$, then $$\begin{array}{lll} \ds \frac{\ds \int_H e^{-\pi \langle \mathcal{Q}x,x\rangle}\prod_{i=1}^m f_i(C_ix)^{p_i}\, dx}{\ds \prod_{i=1}^m \Bigl(\int_{H_i}f_i \Bigl)^{p_i}} \\ \leq \ds \sup_{\substack{0<B_i \leq Q_i \text{ for all } i \in I \\ B_j \geq Q_j \text{ for all } j \in I^c}}\frac{\ds \int_H e^{-\pi \langle \mathcal{Q}x,x\rangle}\prod_{i=1}^m g_{B_i}(C_ix)^{p_i}\, dx}{\ds \prod_{i=1}^m \Bigl(\int_{H_i}g_{B_i} \Bigl)^{p_i}}=\sup_{\substack{0<B_i \leq Q_i \text{ for all } i \in I \\ B_j \geq Q_j \text{ for all } j \in I^c}}\Bigl \{ \frac{\prod_{i=1}^m (\det B_i)^{p_i}}{\det(\mathcal{Q}+\sum_{i=1}^m p_iC_i^* B_iC_i)}\Bigl \}^{1/2}.
    \end{array}$$
By duality, if $g_i$ is more log-concave than $g_{Q_i^{-1}}$ for any $i \in I$ and $g_j$ is more log-convex than $g_{Q_j^{-1}}$ for any $j \in I^c$, then  $$\begin{array}{lll}
    \ds \frac{\ds \int_{H}g_{\mathcal{Q}^{-1}} * \sup \Bigl \{ \prod_{j=1}^m g_j(y_j)^{p_j}: \sum_{j=1}^m p_jC_j^*y_j=x, y_j \in H_j\Bigl \}\, dx}{\ds \prod_{i=1}^m \Bigl( \int_{H_i}g_i\Bigl)^{p_i}} \\ \nm 
    \geq \ds \inf_{\substack{B_i \geq Q_i^{-1} \text{ for all } i \in I \\ 0<B_j \leq Q_j^{-1} \text{ for all } j \in I^c}}\frac{\ds \int_{H}g_{\mathcal{Q}^{-1}} *\sup \Bigl \{ \prod_{j=1}^m g_{B_j}(y_j)^{p_j}: \sum_{j=1}^m p_jC_j^*y_j=x, y_j \in H_j\Bigl \}\, dx}{\ds \prod_{i=1}^m \Bigl( \int_{H_i}g_{B_i}\Bigl)^{p_i}} \\ \nm
    \qquad \ds =\inf_{\substack{0< B_i \leq Q_i \text{ for all } i \in I \\ B_j \geq Q_j \text{ for all } j \in I^c}} \Bigl \{ \frac{\det(\mathcal{Q}+\sum_{i=1}^m p_i C_i^*B_iC_i)}{\prod_{i=1}^m (\det B_i)^{p_i}}  \Bigl \}^{1/2} \end{array}$$ where $$f * g(x)=\sup_{y \in H}f(y)g(x-y)$$ and $g_{\mathcal Q^{-1}}$ is understood for $\mathcal{Q}=0$ to be equal to $$\left\{ \begin{array}{l}
    \ds 1 \text{ if } y=0,\\
    \nm
    0 \text{ if } y\neq 0.
    \end{array}
    \right. $$
    Even more generally, for $\tilde G_i,G_i:H_i \rightarrow H_i$ with $0 \le \tilde G_i \le G_i \le +\infty$ for any $1 \le i \le m$, if $f_i \in L^1(H_i)$ is more log-concave than $g_{\tilde G_i}$ and more log-convex than $g_{G_i}$, then $$\begin{array}{lll} \ds \frac{\ds \int_H e^{-\pi \langle \mathcal{Q}x,x\rangle}\prod_{i=1}^m f_i(C_ix)^{p_i}\, dx}{\ds \prod_{i=1}^m \Bigl(\int_{H_i}f_i \Bigl)^{p_i}} \\ \leq \ds \sup_{\tilde{G_i} \le B_i \le G_i}\frac{\ds \int_H e^{-\pi \langle \mathcal{Q}x,x\rangle}\prod_{i=1}^m g_{B_i}(C_ix)^{p_i}\, dx}{\ds \prod_{i=1}^m \Bigl(\int_{H_i}g_{B_i} \Bigl)^{p_i}}=\sup_{\tilde{G_i} \le B_i \le G_i}\Bigl \{ \frac{\prod_{i=1}^m (\det B_i)^{p_i}}{\det(\mathcal{Q}+\sum_{i=1}^m p_iC_i^* B_iC_i)}\Bigl \}^{1/2}\end{array}$$ and similarly, for the dual inequality. 
\end{theorem}

\hfill \\

Considering $I=\emptyset$ and $I=\{1,\cdots,m\}$ gives the following proposition.

\begin{proposition} \label{Theorem_REG_BL}
    Let $\mathcal{Q}: H \rightarrow H$ be a positive semi-definite transformation and $Q_i: H_i \rightarrow H_i$ be positive definite transformations for $i=1,\cdots,m$. If $f_i \in L^1(H_i)$ is more log-convex than $g_{Q_i}$ for any $1 \leq i \leq m$, then \begin{equation} \label{Reg_BL_Log_Convexity} \begin{array}{lll} \ds \frac{\ds \int_H e^{-\pi \langle \mathcal{Q}x,x\rangle}\prod_{i=1}^m f_i(C_ix)^{p_i}\, dx}{\ds \prod_{i=1}^m \Bigl(\int_{H_i}f_i \Bigl)^{p_i}} \\ \leq \ds \sup_{0<B_i \leq Q_i}\frac{\ds \int_H e^{-\pi \langle \mathcal{Q}x,x\rangle}\prod_{i=1}^m g_{B_i}(C_ix)^{p_i}\, dx}{\ds \prod_{i=1}^m \Bigl(\int_{H_i}g_{B_i} \Bigl)^{p_i}}=\sup_{0<B_i \leq Q_i}\Bigl \{ \frac{\prod_{i=1}^m (\det B_i)^{p_i}}{\det(\mathcal{Q}+\sum_{i=1}^m p_iC_i^* B_iC_i)}\Bigl \}^{1/2}.
    \end{array}\end{equation} By duality, if $g_j$ is more log-concave than $g_{Q_j^{-1}}$ for any $1 \leq j \leq m$, then \begin{equation} \label{Reg_Dual_BL_Log_Concavity} \begin{array}{l}
    \ds \frac{\ds \int_{H}g_{\mathcal{Q}^{-1}} * \sup \Bigl \{ \prod_{j=1}^m g_j(y_j)^{p_j}: \sum_{j=1}^m p_jC_j^*y_j=x, y_j \in H_j\Bigl \}\, dx}{\ds \prod_{i=1}^m \Bigl( \int_{H_i}g_i\Bigl)^{p_i}} \\ 
    \nm 
    \geq \ds \inf_{B_i \geq Q_i^{-1}}\frac{\ds \int_{H}g_{\mathcal{Q}^{-1}} * \sup \Bigl \{ \prod_{j=1}^m g_{B_j}(y_j)^{p_j}: \sum_{j=1}^m p_jC_j^*y_j=x, y_j \in H_j\Bigl \}\, dx}{\ds \prod_{i=1}^m \Bigl( \int_{H_i}g_{B_j}\Bigl)^{p_i}}  \\
    \nm 
    \qquad = \ds \inf_{0<B_i \leq Q_i} \Bigl \{ \frac{\det(\sum_{i=1}^m p_i C_i^*B_iC_i+\mathcal{Q})}{\prod_{i=1}^m (\det B_i)^{p_i}}  \Bigl \}^{1/2}.\end{array} \end{equation} 

    
   
    If $f_i \in L^1(H_i)$ is more log-concave than $g_{Q_i}$ for any $1 \leq i \leq m$, then
    \begin{equation} \label{Reg_BL_Log_Concavity} \begin{array}{lll} 
     \frac{\ds \int_H e^{-\pi \langle \mathcal{Q}x,x\rangle} \prod_{i=1}^m f_i(C_ix)^{p_i}\, dx }{\ds \prod_{i=1}^m \Bigl(\int_{H_i}f_i\Bigl)^{p_i}} 
    \\ \nm
    \leq  \ds \sup_{B_i \geq Q_i}\frac{\ds \int e^{-\pi \langle \mathcal{Q}x,x\rangle}\prod_{i=1}^m g_{B_i}(C_ix)^{p_i}\, dx}{\ds \prod_{i=1}^m \Bigl(\int_{H_i}g_{B_i} \Bigl)^{p_i}}=\sup_{B_i \geq Q_i}\Bigl \{ \frac{\prod_{i=1}^m (\det B_i)^{p_i}}{\det(\mathcal{Q}+\sum_{i=1}^m p_iC_i^* B_iC_i)}\Bigl \}^{1/2}. \end{array} \end{equation} By duality, if $g_j$ is more log-convex than $g_{Q_j^{-1}}$ for any $1 \leq j \leq m$, then \begin{equation} \label{Reg_Dual_BL_Log_Convexity} \begin{array}{lll}
    \ds \frac{\ds \int_{H}g_{\mathcal{Q}^{-1}} * \sup \Bigl \{ \prod_{j=1}^m g_j(y_j)^{p_j}: \sum_{j=1}^m p_jC_j^*y_j=x, y_j \in H_j\Bigl \}\, dx}{\ds \prod_{i=1}^m \Bigl( \int_{H_i}g_i\Bigl)^{p_i}} \\ \nm 
    \geq \ds \inf_{0<B_i \leq Q_i^{-1}}\frac{\ds \int_{H}g_{\mathcal{Q}^{-1}} *\sup \Bigl \{ \prod_{j=1}^m g_{B_j}(y_j)^{p_j}: \sum_{j=1}^m p_jC_j^*y_j=x, y_j \in H_j\Bigl \}\, dx}{\ds \prod_{i=1}^m \Bigl( \int_{H_i}g_{B_i}\Bigl)^{p_i}} \\ \nm
    \qquad \ds =\inf_{B_i \geq Q_i} \Bigl \{ \frac{\det(\mathcal{Q}+\sum_{i=1}^m p_i C_i^*B_iC_i)}{\prod_{i=1}^m (\det B_i)^{p_i}}  \Bigl \}^{1/2}. \end{array} \end{equation} 
\end{proposition}

\hfill \\

We should mention that \cref{Reg_BL_Log_Convexity}, in the special case $$H=\mathbb{R}^N=\bigoplus_{i=1}^m \mathbb{R}^{n_i}, H_i=\mathbb{R}^{n_i} \text{ and } C_i=P_{H_i}$$ for any $1 \le i \le m$, has already appeared in the work of \cite{Milman} and \cref{Reg_Dual_BL_Log_Concavity}, in the case $\mathcal Q=0$, was proven in \cite{Vald1}. 


%For the sake of clarity, we will only study, starting from Section 3, the structure of \cref{Reg_BL_Log_Convexity}. The corresponding theorems treating \cref{Reg_BL_Log_Concavity} have flipped signs compared to the theorems established for \cref{Reg_BL_Log_Convexity}. Regarding \cref{Reformulation_REG_BL}, when conditions of log-convexity and log-concavity are possibly mixed, the technique of proofs and the statements, apart from the finiteness of the Brascamp-Lieb constant, remain the same.    


\hfill \\

When $\mathcal{Q}=0$, let us note that if the non-degeneracy assumption $$\bigcap_{i=1}^m \ker(C_i)=\{0\}$$ is not satisfied, the expressions in \cref{Reg_BL_Log_Concavity,Reg_BL_Log_Convexity} are equal to $+\infty$ and those in \cref{Reg_Dual_BL_Log_Concavity,Reg_Dual_BL_Log_Convexity} are equal to $0.$ Thus, from now on, if $\mathcal{Q}$ is assumed to be equal to $0$, the non-degeneracy assumption is implicitly supposed to be true.

\begin{remark}
Recall that $$\int e^{-\pi \langle Ax,x\rangle}\, dx=\frac{1}{\sqrt{\det A}}$$ for a positive definite $n \times n$ matrix $A.$
This shows that the equality for the suprema both in \cref{Reg_BL_Log_Convexity,Reg_BL_Log_Concavity} holds. We will see in the proof why the equality for infima in both \cref{Reg_Dual_BL_Log_Concavity,Reg_Dual_BL_Log_Convexity} holds. 
Moreover, taking $Q_i=\Lambda \mathrm{id}_{H_i}$ with $\Lambda \uparrow +\infty$ in \cref{Reg_BL_Log_Convexity}, one recovers the usual Brascamp-Lieb inequality (see \cite[Corollary 8.16]{BCCT}), and \cref{Reg_Dual_BL_Log_Convexity} gives back the reverse Brascamp-Lieb formula of \cite{Barthe} by taking $Q_j=\Lambda \mathrm{id}_{H_j}$ and letting $\Lambda \rightarrow +\infty$. 
\end{remark}









\begin{remark}
To prove \cref{Reformulation_REG_BL}, one may reduce to the case $\mathcal{Q}=0$. 
Indeed, first suppose that $\mathcal{Q}>0.$ We take $f_{m+1}=g_{\mathcal Q}, p_{m+1}=1$ and $H_{m+1}=H, C_{m+1}=\mathrm{id}_{H}, Q_{m+1}=\mathcal Q$ (which corresponds to a localised BL data in the terms of \cite{BCCT}). Then, assuming that the forward regularized Brascamp-Lieb inequality with $\mathcal Q=0$ holds true and setting $\ds M=\sum_{i=1}^m p_iC_i^*B_iC_i, I \subset \{1,\cdots,m\}$, one has $$\begin{array}{lll} \ds \frac{\ds \int_H e^{-\pi \langle \mathcal Qx,x\rangle}\prod_{i=1}^m f_i(C_ix)^{p_i}\, dx}{\ds \prod_{i=1}^m \Bigl(\int_{H_i}f_i(x_i) \, dx_i \Bigl)^{p_i}}
\leq \ds \Bigl(\int_{H} f_{m+1}(x) \, dx \Bigl) \sup_{\substack{0<B_i \leq Q_i \text{ for all } i \in I \cup \{m+1\} \\  B_{j} \geq Q_j \text{ for all } j \in I^c}}\frac{\ds \int \prod_{i=1}^{m+1} g_{B_i}(C_ix)^{p_i}\, dx}{\ds \prod_{i=1}^{m+1} \Bigl(\int_{\mathbb{R}^{n_i}}g_{B_i}(x_i) \, dx_i \Bigl)^{p_i}} \\
\nm
\leq \ds \frac{1}{\sqrt{\det\mathcal Q}} \sup_{\substack{0<B_i \leq Q_i \text{ for all } i \in I \cup \{m+1\} \\  B_{j} \geq Q_j \text{ for all } j \in I^c}} \det(B_{m+1}+\sum_{i=1}^m p_iC_i^*B_iC_i)^{-1/2} \ds \prod_{i=1}^{m+1}\det(B_i)^{p_i/2} \\
\nm 
\leq \ds  \frac{1}{\sqrt{\det(\mathcal Q)}}  \sup_{\substack{0<B_i \leq Q_i \text{ for all } i \in I \cup \{m+1\} \\  B_{j} \geq Q_j \text{ for all } j \in I^c}}  \prod_{i=1}^{m} \det(B_i)^{p_i/2} \det(\mathrm{id}_H+M^{1/2}B_{m+1}^{-1}M^{1/2})^{-1/2}.
\end{array}$$
Note that for fixed $B_1,\cdots,B_m$, the quantity $\det(\mathrm{id}_H+M^{1/2}B_{m+1}^{-1}M^{1/2})^{-1/2}$ with $0<B_{m+1} \leq \mathcal Q$ is maximized when $B_{m+1}=\mathcal Q.$ This implies that $$\begin{array}{lll}\ds \frac{\ds \int e^{-\pi \langle \mathcal Qx,x\rangle}\prod_{i=1}^m f_i(C_ix)^{p_i}\, dx}{\ds \prod_{i=1}^m \Bigl(\int_{H_i}f_i(x_i) \, dx_i \Bigl)^{p_i}} &\leq& \ds  \sup_{\substack{0<B_i \leq Q_i \text{ for all } i \in I \\  B_{j} \geq Q_j \text{ for all } j \in I^c}} \det(\mathcal Q+M)^{-1/2}   \prod_{i=1}^m \det(B_i)^{p_i/2} \\ \\
&\leq& \ds  \sup_{\substack{0<B_i \leq Q_i \text{ for all } i \in I \\  B_{j} \geq Q_j \text{ for all } j \in I^c}}\frac{\ds \int e^{-\pi \langle \mathcal Qx,x\rangle}\prod_{i=1}^m g_{B_i}(C_ix)^{p_i}\, dx}{\ds \prod_{i=1}^m \Bigl(\int_{H_i}g_{B_i}(x_i) \, dx_i \Bigl)^{p_i}}.
\end{array}$$ 
Similarly, assuming that the dual regularized Brascamp-Lieb inequality with $\mathcal Q=0$ holds true, by taking $g_{m+1}=g_{\mathcal Q^{-1}}$, one gets $$\begin{array}{lll}
    \ds \frac{\ds \int_{H}g_{\mathcal Q^{-1}} * \sup \Bigl \{ \prod_{j=1}^m g_j(y_j)^{p_j}: \sum_{j=1}^m p_jC_j^*y_j=x, y_j \in H_j\Bigl \}\, dx}{\ds \prod_{i=1}^m \Bigl( \int_{H_i}g_i\Bigl)^{p_i}} \\
    \nm 
    \ds = \frac{\ds \int_{H}\sup \Bigl \{ \prod_{j=1}^{m+1} g_j(y_j)^{p_j}: \sum_{j=1}^{m+1} p_jC_j^*y_j=x, y_j \in H_j\Bigl \}\, dx}{\ds \prod_{i=1}^m \Bigl( \int_{H_i}g_i\Bigl)^{p_i}} \\ \\
    \geq \ds \Bigl( \int_H g_{m+1}(x)\, dx\Bigl) \inf_{\substack{0< B_i \leq Q_i \text{ for all } i \in I \cup \{m+1\} \\ B_j \geq Q_j \text{ for all } j \in I^c}}  \frac{\det(\sum_{i=1}^{m+1} p_i C_i^*B_iC_i)^{1/2}}{\prod_{i=1}^{m+1} (\det B_i)^{p_i/2}}  \\
    \nm 
    = \ds \sqrt{\det(\mathcal Q)} \inf_{\substack{0< B_i \leq Q_i \text{ for all } i \in I \cup \{m+1\} \\ B_j \geq Q_j \text{ for all } j \in I^c}}  \det(\mathrm{id}_H+M^{1/2}B_{m+1}^{-1}M^{1/2})^{1/2}\prod_{i=1}^m (\det B_i)^{-p_i/2}.
\end{array}$$ As for fixed $B_1,\cdots,B_m$, the quantity $\det(\mathrm{id}_H+M^{1/2}B_{m+1}^{-1}M^{1/2})^{1/2}$ with $0<B_{m+1} \leq \mathcal Q$ is minimized when $B_{m+1}=\mathcal Q$, one concludes in a similar fashion: $$\begin{array}{lll}
    \ds \frac{\ds \int_{H}g_{\mathcal Q^{-1}} * \sup \Bigl \{ \prod_{j=1}^m g_j(y_j)^{p_j}: \sum_{j=1}^m p_jC_j^*y_j=x, y_j \in H_j\Bigl \}\, dx}{\ds \prod_{i=1}^m \Bigl( \int_{H_i}g_i\Bigl)^{p_i}} \\
    \geq \ds \inf_{\substack{0< B_i \leq Q_i \text{ for all } i \in I \\ B_j \geq Q_j \text{ for all } j \in I^c}}  \det(\mathcal Q+M)^{1/2}\prod_{i=1}^m (\det B_i)^{-p_i/2} \\
    \nm 
    \geq \ds \inf_{\substack{0< B_i \leq Q_i \text{ for all } i \in I \\ B_j \geq Q_j \text{ for all } j \in I^c}} \Bigl\{\frac{\det(\mathcal Q+\sum_{i=1}^m p_iC_i^*B_iC_i)}{\prod_{i=1}^m \det(B_i)^{p_i}} \Bigl \}^{1/2} \\
    \nm \geq \ds \inf_{\substack{B_i \geq Q_i^{-1} \text{ for all } i \in I \\ 0<B_j \leq Q_j \text{ for all } j \in I^c}}\frac{\ds \int_H g_{\mathcal Q^{-1}} * \sup \Bigl \{ \prod_{j=1}^m g_{B_j}(y_j)^{p_j}: \sum_{j=1}^m p_jC_j^*y_j=x, y_j \in H_j\Bigl \}\, dx}{\ds \prod_{i=1}^m \Bigl( \int_{H_i}g_{B_i}\Bigl)^{p_i}}.
    \end{array}$$




By approximation for matrices, the previous statements are proven for any $\mathcal Q \geq 0.$ 
\end{remark}

\bigskip

We follow the strategy of Valdimarsson's proof (see \cite{Vald1}). 


\begin{proof}[Proof of \cref{Theorem_REG_BL}]
We can suppose that $H=\mathbb{R}^n, H_i=\mathbb{R}^{n_i}$ for $1 \leq i \leq m$ since any Euclidean space $E$ is isomorphic to $\mathbb{R}^{\mathrm{dim}(E)}.$ \hfill \\
Let us first show \cref{Reg_BL_Log_Convexity,Reg_Dual_BL_Log_Concavity}. 
Set $$J((f_i)_{i=1}^m)=\frac{\int_{\mathbb{R}^n}\prod_{i=1}^m f_i(C_i x)^{p_i}\, dx}{\prod_{i=1}^m (\int_{\mathbb{R}^{n_i}} f_i)^{p_i}}$$ and $$I((g_j)_{j=1}^m)=\frac{\int_{\mathbb{R}^n}\sup \Bigl \{ \prod_{j=1}^m g_j(y_j)^{p_j}: \sum_{j=1}^m p_jC_j^*y_j=x, y_j \in \mathbb{R}^{n_j}\Bigl \}\, dx}{\prod_{j=1}^m (\int_{\mathbb{R}^{n_j}}g_j)^{p_j}}.$$ Define $$F=\sup \Biggl \{ J((f_i)_{i=1}^m) : \substack{\ds f_i \text{ more log-convex than } g_{Q_i} \text{ for any } i \in I, \\ \ds f_i \text{ more log-concave than } g_{Q_j} \text{ for any } j \in I^c} \Biggl \},$$ \hfill \\  $$F_g=\sup_{\substack{0<B_i \leq Q_i \text{ for all } i \in I \\  B_j \geq Q_j \text{ for all } j \in I^c}}J((g_{B_i})_{i=1}^m),$$ and $$E=\inf \Biggl \{ I((g_j)_{j=1}^m): \substack{\ds g_i \text{ more log-concave than } g_{Q_i^{-1}} \text{ for any } i \in I,\\ \ds g_j \text{ more log-convex than } g_{Q_j^{-1}} \text{ for any } j \in I^c  \\ } \Biggl\},$$ \hfill \\ $$E_g=\inf_{\substack{B_i \geq Q_i^{-1} \text{ for all } i \in I \\ 0<B_j \leq  Q_j^{-1} \text{ for all } j \in I^c}}  I((g_{B_j})_{j=1}^m).$$

The goal is to show that $F=F_g, E=E_g$ and $E_g$ is equal to $F_g^{-1}.$ 

\hfill \\



As the infimum for $E_g$ is taken over a smaller class of functions than the infimum for $E$, $E_g \geq E$ and similarly, $F \geq F_g.$ From the previous remark, $$F_g=\frac{1}{\sqrt{D_g}}$$ where $$D_g=\inf_{\substack{0<B_i \leq Q_i \text{ for all } i \in I \\ B_j \ge Q_j \text{ for all } j \in I^c}} \Bigl \{ \frac{\det(\sum_{i=1}^m p_i C_i^*B_iC_i)}{\prod_{i=1}^m (\det B_i)^{p_i}}  \Bigl \}.$$ If $D_g=0$, there is nothing to show. Hence, we may assume that $D_g>0$ and denote $$Q(y)=\Bigl \langle \sum_{i=1}^m p_iC_i^*B_iC_iy,y \Bigl \rangle$$ for all $y \in H.$
Let us show that $$E_gF_g=1.$$ Take $0<B_i \leq Q_i$ for all $i \in I, B_j \ge Q_j$ for all $j \in I^c.$ 
Notice that $$I((g_{B_j^{-1}})_{j=1}^m)=\int e^{-\pi R(x)}\, dx$$ where $$R(x)=\inf \Bigl \{ \sum_{j=1}^m p_j \langle B_j^{-1}x_j,x_j\rangle : x=\sum_{j=1}^m p_jC_j^*x_j \text{ where } x_j \in \mathbb{R}^{n_j} \Bigl \}.$$ Let us recall that the dual of a quadratic form $\tilde Q$ is defined by $$\tilde Q^*(x)=\sup \{ |\langle x,y \rangle|^2 : \tilde Q(y) \leq 1\}.$$ It is proven in \cite[Lemma 2]{Barthe} that $R$ is a quadratic form and $R=Q^*.$ This implies that $$\frac{I((g_{B_j^{-1}})_{j=1}^m)}{\prod_{j=1}^m (\int_{\mathbb{R}^{n_j}}g_{B_j^{-1}}(x_j) \, dx_j)^{p_j}}= \Bigl (\frac{\prod_{j=1}^m (\det B_j^{-1})^{p_j} }{\det(R)} \Bigl)^{1/2}$$ and since $\det R=\det Q^*=(\det Q)^{-1}$, $$E_g=\sqrt{D_g}$$ as wished. 

If one manages to prove that $E \geq D_gF$, the following chain of inequalities allows to conclude: $$\sqrt{D_g}=E_g \geq E \geq D_gF \geq D_gF_g=\sqrt{D_g}.$$

Hence, it remains to show that $$I(g_1,\cdots,g_m) \geq D_g J(f_1,\cdots,f_m)$$ for $f_i$ more log-convex than $g_{Q_i}$ for any $i \in I$, $f_j$ more log-concave than $g_{Q_j}$ for any $j \in I^c$, $g_i$ more log-concave than $g_{Q_i^{-1}}$ for any $i \in I$ and $g_j$ more log-convex than $g_{Q_j^{-1}}$ for any $j \in I^c$ and $\int f_i=\int g_j=1.$ 

\hfill \\ We write $f_i=\exp(-V_i), g_i=\exp(-W_i)$ and by assumptions, $$\nabla^2 V_i \leq 2\pi Q_i, \nabla^2 W_i \geq 2\pi Q_i^{-1} \quad \text{ for all } i \in I$$ and $$\nabla^2 V_j \ge 2\pi Q_j, \nabla^2 W_j \le 2\pi Q_j^{-1} \quad \text{ for all } j \in I^c.$$

\hfill \\

Consider the optimal transport map $T_i=\nabla \phi_i$ between the probability measures $f_i$ and $g_i.$ 
Then, by \cref{General_Caff_Contraction}, it follows that $$0 \leq \nabla^2 \phi_i \leq Q_i \text{ for any } i \in I,$$ and $$\nabla^2 \phi_j \ge Q_j \text{ for any } j \in I^c.$$

In addition, the Monge-Ampère equation is satisfied with $$g_i(\nabla \phi_i(x))\det(\nabla^2\phi_i(x))=f_i(x).$$ Define $$\Theta(y)=\sum_{i=1}^m p_iC_i^*\nabla\phi_i(C_iy)$$ and notice that the Jacobian of $\Theta$ at $x$ is $$\det \Bigl(\sum_{i=1}^m p_iC_i^*\nabla^2 \phi_i(C_ix)C_i \Bigl)>0$$ by the assumption that $D_G>0.$ Using the Monge-Ampère equation, $$\begin{array}{lll}
J((f_i)_{i=1}^m) &=& \ds \int_{\mathbb{R}^n}\prod_{i=1}^m f_i^{p_i}(C_ix)\, dx \\
\nm 
&=& \ds \int_{\mathbb{R}^n}\prod_{i=1}^m g_i^{p_i}(\nabla \phi_i(C_ix))\prod_{i=1}^m (\det \nabla^2 \phi_i(C_ix))^{p_i}\, dx \\
\nm 
&\leq& \ds \frac{1}{D_g}\int_{\mathbb{R}^n}\prod_{i=1}^m g_i^{p_i}(\nabla \phi_i(C_ix)) \det \Bigl( \sum_{i=1}^m p_iC_i^* \nabla^2\phi_i(C_ix)C_i \Bigl ) \, dx \\
\nm 
&\leq& \ds \frac{1}{D_g}\int_{\mathbb{R}^n}\sup_{\Theta(x)=\sum p_iC_i^*x_i} \Bigl( \prod_{i=1}^m g_i^{p_i}(x_i) \Bigl) \det(d\Theta(x)) \, dx.
\end{array}$$ By the change of variables $y=\Theta(x)$, we get the desired claim. 
\end{proof}

\hfill \\


Notice that this strategy of proof may be applied to the work of Barthe and Wolff (see \cite{Barthe_Wolff}) to derive an inverse regularized Brascamp-Lieb inequality. More precisely, in this context, assuming that $\mathcal{Q}:H \rightarrow H$ is self-adjoint, $C_j:H \rightarrow H_j$ is surjective and $p$ is arranged such that $$p_1,\cdots,p_{m_+}>0>p_{m_++1},\cdots,p_m \text{ for some } 0\le m_+ \le m,$$  we say that $(\textbf{C},p,\mathcal Q,\tilde{\textbf{G}},\textbf{G})$ is non-degenerate if $$\begin{array}{lll}
\ds \mathcal Q|_{\ker(C_+)}>0, \quad \dim(H) \geq s^+(\mathcal Q)+\sum_{j=1}^{m_+}\dim(H_j) \\
\ds \mathcal Q+\sum_{j=1}^{m_+} p_jC_j^*G_jC_j>\sum_{j=m_++1}^m |p_j|C_j^*\tilde{G_j}C_j
\end{array}$$ where $$C_+:H \rightarrow \bigoplus_{j=1}^{m_+}H_j, \; C_+x=(C_1x,\cdots,C_{m^+}x).$$ 

Recall that one may decompose $\mathcal Q$ as $$\mathcal Q=p_0C_0^*\mathcal Q_+C_0+p_{m+1}C_{m+1}^*\mathcal Q_-C_{m+1}$$ where $p_0=-p_{m+1}=1$, $C_0: H \rightarrow H_0$ and $C_{m+1}: H \rightarrow H_{m+1}$ are surjective linear maps onto Euclidean spaces $H_0$ and $H_{m+1}$ such that $\mathcal Q_+, \mathcal Q_-$ are positive definite quadratic forms on $H_0$ and $H_{m+1}$ respectively. 



\begin{theorem} 
    Let $(\textbf{C},p,\mathcal Q,\textbf{G})$ a non-degenerate generalized Brascamp-Lieb datum. If $f_i \in L^1(H_i)$ is more log-concave than $g_{\tilde G_i}$ and more log-convex than $g_{G_i}$ for any $1 \le i \le m$, then
    \begin{equation} \label{Reg_Inverse_BL} 
    \frac{\ds \int_H e^{-\pi\langle Qx,x\rangle}\prod_{i=1}^m f_i(C_ix)^{p_i}\, dx}{\ds \prod_{i=1}^m \Bigl ( \int f_i \Bigl)^{p_i}}\ge \inf_{\tilde G_i \le B_i \le G_i}\Bigl\{ \frac{\prod_{i=1}^m (\det B_i)^{p_i}}{\det(\mathcal Q+\sum_{i=1}^m p_iC_i^*B_iC_i)} \Bigl\}^{1/2}.\end{equation}
    By duality, if $g_j \in L^1(H_j)$ is more log-concave than $g_{G_j^{-1}}$ and more log-convex than $g_{\tilde{G_j}^{-1}}$ for any $1 \le j \le m$, then $$\begin{array}{lll}
    \ds \frac{\ds \int_H \inf \left\{e^{-\pi \langle \mathcal Q_+^{-1}y_0,y_0\rangle}e^{\pi \langle \mathcal{Q}_-^{-1} y_{m+1},y_{m+1} \rangle} \prod_{j=1}^m g_j(y_j)^{p_j} : \sum_{j=0}^{m+1} p_jC_j^*y_j=y \right\}\, dy}{\ds \prod_{j=1}^m \Bigl( \int_{H_j}g_j\Bigl)^{p_j}} \\ 
    \ds \le \sup_{G_i^{-1} \le B_i \le \tilde{G_i}^{-1}} \frac{\ds \int_H \inf \left\{e^{-\pi \langle \mathcal Q_+^{-1}y_0,y_0\rangle}e^{\pi \langle \mathcal{Q}_-^{-1} y_{m+1},y_{m+1} \rangle} \prod_{j=1}^m g_{B_j}(y_j)^{p_j} : \sum_{j=0}^{m+1} p_jC_j^*y_j=y \right\}\, dy}{\ds \prod_{j=1}^m \Bigl( \int_{H_j}g_{B_j}\Bigl)^{p_j}} \\
    \ds \le \sup_{\tilde G_i \le B_i \le G_i} \Bigl\{ \frac{\det(\sum_{i=1}^m p_iC_i^*B_iC_i+\mathcal Q)}{\prod_{i=1}^m (\det B_i)^{p_i}} \Bigl\}^{1/2}.
    \end{array}$$ 
\end{theorem}

\hfill \\

Equation \cref{Reg_Inverse_BL} generalizes Theorem 2.1 of \cite{Nakamura} where the previous statement is established with $\tilde G_i=0$ and for type $G_i$ functions or in other words, convolutions with $N(0,G_i^{-1})$, \textit{i.e.} $$(\det G_i)^{1/2}\int_{\mathbb{R}^n}e^{-\pi \langle x-y,G_i(x-y)\rangle}\, d\mu(y)$$ for any $\mu$ positive finite Borel measure with non-zero mass. 
   




\hfill \\


\section{Finiteness and Structure of Regularized Brascamp-Lieb Inequalities}

\hfill \\

Now that the desired inequalities have been proved, it is natural to ask when the regularized Brascamp-Lieb constants are finite or even whether they are attained by a family of transformations $B_i: H_i \rightarrow H_i$ for $i=1,\cdots,m.$ \hfill \\

By convention, we will denote $n=\dim(H), n_i=\dim(H_i)$ for any $1 \leq i \leq m.$



\begin{definition}
    We define $$\mathrm{BL}(\textbf{C},p,\textbf{Q})=\sup_{0<B_i \leq Q_i} \Bigl \{ \frac{\prod_{i=1}^m (\det B_i)^{p_i}}{\det(\sum_{i=1}^m p_iC_i^*B_iC_i)} \Bigl \}^{1/2}$$ where $$\textbf{C}=\Bigl (H, (H_i)_{1 \leq i \leq m},(C_i)_{1 \leq i \leq m} \Bigl)$$ with $H,H_1,\cdots,H_m$ Euclidean spaces and for each $i$, $C_i: H \rightarrow H_i$ linear transformations (with $C$ non-degenerate \textit{i.e.} $C_i$ surjective and $\bigcap \ker(C_i)=\{0\}$), $p=(p_i)_{1 \leq i \leq m} \subset \mathbb{R}_+^m$ and $\textbf{Q}=(Q_i)_{1 \leq i \leq m}$ with $Q_i: H_i \rightarrow H_i, Q_i>0.$ Similarly, one may define \begin{equation} \label{Full_LogConvex_BL_Constant} \mathrm{BL}(\textbf{C},p,\textbf{Q},\mathcal{Q})=\sup_{0<B_i \leq Q_i}\Bigl \{ \frac{\prod_{i=1}^m (\det B_i)^{p_i}}{\det(\mathcal{Q}+\sum_{i=1}^m p_iC_i^*B_iC_i)} \Bigl \}^{1/2} \end{equation} and \begin{equation} \label{Mixed_BL_Constant} \mathrm{BL}(\textbf{C},p,\textbf{Q},I,\mathcal Q)=\sup_{\substack{0<B_i \le Q_i \; \forall i \in I \\ B_j \ge Q_j \; \forall j \in I^c}}\left\{ \frac{\prod_{i=1}^m (\det B_i)^{p_i}}{\det(\mathcal Q+\sum_{i=1}^m p_iC_i^*B_iC_i)} \right\}^{1/2}.\end{equation}
\end{definition}

In the latter, $(\textbf{C},p,\textbf{Q},I,\mathcal Q)$ (if $\mathcal Q=0$, we will write $(\textbf{C},p,\textbf{Q},I)$ and further, if $I=\{1,\cdots,m\}$, we will write $(\textbf{C},p,\textbf{Q})$) will be said to be a generalized Brascamp-Lieb datum and $\mathrm{BL}(\textbf{C},p,\textbf{Q},I,\mathcal Q)$ a generalized Brascamp-Lieb constant. 

\hfill \\

The conditions of finiteness, which will only be considered in the case $I=\{1,\cdots,m\}$, can immediately be treated for the case $\mathcal Q>0$. The supremum may be extended continuously to semi-definite $B_i$ and as a supremum over a compact interval, it is finite and attained. More precisely, $$\sup_{0<B_i \leq Q_i}\Bigl \{ \frac{\prod_{i=1}^m (\det B_i)^{p_i}}{\det(\mathcal Q+\sum_{i=1}^m p_iC_i^* B_iC_i)}\Bigl \}^{1/2}=\max_{0 \leq B_i \leq Q_i}\Bigl \{ \frac{\prod_{i=1}^m (\det B_i)^{p_i}}{\det(\mathcal Q+\sum_{i=1}^m p_iC_i^* B_iC_i)}\Bigl \}^{1/2}<+\infty.$$ Thus, an extremizer exists in the range $0 \leq B_i \leq Q_i$ but the expression vanishes when $\det(B_i)=0$ for some $i$ and hence at the extremum, one has $B_i>0$ for all $1 \leq i \leq m.$ 


\hfill \\ 


If $\ker(\mathcal Q) \neq \{0\}$, the sufficient and necessary conditions for the supremum in \cref{Reg_BL_Log_Convexity} to be finite can be formulated as follows.








\begin{theorem} \label{Necessary_suff_cond_finite_BL}
    The regularized BL constant $\mathrm{BL}(\textbf{C},p,\textbf{Q})$ is finite if and only if for any subspace $V \subset H$, \begin{equation} \label{Condition_BL_finite}\dim(V) \leq \sum_{j=1}^m p_j \dim(C_jV).\end{equation}
    More generally, if $\mathcal Q \geq 0$ and is not positive definite, $$\mathrm{BL}(\textbf{C},p,\textbf{Q},\mathcal Q)<+\infty$$ if and only if for any subspace $V \subset \ker(\mathcal Q),$ $$\dim(V) \leq \sum_{j=1}^m p_j\dim(C_jV).$$
\end{theorem}


\begin{remark}
    We will only write the proof for the case $\mathcal Q=0.$ 
    Indeed, if $\{0\} \subsetneq \ker(\mathcal Q) \subsetneq H$, the supremum may be extended to $0 \leq B_i \leq Q_i$ with $B_i|_{H_i/(C_i\ker(\mathcal Q))}>0$ and then, one can apply the same arguments that will be presented on $\ker(\mathcal Q)$ where obviously $\mathcal Q=0.$
\end{remark}

\begin{proof}
    We can prove easily the direction $\Longrightarrow.$ Indeed, take $V$ to be any subspace of $H, 0<\epsilon<1.$ Let us define $A^{(\epsilon)}=(A^{(\epsilon)}_j)_{1 \leq j \leq m}$ with $$A_j^{(\epsilon)}=\frac{\epsilon}{\lambda}\mathrm{id}_{C_jV} \, \oplus \, \frac{1}{\lambda}\mathrm{id}_{(C_jV)^{\perp}} \quad \text{ with } \lambda>\frac{1}{\min_{1 \leq j \leq m}\lambda_{\text{min}}(Q_j)}.$$ By the choice of $\lambda$ and as $0<\epsilon<1$, $0<A_j^{(\epsilon)} \leq Q_j$ for all $1 \leq j \leq m.$ It is clear that $\det(A_j^{(\epsilon)})$ decays like $\lambda^{-n_j}\epsilon^{\dim(C_jV)}$ as $\epsilon \rightarrow 0$ and $$\det \Bigl(\sum_{j=1}^m C_j^* A_j^{(\epsilon)}C_j \Bigl) \text{ decays at least as fast as } \lambda^{-n}\epsilon^{\dim(V)}$$ since $\sum_{j=1}^m C_j^*A_j^{(\epsilon)}C_j$ is bounded uniformly in $\epsilon$ and when restricted to $V$, decays linearly in $\epsilon$. Thus, as $\mathrm{BL}(\textbf{C},p,\textbf{Q})<+\infty$, we conclude that $$\lambda^{(n-\sum_{j=1}^m p_jn_j)/2}\epsilon^{(\sum_{j=1}^m p_j \dim(C_jV)-\dim(V))/2}<+\infty \text{ as } \epsilon \rightarrow 0$$ which means that $$\sum_{j=1}^m p_j\dim(C_jV) \geq \dim(V).$$
\end{proof}

To prove the opposite direction (which is an immediate consequence of \cref{Prop_Simple_GE}), some additional work needs to be done. We recall a very important definition to be able to study finiteness and extremizers.
 
\begin{definition}
We say that a subspace $V$ of $H$ is critical for a generalized Brascamp-Lieb datum $(\textbf{C},p,\textbf{Q})$  if $$\dim(V)=\sum_{j=1}^m p_j \dim(C_jV).$$
\end{definition}

For the following lemma, \cite[Definition 4.2]{BCCT} is employed; more precisely, $$\mathbf{C}|_V=(V,(C_iV)_{1 \le i \le m},(C_i|_V)_{1 \le i \le m}).$$ In addition, when writing $\textbf{Q}|_V$, we mean that $\textbf{Q}|_V=(Q_i|_{C_iV})_{1 \leq i \leq m}.$ 

\begin{lemma}
    Let $(\textbf{C},p,\textbf{Q})$ be a Brascamp-Lieb datum and let $V_c$ be a critical subspace. Then, $(\textbf{C},p,\textbf{Q})$ satisfies \cref{Condition_BL_finite} if and only if both $(\textbf{C}_{V_c},p,\textbf{Q}|_{V_c})$ and $(\textbf{C}_{H/V_c},p,\textbf{Q}|_{H/V_c})$ obey \cref{Condition_BL_finite}. 
\end{lemma}


\begin{proof}
    On the one hand, suppose that $(\textbf{C},p,\textbf{Q})$ satisfies \cref{Condition_BL_finite}. Then, it is clear that $(\textbf{C}_{V_c},p,\textbf{Q}|_{V_c})$ obeys \cref{Condition_BL_finite} by restricting the spaces $V \subseteq H$ to subspaces of $V_c.$ Let now $V$ be a subspace of $H/V_c.$ We notice that $$\begin{array}{lll}
    \dim(V) &=& \dim(V+V_c)-\dim(V_c) \\
    \nm 
    &\leq& \ds \sum_{j=1}^m p_j\dim(C_j(V+V_c))-\dim(V_c) \\
    \nm 
    &=& \ds \sum_{j=1}^m p_j(\dim(C_jV+C_jV_c)-\dim(C_jV_c)) \\
    \nm 
    &=& \ds \sum_{j=1}^m p_j \dim(C_{j,H/V_c}V)
    \end{array}$$ as desired. On the other hand, assume that $(\textbf{C}_{V_c},p,\textbf{Q}|_{V_c})$ and $(\textbf{C}_{H/V_c},p,\textbf{Q}|_{H/V_c})$ obey \cref{Condition_BL_finite}. For any $V \subseteq H$, writing $U=V \cap V_c, W=(V+V_c)/V_c$, one has $$\begin{array}{lll}
    \dim(V) &=& \dim(V \cap V_c)+\dim(V+V_c)-\dim(V_c) \\
    \nm 
    &=& \dim(U)+\dim(W) \\
    \nm
    &\leq& \ds \sum_{j=1}^m p_j \dim(C_{j,V}U)+p_j \dim(C_{j,H/V}W) \\
    \nm 
    &=& \ds \sum_{j=1}^m p_j \dim(C_jU)+p_j(\dim(C_jW+C_jV_c)-\dim(C_jV_c)) \\
    \nm 
    &=& \ds \sum_{j=1}^m p_j \Bigl[ \dim(C_jU)+\dim(C_jV+C_jV_c)-\dim(C_jV_c) \Bigl]
    \end{array}$$ since $W+V_c=V+V_c.$ As $$\dim(C_jU)=\dim(C_j(V \cap V_c)) \leq \dim(C_jV \cap C_jV_c)$$ and $$\dim(C_jV \cap C_jV_c)+\dim(C_jV+C_jV_c)-\dim(C_jV_c)=\dim(C_jV),$$ we see that $(\textbf{C},p,\textbf{Q})$ satisfies \cref{Condition_BL_finite}.  
\end{proof}

\begin{lemma}
    If $(\textbf{C},p,\textbf{Q})$ is a generalized Brascamp-Lieb datum, then $$\mathrm{BL}(\textbf{C},p,\textbf{Q}) \leq \mathrm{BL}(\textbf{C}_V,p,\textbf{Q}|_V)\mathrm{BL}(\textbf{C}_{H /V},p,\textbf{Q}|_{H/V}).$$
\end{lemma}


\begin{proof}
    We may assume that $\mathrm{BL}(\textbf{C}_V,p,\textbf{Q}|_V)$ and $\mathrm{BL}(\textbf{C}_{H/V},p,\textbf{Q}|_{H \setminus V})$ are finite. Take $f=(f_i)_{1 \leq i \leq m}$ with $f_i$ more log-convex than $g_{Q_i}$ for any $1 \leq i \leq m.$ We may normalize $\int_{H_j}f_j=1.$ The goal is to show that $$\int_H \prod_{j=1}^m (f_j \circ C_j)^{p_j}\, dx \leq \mathrm{BL}(\textbf{C}_V,p,\textbf{Q}|_V)\mathrm{BL}(\textbf{C}_{H/V},p,\textbf{Q}|_{H/V}).$$ By Fubini-Tonelli, the left-hand side can be rewritten as $$\int_{V^\perp}\Bigl (\int_V \prod_{j=1}^m (f_j \circ C_j)^{p_j}(v+w)\, dv \Bigl)\, dw.$$ However, one can write $f_j \circ C_j(v+w)=f_{j,w}\circ C_{j,V}(v)$ where $f_{j,w}: C_jV \rightarrow [0;+\infty)$ is the function $$f_{j,w}(v_j)=f_j(v_j+C_jw) \quad \forall v_j \in C_jV.$$ We claim that $f_{j,w}$ (respectively $f_{j,H/V}$, see the definition below) is more log-convex than $g_{Q_j|_{C_jV}}$ (resp. $g_{Q_j|_{C_j(H/V)}}.$) Suppose the claim is proven.
    
    Then, since $f_{j,w}$ is more-log convex than $g_{Q_j|_V}$, we conclude that $$\int_V \prod_{j=1}^m (f_j \circ C_j)^{p_j}(v+w)\, dv \leq \mathrm{BL}(\textbf{C}_V,p,\textbf{Q}|_V)\prod_{j=1}^m \Bigl (\int_{C_jV}f_{j,w}\Bigl)^{p_j}.$$ Therefore, it suffices to show that \begin{equation} \label{Split_eq}\int_{V^\perp}\prod_{j=1}^m \Bigl( \int_{C_jV}f_{j,w}\Bigl)^{p_j}\, dw \leq \mathrm{BL}(\textbf{C}_{H/V},p,\textbf{Q}|_{H/V}).\end{equation} Identifying $V^\perp$ with $H/V$, we observe that $\ds \int_{C_jV}f_{j,w}=f_{j,H/V}\circ C_{j,H/V}(w)$ where $f_{j,H/V}:H_j/(C_jV)\ \longrightarrow [0;+\infty)$ is the function $$f_{j,H/V}(w_j)=\int_{w_j+C_jV}f_j.$$ 
    Since $f_{j,{H/V}}$ is more log-convex than $g_{Q_j|_{H/V}}$ for every $1 \leq j \leq n$, by Fubini-Tonelli and the normalization $\int_{H_j}f_j=1$, \cref{Split_eq} holds true. \\ \hfill \\ Let us now prove the claim. In fact, setting $\tilde w=(1-\alpha)w_j+\alpha \tilde w_j, \tilde v=(1-\alpha)v_j+\alpha \tilde v_j$, one has using that $\log(f_j)+\pi \langle Q_j \cdot,\cdot \rangle$ is convex, the following for $\alpha \in (0;1)$:  \begin{equation} \label{Log_Convexity_Superposition} \begin{array}{lll}
    \\ \log(f_j(\tilde v+\tilde w)) = \log(f_j((1-\alpha)(v_j+w_j)+\alpha(\tilde v_j+\tilde w_j)) \\ \\
    \nm 
    = \log[f_j((1-\alpha)(v_j+w_j)+\alpha(\tilde v_j+\tilde w_j))] +\pi \langle Q_j (\tilde v+\tilde w),\tilde v+\tilde w \rangle \\ \\
    \nm 
    -\pi(1-\alpha)^2 \langle Q_j(v_j+w_j),v_j+w_j \rangle-2\pi \alpha(1-\alpha)\langle Q_j(v_j+w_j),\tilde v_j+\tilde w_j \rangle   \\ \\
    \nm 
    -\pi \alpha^2 \langle Q_j(\tilde v_j+\tilde w_j),\tilde v_j+\tilde w_j \rangle \leq (1-\alpha)\Bigl [ \log f_j(v_j+w_j)+\pi \langle Q_j(v_j+w_j),v_j+w_j \rangle \Bigl] \\ \\
    \nm 
    +\alpha \Bigl [ \log f_j(\tilde v_j+\tilde w_j)+\pi \langle Q_j(\tilde v_j+\tilde w_j),\tilde v_j+\tilde w_j\rangle \Bigl]-\pi(1-\alpha)^2 \langle Q_j(v_j+w_j),v_j+w_j \rangle \\ \\
    \nm 
    -2\pi \alpha(1-\alpha)\langle Q_j(v_j+w_j),\tilde v_j+\tilde w_j \rangle   -\pi \alpha^2 \langle Q_j(\tilde v_j+\tilde w_j),\tilde v_j+\tilde w_j  \rangle = (1-\alpha) \log f_j(v_j+w_j) \\ \\
    \nm 
    +\alpha \log f_j(\tilde v_j+\tilde w_j)+\pi \alpha(1-\alpha) \Bigl[ \langle Q_j(v_j+w_j),v_j+w_j\rangle-2\langle Q_j(v_j+w_j), \tilde v_j+\tilde w_j \rangle \\ \\
    \nm 
    +\langle Q_j(\tilde v_j+\tilde w_j),\tilde v_j+\tilde w_j \rangle \Bigl]=(1-\alpha) \log f_j(v_j+w_j)+\alpha \log f_j(\tilde v_j+\tilde w_j) \\ \\
    \nm 
    +\pi \alpha(1-\alpha) \langle Q_j(v_j+w_j-\tilde v_j-\tilde w_j),v_j+w_j-\tilde v_j-\tilde w_j \rangle.
    \end{array} \end{equation}
    In particular, taking $v_j=\tilde v_j$ in \cref{Log_Convexity_Superposition}, one obtains $$\begin{array}{lll}\log f_j(v_j+(1-\alpha)w_j+\alpha \tilde w_j) &\leq& (1-\alpha)\log f_j(v_j+w_j)+\alpha \log f_j(v_j+\tilde w_j) \\
    \nm 
    && +\pi \alpha(1-\alpha) \langle Q_j(w_j-\tilde w_j),w_j-\tilde w_j\rangle \end{array}$$ and taking the exponential, $$f_j(v_j+(1-\alpha)w_j+\alpha \tilde w_j) \leq f_j(v_j+w_j)^{1-\alpha}f_j(v_j+\tilde w_j)^{\alpha} e^{\pi \alpha (1-\alpha) \langle Q_j(w_j-\tilde w_j),w_j-\tilde w_j\rangle}.$$ By Hölder's inequality, $$\begin{array}{lll}
    \ds \int_{C_jV}f_j(v_j+(1-\alpha)w_j+\alpha \tilde w_j) \, dv_j &\leq& \ds e^{\pi \alpha(1-\alpha) \langle Q_j(w_j-\tilde w_j),w_j-\tilde w_j\rangle} \Bigl( \int_{C_jV}f_j(v_j+w_j)\, dv_j \Bigl)^{1-\alpha} \\
    \nm 
    && \ds \Bigl ( \int_{C_jV}f_j(v_j+\tilde w_j)\, dv_j\Bigl)^{\alpha}.
    \end{array}$$ This implies that $$\begin{array}{lll}
    \log f_{j,H/V}((1-\alpha)w_j+\alpha \tilde w_j)+\pi \langle Q_j((1-\alpha)w_j+\alpha \tilde w_j),(1-\alpha)w_j+\alpha \tilde w_j \rangle \\ \\
    \nm 
    \leq (1-\alpha)\log f_{j,H/V}(w_j)+\alpha \log f_{j,H/V}(\tilde w_j)+\pi \alpha(1-\alpha) \langle Q_j(w_j-\tilde w_j),w_j-\tilde w_j \rangle \\ \\
    \nm 
    +\pi(1-\alpha)^2 \langle Q_jw_j,w_j \rangle+2\pi \alpha(1-\alpha) \langle Q_jw_j,\tilde w_j\rangle+\pi \alpha^2 \langle Q_j \tilde w_j, \tilde w_j \rangle \\ \\
    \nm 
    = (1-\alpha) \Bigl[ \log f_{j,H/V}(w_j)+\pi \langle Q_jw_j,w_j \rangle \Bigl]+ \alpha \Bigl[\log f_{j,H/V}(\tilde w_j)+\pi \langle Q_j\tilde w_j,\tilde w_j \rangle \Bigl].
    \end{array}$$ Setting $\tilde w=w_j=\tilde w_j=C_jw$ in \cref{Log_Convexity_Superposition}, one gets $$\begin{array}{lll}
    \log f_{j,w}((1-\alpha)v_j+\alpha \tilde v_j) &=& \log f_j((1-\alpha)v_j+\alpha \tilde v_j+C_jw) \\
    \nm 
    &\leq& (1-\alpha)\log f_j(v_j+C_jw)+\alpha \log f_j(\tilde v_j+C_jw) \\
    \nm 
    && +\pi \alpha (1-\alpha)\langle Q_j(v_j-\tilde v_j),v_j-\tilde v_j\rangle 
    \end{array}$$ which implies that $$\begin{array}{lll}
    \log f_{j,w}((1-\alpha)v_j+\alpha \tilde v_j)+\pi \langle Q_j[(1-\alpha )v_j+\alpha \tilde v_j],(1-\alpha)v_j+\alpha \tilde v_j\rangle \leq \\
    \nm 
    \leq (1-\alpha)\Bigl[\log f_{j,w}(v_j)+\pi \langle Q_jv_j,v_j \rangle \Bigl] +\alpha \Bigl[ \log f_{j,w}(\tilde v_j)+\pi \langle Q_j\tilde v_j,\tilde v_j\rangle \Bigl]
    \end{array}$$ and concludes the proof of the claim. 
\end{proof} 

\begin{lemma}
    In fact, if $V$ is critical and $Q$ splits with respect to $V$, \textit{i.e.} $$\quad Q_j(C_j V) \subset C_jV,\quad Q_j(H_j/(C_j V)) \subset H_j/(C_jV) \quad \forall 1 \leq j \leq m,$$ then $$\mathrm{BL}(\textbf{C},p,\textbf{Q})=\mathrm{BL}(\textbf{C}_V,p,\textbf{Q}|_V)\mathrm{BL}(\textbf{C}_{H/V},p,\textbf{Q}|_{H/V}).$$
\end{lemma}

\begin{proof}
    By the previous lemma, we only need to show that $\geq$ holds.
    If \cref{Condition_BL_finite} does not hold for $(\textbf{C},p)$, there is nothing to show as $\mathrm{BL}(\textbf{C},p,\textbf{Q})=+\infty.$ \\ \hfill \\ Let $0<C_V<\mathrm{BL}(\textbf{C}_V,p,\textbf{Q}|_V), 0<C_{H/V}<\mathrm{BL}(\textbf{C}_{H/V},p,\textbf{Q}|_{H/V})$ be arbitrary constants. By \cref{Theorem_REG_BL}, there exist $0<B_i \leq Q_i|_{C_iV}$ and $0<B_i' \leq Q_i|_{H_i/(C_iV)}$ for any $1 \leq i \leq m$ such that $$\frac{\prod_{i=1}^m (\det B_i)^{p_i}}{\det(\sum_{i=1}^m p_iC_i^*B_iC_i)} \geq C_V^2, \quad \frac{\prod_{i=1}^m (\det B_i')^{p_i}}{\det(\sum_{i=1}^m p_iC_i^*B_i'C_i)} \geq C_{H/V}^2.$$ \hfill \\ Identifying $H_i/(C_iV)$ with $(C_iV)^\perp$ as a subset of $H_i$, we define $P_i: H_i \rightarrow C_iV$ orthogonal projection and $I-P_i: H_i \rightarrow H_i/(C_iV)$, which leads to considering $$\tilde B_i=P_i^* B_iP_i+(I-P_i)^*B_i'(I-P_i) \quad \forall 1 \leq i \leq m.$$ Then, $0<\tilde B_i \leq Q_i$ as $\textbf{Q}$ splits with respect to $V$ and thus $\langle Q_iP_ix,(I-P_i)x\rangle=0$ for any $x \in H_i$ which implies that $$\begin{array}{lll} \langle \tilde B_ix,x\rangle &=& \langle B_iP_ix,P_ix\rangle+\langle B_i'(I-P_i)x,(I-P_i)x\rangle \\
    \nm 
    &\leq& \langle Q_iP_ix,P_ix\rangle+\langle Q_i(I-P_i)x,(I-P_i)x\rangle \\
    \nm 
    &=& \langle Q_ix,x\rangle \quad \forall x \in H_i.
    \end{array}$$ 
    Furthermore, $$\begin{array}{lll} 
    \mathrm{BL}(\textbf{C},p,\textbf{Q})^2 &\geq& \ds \frac{\prod_{i=1}^m (\det \tilde B_i)^{p_i}}{\det(\sum_{i=1}^m p_iC_i^*\tilde B_iC_i)}=\frac{\prod_{i=1}^m (\det B_i)^{p_i}}{\det(\sum_{i=1}^m p_iC_i^*B_iC_i)} \times \frac{\prod_{i=1}^m (\det B_i')^{p_i}}{\det(\sum_{i=1}^m p_iC_i^*B_i'C_i)} \\ \\
    \nm 
    &\geq& C_V^2C_{H/V}^2. \end{array}$$
\end{proof}

\begin{remark}
    We do not expect this to hold if there exists $1 \leq j \leq m$ such that $Q_j$ does not split with respect to $V.$ As one is imposing a quite strong condition on the structure of $Q_j$ for any $1 \leq j \leq m$, \textit{i.e.} $$Q_j=\left(\begin{array}{ c | c } 
    C_jV & 0 \\
    \hline 
    0 & (C_jV)^\perp
    \end{array} \right),$$ the property of multiplicativity of Brascamp-Lieb constants is not completely preserved by regularizing.
\end{remark}


\begin{lemma}
    If $$\dim(V) \leq \sum_{j=1}^m p_j \dim(C_jV) \text{ for all subspaces } V \text{ of } H,$$ then there exists a positive real number $c>0$ such that for every orthonormal basis $e_1,\cdots,e_n$ of $H$, there exists a set $I_j \subseteq \{1,\cdots,n\}$ for each $1 \leq j \leq m$ with $|I_j|=\dim(H_j)$ so that \begin{equation} \label{EQ_cardinality}\sum_{j=1}^m p_j|I_j \, \cap \{k+1,\cdots,n\}| \geq n-k \quad \forall 0 \leq k \leq n \end{equation} and \begin{equation} \label{EQ_wedge} \Bigl \| \bigwedge_{i \in I_j}C_je_i \Bigl \|_{H} \geq c \quad \forall 1 \leq j \leq m.\end{equation} If there are no critical subspaces, we can enforce strict inequality in \cref{EQ_cardinality} for all $0 \leq k<n.$ 
\end{lemma}

\begin{proof}
    Recall that the space of all orthonormal bases is compact and the number of possible $I_j$ is finite. By continuity and compactness, it suffices to show that $$\bigwedge_{i \in I_j}C_je_i \neq 0 \text{ for all } 1 \leq j \leq m$$ to deduce \cref{EQ_wedge}. In other words, we want the vectors $(C_je_i)_{i \in I_j}$ to be linearly independent in $H_j$ for each $j.$ \\ In fact, take $I_j$ to be the set of indices $i$ for which $C_je_i$ is not in the linear span of $\{C_je_{i'}: i<i' \leq n\}.$ As $C_j$ is surjective, $|I_j|=\dim(H_j)$ for every $1 \leq j \leq m.$ Applying the assumption with $V$ equal to the span of $\{e_{k+1},\cdots,e_n\}$, one has $$\sum_{j=1}^m p_j \dim(C_jV) \geq n-k.$$ But by construction of $I_j$, $\dim(C_jV)=|I_j \cap \{k+1,\cdots,n\}|.$ If there are no critical subspaces, one can use that $$\sum_{j=1}^m p_j\dim(C_jV)>\dim(V) \text{ for any non-zero subspace } V \text{ of } H.$$ 
\end{proof}

\begin{proposition} \label{Prop_Simple_GE}
    Let $(\textbf{C},p,\textbf{Q})$ be a Brascamp-Lieb datum such that \cref{Condition_BL_finite} holds. Then, $\mathrm{BL}(\textbf{C},p,\textbf{Q})$ is finite. 
    Furthermore, if $(\textbf{C},p,\textbf{Q})$ is simple (in other words, there is no critical subspace), then $(\textbf{C},p,\textbf{Q})$ is Gaussian-extremizable. 
\end{proposition}

\begin{proof}
    Fix $A=(A_j)_{1 \leq j \leq m}$ with $0<A_j \leq Q_j.$ Set $M=\sum_{j=1}^m p_jC_j^*A_jC_j>0$ as a positive definite self-adjoint transformation. By choosing an appropriate orthonormal basis $\{e_1,\cdots,e_n\}$ of $H$, we can assume that $$M=\mathrm{diag}(\lambda_1,\cdots,\lambda_n) \text{ with } \lambda_1 \geq \cdots \geq \lambda_n>0.$$ Moreover, $M \leq \tilde M=\sum_{j=1}^m p_jC_j^*Q_jC_j$ and in particular, $$\lambda_1=\langle Me_1,e_1\rangle \leq \langle \tilde Me_1,e_1\rangle \leq \sup_{\|e\|=1}\langle \tilde Me,e\rangle=\lambda_{\max}(\tilde M).$$ By the previous lemma, we find $I_j \subset \{1,\cdots,n\}$ for each $1 \leq j \leq m$ of cardinality $\dim(H_j)$ obeying \cref{EQ_cardinality} and \cref{EQ_wedge}. Thus, for $i \in I_j$, $$\langle A_jC_je_i,C_je_i \rangle=\langle e_i, C_j^*A_jC_je_i \rangle \leq \frac{1}{p_j}\langle e_i,Me_i \rangle=\frac{\lambda_i}{p_j}.$$ As $(C_je_i)_{i \in I_j}$ is a basis of $H_j$ with a lower bound on the degeneracy, $$\det(A_j) \leq C\prod_{i \in I_j}\lambda_i$$ for some constant $C>0$ depending on $(\textbf{C},p).$ This implies that $$\prod_{j=1}^m \det(A_j)^{p_j} \leq C\prod_{i=1}^n \lambda_i^{\sum_{j=1}^n p_j|I_j \cap \{i\}|}.$$ Telescoping the right hand side, one gets $$\prod_{j=1}^m (\det A_j)^{p_j} \leq C\lambda_1^{\sum_{j=1}^m p_jn_j} \; \prod_{k=0}^{n-1}(\lambda_{k+1}/\lambda_k)^{\sum_{j=1}^m p_j|I_j \cap \{k+1,\cdots,n\}|}$$ with the convention $\lambda_0=\lambda_1.$ As $\lambda_{k+1} \leq \lambda_k$ for any $0 \leq k \leq n-1$ and $$\sum_{j=1}^m p_j|I_j \cap \{k+1,\cdots,n\}| \geq n-k,$$ one concludes that $$\begin{array}{lll}
    \ds \prod_{j=1}^m (\det A_j)^{p_j} &\leq& \ds C\lambda_1^{\sum_{j=1}^m p_jn_j}\prod_{k=0}^{n-1}(\lambda_{k+1}/\lambda_k)^{n-k} \\
    \nm 
    &\leq& \ds C\lambda_1^{\sum_{j=1}^m p_jn_j-n}\prod_{k=1}^{n}\lambda_k. \\
    \nm 
    &\leq& \ds C\lambda_{\max}(\tilde M)^{\sum_{j=1}^m p_jn_j-n}\det(M)=\tilde{C}\det(M)
    \end{array}$$ for some constant $\tilde C>0$ depending on $(\textbf{C},p,\textbf{Q})$, using that $\sum_{j=1}^m p_jn_j \geq n.$ 



    
    Let us assume that $(\textbf{C},p,\textbf{Q})$ is simple. We then have a strict inequality in \cref{EQ_cardinality} which gives that $$\prod_{j=1}^m (\det A_j)^{p_j} \leq C \det(M) \prod_{k=1}^{n-1}(\lambda_{k+1}/\lambda_k)^c=C \det(M) (\lambda_n /\lambda_1)^c$$ for some $c>0$ depending on the Brascamp-Lieb data. Thus, $\prod_{j=1}^m (\det A_j)^{p_j}/\det(M)$ converges to $0$ whenever $\lambda_n/\lambda_1$ goes to $0.$ This proves that it suffices to evaluate the supremum in the region $\lambda_1(M) \leq C\lambda_n(M).$ \\ Assume by contradiction that $(\textbf{C},p,\textbf{Q})$ is not Gaussian-extremizable. This means there exists a sequence $(A_j^{(\epsilon)})_{1 \leq j \leq m}$ such that $0<A_j^{(\epsilon)} \leq Q_j \quad \forall 1 \leq j \leq m$, so that $$\frac{\prod_{j=1}^m \det(A_j^{(\epsilon)})^{p_j}}{\det(\sum_{j=1}^m p_jC_j^*A_j^{(\epsilon)}C_j)} \rightarrow_{\epsilon \downarrow 0} \mathrm{BL}(\textbf{C},p,\textbf{Q})>\frac{\prod_{j=1}^m \det(B_j)^{p_j}}{\det(\sum_{j=1}^m p_jC_j^*B_jC_j)}$$ for any $0<B_j \leq Q_j$ and the determinant of $M^{(\epsilon)}=\sum_{j=1}^m p_jC_j^* A_j^{(\epsilon)}C_j$  converges to $0.$ Since $M^{(\epsilon)}$ is symmetric, one can write $$M^{(\epsilon)}=U^{(\epsilon)}\mathrm{diag}(\lambda_1^{(\epsilon)},\cdots,\lambda_n^{(\epsilon)})(U^{(\epsilon)})^T \text{ for } U^{(\epsilon)} \in O(n), \text{ with } \lambda_1^{(\epsilon)} \geq \cdots \geq \lambda_n^{(\epsilon)}>0.$$ Then, $\lambda_n^{(\epsilon)} \rightarrow_{\epsilon \rightarrow 0} 0$ and by what was previously shown, $C\lambda_n^{(\epsilon)} \geq \lambda_1^{(\epsilon)}$. This means that $\lambda_1^{(\epsilon)} \rightarrow_{\epsilon \rightarrow 0} 0$ with the same rate of convergence than $\lambda_n^{(\epsilon)}$, which will be denoted by $f(\epsilon).$ Consequently, $\lambda_k^{(\epsilon)} \rightarrow_{\epsilon \rightarrow 0} 0$ for any $1 \leq k \leq n$ with rate of convergence $f(\epsilon).$   
    Moreover, by compactness of $O(n)$, up to a subsequence, $U^{(\epsilon)} \rightarrow U \in O(n)$ which proves that $M^{(\epsilon)} \rightarrow_{\epsilon \rightarrow 0} 0.$ Hence, $$\sum_{j=1}^m p_j\langle A_j^{(\epsilon)}C_jx,C_jy\rangle=\langle M^{(\epsilon)}x,y\rangle \rightarrow 0 \quad \forall x,y \in H$$ and by surjectivity of $C_j$ and $p_j>0$ for $j=1,\cdots,m$, $$\langle A_j^{(\epsilon)}v_j,w_j \rangle \rightarrow_{\epsilon \rightarrow 0} 0 \quad \forall v_j, w_j \in H_j \Longrightarrow A_j^{(\epsilon)} \rightarrow 0 \quad \forall 1 \leq j \leq m.$$ 
Now, write $$\ds \frac{\prod_{j=1}^m (\det A_j^{(\epsilon)})^{p_j}}{\det(M^{(\epsilon)})}=f(\epsilon)^{\sum_{j=1}^m p_jn_j-n}\frac{\prod_{j=1}^m \det \Bigl(\frac{A_j^{(\epsilon)}}{f(\epsilon)} \Bigl)^{p_j}}{\det \Bigl( \frac{M^{(\epsilon)}}{f(\epsilon)}\Bigl)}$$ and notice that $M^{(\epsilon)}/f(\epsilon) \rightarrow_{\epsilon \rightarrow 0}U\mathrm{diag}(\tilde \lambda_1,\cdots,\tilde \lambda_n)U^T=\tilde M>0$ up to a subsequence. Furthermore, for $v_j^{(k,\epsilon)}$ a unit norm eigenvector of $A_j^{(\epsilon)}$ associated to its $k$-th eigenvalue $\lambda_k(A_j^{(\epsilon)})$, $$\frac{\lambda_k(A_j^{(\epsilon)})}{f(\epsilon)}=\Bigl\langle \frac{A_j^{(\epsilon)}v_j^{(k,\epsilon)}}{f(\epsilon)},v_j^{(k,\epsilon)}\Bigl\rangle \leq \frac{1}{p_j} \Bigl\langle \frac{M^{(\epsilon)}}{f(\epsilon)}v^{(\epsilon)},v^{(\epsilon)} \Bigl\rangle$$ where $v^{(\epsilon)}=D_jv_j^{(k,\epsilon)}$ for $D_j$ a right-inverse of $C_j$, which exists since $C_j$ is surjective (\textit{i.e.} $C_jD_j=\mathrm{id}_{H_j}$ so that $C_jv^{(\epsilon)}=v_j^{(k,\epsilon)}$). As $v^{(\epsilon)}$ has a bounded norm and consequently has a convergent subsequence by compactness, $\frac{\lambda_k(A_j^{(\epsilon)})}{f(\epsilon)}$ is bounded in $\epsilon$ up to a subsequence. Thus, we can conclude with a contradiction since $\sum_{j=1}^m p_jn_j-n>0$, $f(\epsilon) \rightarrow_{\epsilon \rightarrow 0} 0$ and $$\frac{\prod_{j=1}^m \det\Bigl( \frac{A_j^{(\epsilon)}}{f(\epsilon)}\Bigl)^{p_j}}{\det \Bigl( \frac{M^{(\epsilon)}}{f(\epsilon)}\Bigl)}$$ is uniformly bounded in $\epsilon$, which would mean that $\mathrm{BL}(\textbf{C},p,\textbf{Q})=0$ (which is, of course, impossible).  
\end{proof}


\if0
Let us now study the structure of the mixed BL of \cref{Reformulation_REG_BL}. 

For $I \subset \{1,\cdots,m\}$, we define $$\mathrm{BL}(\mathbf{C},p,\mathbf{Q},I,\mathcal Q)=\sup \left\{ \frac{\prod_{k=1}^m (\det B_k)^{p_k}}{\det(\mathcal Q+\sum_{k=1}^m p_kC_k^*B_kC_k)} : 0<B_i \le Q_i \quad \forall i \in I, B_i \ge Q_i \quad \forall i \in I^c\right\}^{1/2}.$$ 

\hfill \\


\begin{theorem}
If $I \subset \{1,\cdots,m\}$, $$\mathrm{BL}(\mathbf{C},p,\mathbf{Q},I,\mathcal Q)<+\infty$$ if and only if $$\dim(V) \ge \sum_{j \in I^c}p_j\dim(C_jV) \quad \forall V \text{ subspace of } H$$ and $$\dim(V) \le \sum_{j \in I}p_j\dim(C_jV) \quad \forall V \text{ subspace of } \ker \Bigl(\mathcal Q+\sum_{j \in I^c}p_jC_j^*Q_jC_j \Bigl).$$
In particular, if $\mathcal{Q}>0$, $$\mathrm{BL}(\textbf{C},p,\textbf{Q},I,\mathcal Q)<+\infty$$ if and only if $$\dim(V) \ge \sum_{j \in I^c}p_j\dim(C_jV) \quad \forall V \text{ subspace of } H.$$
\end{theorem}

%{\color{red} This holds for sure when $\mathcal Q>0$ by the next argument, $\mathcal Q=0$ is not yet completely clear to me. Maybe, we also need to add $$\dim(V) \le \sum_{j \in I}p_j\dim(C_jV) \quad \forall V \text{ subspace of } H?$$ This would be great and would allow to unify completely the structural properties of the BL.}



\begin{proof}
Indeed, for the implication $\Longrightarrow$ , note that $\mathrm{BL}(\textbf{C},p,\textbf{Q},I,\mathcal Q)<+\infty$ implies in particular that $$ \prod_{j \in I^c}(\det Q_j)^{p_j} \sup_{0<B_i \le Q_i\; \forall i \in I} \left\{ \frac{\prod_{i \in I}(\det B_i)^{p_i}}{\det(\mathcal Q+\sum_{j \in I^c}p_jC_j^*Q_jC_j+\sum_{i \in I}p_iC_i^*B_iC_i)} \right\}<+\infty.$$ By \cref{Necessary_suff_cond_finite_BL}, this implies that $$\dim(V) \le \sum_{j \in I^c}p_j\dim(C_jV) \quad \forall V \text{ subspace of } \ker \Bigl(\mathcal Q+\sum_{j \in I^c}p_jC_j^*Q_jC_j \Bigl).$$



Moreover, $$\mathrm{BL}(\textbf{C},p,\textbf{Q},I,\mathcal Q)<+\infty$$ also implies that $$\sup_{B_j \ge Q_j \; \forall j \in I^c} \left\{ \frac{\prod_{j \in I^c}(\det B_j)^{p_j}}{\det(\mathcal Q+\sum_{i \in I}p_iC_i^*Q_iC_i+\sum_{j \in I^c}p_jC_j^*B_jC_j)} \right\}<+\infty.$$ 
Take $V$ to be any subspace of $H$ and for $0<\epsilon<1$, define $$B_j^{(\epsilon)}=\frac{\lambda}{\epsilon}\mathrm{id}_{C_jV} \oplus \lambda \mathrm{id}_{(C_jV)^\perp} \text{ with } \lambda>\max_{j \in I^c}\lambda_{\max}(Q_j).$$ 

Then, it is clear that $\det(B_j^{(\epsilon)})$ grows like $\lambda^{n_j}\epsilon^{-\dim(C_jV)}$ as $\epsilon \rightarrow 0$ and notice that $$\begin{array}{lll}
\ds \epsilon^n \det\Bigl(\mathcal Q+\sum_{i \in I}C_i^*Q_iC_i+\sum_{j \in I^c} C_j^*B_j^{(\epsilon)}C_j\Bigl) &=& \ds \det \Bigl(\epsilon \tilde{\mathcal Q}+\sum_{j \in I^c}C_j^*(\lambda \mathrm{id}_{C_jV} \oplus \epsilon \lambda \mathrm{id}_{(C_jV)^\perp})C_j \Bigl) \\
&\le& \ds C\epsilon^{\dim(V^\perp)}
\end{array}$$ where $\tilde{\mathcal Q}=\mathcal Q+\sum_{i \in I}C_i^*Q_iC_i$ and $C$ is a constant independent of $\epsilon.$ In other words, $$\det\Bigl(\mathcal Q+\sum_{i \in I}C_i^*Q_iC_i+\sum_{j \in I^c} C_j^*B_j^{(\epsilon)}C_j\Bigl) \le C\epsilon^{-\dim(V)}.$$ This means that $$\epsilon^{\dim(V)-\sum_{j \in I^c}p_j\dim(C_jV)}<+\infty \text{ as } \epsilon \rightarrow 0$$ and hence, $$\dim(V) \ge \sum_{j \in I^c}p_j\dim(C_jV).$$
\hfill \\
For the reverse direction ({\color{red} For the moment, I am only sure of how to prove this direction with $\mathcal Q>0$}), we may use an adapted version of Lemma 3.9 which will say that if $$\dim(V) \ge \sum_{j \in I^c}p_j\dim(C_jV) \text{ for all subspaces } V \text{ of } H,$$ then there exists a positive real number $c>0$ such that for every orthonormal basis $e_1,\cdots,e_n$ of $H$, there exists a set $I_j \subseteq \{1,\cdots,n\}$ for each $j \in I^c$ with $|I_j|=\dim(H_j)$ so that $$\sum_{j \in I^c} p_j|I_j \cap \{k+1,\cdots,n\}| \le n-k \quad \forall 0 \le k \le n$$ and $$\Bigl\| \bigwedge_{i \in I_j}C_je_i \Bigl\|_H \ge c \quad \forall j \in I^c.$$ Fix $B=(B_j)_{1 \le j \le m}$ with $0<B_j \le Q_j$ for all $j \in I$ and $B_j \ge Q_j$ for all $j \in I^c.$
Set $M=\sum_{j=1}^m p_jC_j^*B_jC_j+\mathcal Q.$
As in the proof of \cref{Prop_Simple_GE}, $$M=\mathrm{diag}(\lambda_1,\cdots,\lambda_n) \text{ with } 0<\lambda_1 \le \cdots \le \lambda_n.$$ Moreover, $$M \ge \sum_{j \in I^c}p_jC_j^*Q_jC_j+\mathcal Q \ge \mathcal Q>0$$ which implies that $\lambda_1 \ge \lambda_{\mathrm{min}}(\mathcal Q).$ Then, $$\begin{array}{lll}
\ds \prod_{j \in I^c}(\det B_j)^{p_j} &\le& \ds C\lambda_1^{\sum_{j \in I^c}p_j n_j} \prod_{k=0}^{n-1}(\lambda_{k+1}/\lambda_k)^{\sum_{j \in I^c}p_j|I_j \cap \{k+1,\cdots,n\}|} \\
&\le& \ds C\lambda_1^{\sum_{j \in I^c}p_jn_j}\prod_{k=0}^{n-1}(\lambda_{k+1}/\lambda_k)^{n-k} \\
&\le& \ds C\lambda_1^{\sum_{j \in I^c}p_jn_j-n}\prod_{k=1}^n \lambda_k \\
&\le& \ds C\lambda_{\min}(\mathcal Q)^{\sum_{j \in I^c}p_jn_j-n}\det(M)=\tilde C \det(M)
\end{array}$$ using that $\sum_{j \in I^c}p_jn_j \ge n.$
\end{proof}
\fi


\bigskip

\begin{remark}
 Recall the following classical fact. When $f$ is more log-convex (resp. more log-concave) than $g_A$, g more log-convex (resp. more log-concave) than $g_B$ for some $A,B>0$, $f \star g$ is more log-convex (resp. more log-concave) than $g_C$ for $C=A-A(A+B)^{-1}A$. Indeed, since $f=g_Ah$ and $g=g_B\tilde h$ for log-convex (resp. log-concave) functions $h, \tilde h$, one has $$\begin{array}{lll}
\ds e^{\pi \langle Cx,x \rangle}(f \star g)(x) &=& \ds \int_{\mathbb{R}^n} h(x-y)\tilde h(y) e^{\pi \langle Cx,x\rangle}e^{-\pi \langle A(x-y), x-y\rangle}e^{-\pi \langle By,y\rangle}\, dy \\
\nm 
&=& \ds \int_{\mathbb{R}^n}h(x-y)\tilde h(y) e^{-\pi \langle M(x,y),(x,y)\rangle}\, dy
\end{array}$$ where $$M=\begin{pmatrix}
    A(A+B)^{-1}A &  -A \\ 
   -A & A+B
\end{pmatrix}.$$ The log-convexity (resp. log-concavity) of $h$ and $\tilde h$ implies the log-convexity (resp. log-concavity) of $(x,y) \mapsto h(x-y)\tilde h(y)$, which gives, by \cite[Theorem 4.3]{Brascamp_Lieb}, that $f \star g$ is more log-convex (resp. log-concavity) than $g_C.$ 
\\ 
Define $$\mathrm{Conv}(f,g)=2^{n/2}(f \star g)(\sqrt{2}\;\cdot)$$ for $f,g \in L^1(\mathbb{R}^n,\mathbb{R}_+).$ From what was proven, if $f$ and $g$ are more log-convex (resp. log-concave) than $g_A$, $\mathrm{Conv}(f,g)$ is also more log-convex (resp. log-convex) than $g_A.$ 
\end{remark}

\hfill \\

Using convolutions and the central limit theorem allows to prove the following: 

\begin{theorem} \label{Ext_implies_Gaussian_ext}
If the Brascamp-Lieb constant $\mathrm{BL}(\textbf{C},p,\textbf{Q},I)$ is extremizable (\textit{i.e.} there is equality for some functions $f_i$ that are more log-convex than $g_{Q_i}$, $i \in  I$ and more log-concave than $g_{Q_i}$, $i \in I^c$), it is Gaussian extremizable (\textit{i.e.} the supremum in the definition of $\mathrm{BL}(\textbf{C},p,\textbf{Q},I)$ is attained). 
\end{theorem}


\begin{proof}
The computation will first start by considering the more general $\mathrm{BL}(\textbf{C},p,\textbf{Q},I,\mathcal{Q})$, that is defined in \cref{Mixed_BL_Constant}. One can assume as usual that $H=\mathbb{R}^n, H_i=\mathbb{R}^{n_i}$ for some $n_i \leq n, \; 1 \leq i \leq m.$ Let $f_i, g_i$ be more log-convex than $g_{Q_i}$ for any $i \in I$, more log-concave than $g_{Q_i}$ for any $i \in I^c$ and $\mathcal{Q} \geq 0.$ Assume without loss of generality that $\int_{\mathbb{R}^{n_i}}f_i=\int_{\mathbb{R}^{n_i}}g_i=1$ and set $$F(x)=e^{-\pi \langle \mathcal{Q}x,x\rangle}\prod_{i=1}^m f_i(C_ix)^{p_i}, G(y)=e^{-\pi \langle \mathcal Qy,y \rangle}\prod_{i=1}^m g_i(C_iy)^{p_i}.$$ Then, \begin{equation} \label{Ext_implies_Gaussianext_eq}\begin{array}{lll}
\ds \int_{\mathbb{R}^n}F \int_{\mathbb{R}^n}G &=& \ds \int_{\mathbb{R}^n}F \star G(x)\, dx \\
\nm 
&=& \ds 2^{n/2} \int_{\mathbb{R}^n}\int_{\mathbb{R}^n} e^{-\pi \langle \mathcal{Q}y,y\rangle} e^{-\pi \langle \mathcal Q(\sqrt{2}x-y),\sqrt{2}x-y\rangle} \prod_{i=1}^m \Bigl(f_i(\sqrt{2}C_ix-C_iy)g_i(C_iy)\Bigl)^{p_i} \, dx dy \\
\nm 
&=& \ds \int_{\mathbb{R}^n} \int_{\mathbb{R}^n} e^{-\pi \langle \mathcal Q \frac{x+y}{\sqrt{2}},\frac{x+y}{\sqrt{2}}\rangle}e^{-\pi \langle \mathcal Q \frac{x-y}{\sqrt{2}},\frac{x-y}{\sqrt{2}} \rangle} \prod_{i=1}^m \Bigl(f_i \bigl(\frac{C_ix-C_iy}{\sqrt{2}} \bigl)g_i \bigl(\frac{C_ix+C_iy}{\sqrt{2}} \bigl)\Bigl)^{p_i} \, dx dy \\
\nm 
&=& \ds \int_{\mathbb{R}^n} e^{-\pi \langle \mathcal Qx,x \rangle} \Bigl( \int_{\mathbb{R}^n} e^{-\pi \langle \mathcal Qy,y \rangle} \prod_{i=1}^m F_i^{(C_ix)}(C_iy)^{p_i} \, dy \Bigl) \, dx
\end{array}\end{equation} where $F_i^{(C_ix)}(y_i)=f_i(\frac{C_ix-y_i}{\sqrt{2}})g_i(\frac{C_ix+y_i}{\sqrt{2}}).$
In the third line, the change of variables is $y \mapsto \frac{x+y}{\sqrt 2}$ and in the last line, it was used that $$\Bigl \langle \mathcal Q\frac{x+y}{\sqrt{2}},\frac{x+y}{\sqrt 2}\Bigl \rangle+\Bigl \langle \mathcal Q\frac{x-y}{\sqrt{2}},\frac{x-y}{\sqrt 2}\Bigl \rangle=\langle \mathcal Qx,x\rangle+\langle \mathcal Qy,y\rangle \quad \forall x,y \in \mathbb{R}^n.$$


As $f_i$ and $g_i$ are more log-convex than $g_{Q_i}$ for any $i \in I$ (resp. more log-concave than $g_{Q_i}$ for any $i \in I^c$), there exist log-convex functions (resp. log-concave) $h_i, \tilde h_i$ such that $f_i(x_i)=e^{-\pi \langle Q_ix_i,x_i\rangle}h_i(x_i), g_i(x_i)=e^{-\pi \langle Q_ix_i,x_i \rangle}\tilde h_i(x_i).$ Hence, we deduce that $F_i^{(C_ix)}$ is also more log-convex (resp. more log-concave) than $g_{Q_i}$ for any $i \in I$ (resp. $i \in I^c$) as $$F_i^{(C_ix)}(y_i)=e^{-\pi \langle Q_iC_ix,C_ix\rangle}e^{-\pi \langle Q_iy_i,y_i \rangle}h_i \bigl(\frac{C_ix-y_i}{\sqrt{2}} \bigl)\tilde h_i \bigl(\frac{C_ix+y_i}{\sqrt{2}} \bigl)$$ and notice that $$\mathrm{Conv}(f_i,g_i)(C_ix)=2^{n_i/2}\int_{\mathbb{R}^{n_i}} f_i(\sqrt{2}C_ix_i-y_i)g_i(y_i)\, dy_i= \int_{\mathbb{R}^{n_i}}F_i^{(C_ix)}(y_i)\, dy_i$$ by the change of variables from $y_i$ to $\frac{C_ix_i+y_i}{\sqrt{2}}.$  

Thus, denoting $\mathrm{Conv}(f,g)=(\mathrm{Conv}(f_i,g_i))_{1 \leq i \leq m}$, \begin{equation} \label{J_Q_increases_under_CONV} \begin{array}{lll}
\ds J_{\mathcal Q}(f)J_{\mathcal Q}(g)=\int_{\mathbb{R}^n}F \int_{\mathbb{R}^n}G &\leq& \ds \mathrm{BL}(\textbf{C},p,\textbf{Q},I,\mathcal Q)\int_{\mathbb{R}^n}e^{-\pi \langle \mathcal Qx,x \rangle} \prod_{i=1}^m \Bigl ( \int_{\mathbb{R}^{n_i}}F_i^{(C_ix)}(y_i)\, dy_i\Bigl)^{p_i} \, dx \\
\nm 
&\leq& \mathrm{BL}(\textbf{C},p,\textbf{Q},I,\mathcal Q) J_{\mathcal Q}(\mathrm{Conv}(f,g)).
\end{array} \end{equation} Here, we set $$J_{\mathcal Q}((h_i)_{i=1}^m)=\frac{\ds \int e^{-\pi \langle \mathcal Qx,x\rangle}\prod_{i=1}^m h_i(C_ix)^{p_i}\, dx}{\ds \prod_{i=1}^m \Bigl( \int h_i \Bigl)^{p_i}}.$$ \hfill \\ %For the rest of the argument, we can take $\mathcal Q=0$ (and we will write $J=J_0$) as we already know that $\mathrm{BL}(\textbf{C},p,\textbf{Q},\mathcal Q)$ is Gaussian extremizable when $\mathcal Q>0$. 
Now, from the previous inequality, if $f$ is an extremizer for $J_{\mathcal Q}$ with $\int f_i=1$ for any $1 \leq i \leq m$, $\mathrm{Conv}(f):=\mathrm{Conv}(f,f)$ is also an extremizer as $$\begin{array}{lll}J_{\mathcal Q}(\mathrm{Conv}(f)) &\geq& \ds \frac{J_{\mathcal Q}(f)^2}{\mathrm{BL}(\textbf{C},p,\textbf{Q},I,\mathcal Q)}=\mathrm{BL}(\textbf{C},p,\textbf{Q},I,\mathcal Q)\\
\nm 
&=& \ds \mathrm{BL}(\textbf{C},p,\textbf{Q},I,\mathcal Q)\prod_{i=1}^m \Bigl( \int \mathrm{Conv}(f_i) \Bigl)^{p_i}\end{array}$$ with $\mathrm{Conv}(f_i)$ being more log-convex than $g_{Q_i}$ for any $i \in I$ and more log-concave than $g_{Q_i}$ for any $i \in I^c$, according to the previous remark. Iterating this process, we deduce that $\mathrm{Conv}^k(f)$ is also an extremizer for all $k \geq 1.$ By the central limit theorem (for example, see \cite[Thm 1.1]{Bobkov}), each $\mathrm{Conv}^k(f_i)$ converges in $L^1(\mathbb{R}^{n_i})$ and almost everywhere as $k \rightarrow +\infty$ to a centered Gaussian $\frac{g_{B_i}}{Z_i}$, which is still more log-convex than $g_{Q_i}$ for any $i \in I$ and more log-concave than $g_{Q_i}$ for any $i \in I^c.$
By Fatou's lemma and \cref{Closure_prop_extremizers} (ii),(iv) with $g=(g_{B_1},\cdots,g_{B_m})$ and $0<B_i \leq Q_i$ for all $i \in I$, $B_i \ge Q_i$ for all $i \in I^c$, $$J_{\mathcal Q}(g) \geq \mathrm{BL}(\textbf{C},p,\textbf{Q},I,\mathcal Q) \Bigl(\int_{\mathbb{R}^{n_i}}g_{B_i} \Bigl)^{p_i}$$ and this proves that the BL constant is Gaussian-extremizable, as desired. 
\end{proof}



\hfill \\


Thanks to the previous lemma, we may easily obtain the following closure properties of extremizers of $J_Q$: 

\begin{lemma} \label{Closure_prop_extremizers}
    Let $(\textbf{C},p,\textbf{Q},I,\mathcal Q)$ with $\mathcal Q \geq 0$ be a generalized Brascamp-Lieb datum with $\mathrm{BL}(\textbf{C},p,\textbf{Q},I,\mathcal Q)$ finite.
    \begin{itemize}
        \item [(i)] If $f=(f_j)_{1 \leq j \leq m}$ is an extremizer, then so is $(f_j(-\cdot))_{1 \leq j \leq m}.$
        \item[(ii)] If $f=(f_j)_{1 \leq j \leq m}$ is an extremizer, then so is $(c_jf_j)_{1 \leq j \leq m}$ for any positive real numbers $c_1,\cdots,c_m.$
        \item[(iii)] If $f=(f_j)_{1 \leq j \leq m}$ is an extremizer, then so is $(f_j(\cdot-C_jx_0))_{1 \leq j \leq m}$ for any $x_0 \in H.$ 
        \item[(iv)] If $f^{(k)}=(f_j^{(k)})_{1 \leq j \leq m}$ is a sequence of extremizers converging in $L^1$ to $f$, \textit{i.e.} $\lim_{k \rightarrow +\infty}\|f_j^{(k)}-f_j\|_{L^1(H_j)}=0$ for all $1 \leq j \leq m$ and such that $f_j$ is more log-convex than $g_{Q_j}$ for any $j \in I$, more log-concave than $g_{Q_j}$ for any $j \in I^c$, then $f=(f_j)_{1 \le j \le m}$ is also an extremizer.
        \item[(v)] If $f$ and $g$ are extremizers, then $\mathrm{Conv}(f,g)$ is also extremizer.
        \item[(vi)] If $f$ and $g$ are extremizers, then $$\Bigl(f_j\Bigl(\frac{\cdot-C_jx_0}{\sqrt{2}}\Bigl)g_j\Bigl(\frac{\cdot+C_jx_0}{\sqrt{2}}\Bigl)\Bigl)_{1 \leq j \leq m}$$ is also an extremizer for almost every $x_0 \in H.$ In addition, if $f$ and $g$ are bounded, then the latter holds for every $x_0 \in H.$ 
    \end{itemize}
\end{lemma}


\begin{proof}
The first three statements are trivial.
Statement (iv) follows from standard arguments since one can use the definition of $\mathrm{BL}(\textbf{C},p,\textbf{Q},\mathcal{Q})$ and the assumption that the Brascamp-Lieb constant is finite to show that $$\int_H e^{-\pi \langle \mathcal Qx,x\rangle}\prod_{j=1}^m (f_j^{(k)}\circ C_j)^{p_j} \rightarrow_{k \rightarrow +\infty }\int_H e^{-\pi \langle \mathcal Q x,x\rangle}\prod_{j=1}^m(f_j \circ C_j)^{p_j}.$$
Statement (v) is a consequence of \cref{J_Q_increases_under_CONV} and $$\int_{H_i}\mathrm{Conv}(f_i,g_i)=\int_{H_i}f_i \int_{H_i}g_i.$$ Finally for statement (vi), we use \cref{Ext_implies_Gaussianext_eq} and the pairs $(\tilde f_j)_{1 \leq j \leq m}, (g_j)_{1 \leq j \leq m}$ where \hfill \\ $\tilde f_j(x)=f_j(-x)$ for any $1 \leq j \leq m.$ As usual, we can assume that $\int f_i=\int g_i=1$ for any $1 \leq i \leq m$ by (ii). Thus, $$\begin{array}{lll}
\ds \mathrm{BL}(\textbf{C},p,\textbf{Q},I,\mathcal Q)^2=\int_H e^{-\pi \langle \mathcal Qx,x\rangle}\Bigl( \int_H e^{-\pi \langle \mathcal Qy,y\rangle} \prod_{i=1}^m (F_i^{(C_ix)}( C_iy))^{p_i} \, dy \Bigl)\, dx \\
\nm 
\leq \ds \mathrm{BL}(\textbf{C},p,\textbf{Q},I,\mathcal Q)\int_H e^{-\pi\langle \mathcal Qx,x\rangle} \prod_{i=1}^m \Bigl(\int_{H_i} F_i^{(C_ix)} \Bigl)^{p_i} \, dx \leq \mathrm{BL}(\textbf{C},p,\textbf{Q},I,\mathcal Q)^2 \end{array}$$ which means that there must be equality $$\int_H e^{-\pi \langle \mathcal Qx,x\rangle}\prod_{i=1}^m (F_i^{(C_ix)}\circ C_i)^{p_i}\, dx=\mathrm{BL}(\textbf{C},p,\textbf{Q},I,\mathcal Q)\prod_{i=1}^m \Bigl( \int_{H_i}F_i^{(C_ix)}\Bigl)^{p_i}$$ for almost every $x.$ To remove the almost in case $f$ and $g$ are bounded, the argument is the same as in \cite{BCCT}. 
\end{proof}


Specializing the above lemma to centered Gaussian extremizers, we conclude: 

\begin{corollary}
    Let $(\textbf{C},p,\textbf{Q},I,\mathcal Q)$ with $\mathcal Q \geq 0$ be a generalized Brascamp-Lieb datum with $\mathrm{BL}(\textbf{C},p,\textbf{Q},I,\mathcal Q)$ finite. \begin{itemize}
        \item [(i)] If $A^{(n)}$ is a sequence of extremizing Gaussian inputs which converge to a Gaussian input $A=(A_i)_{1 \leq i \leq m}$ with $0<A_i \leq Q_i$ for any $i \in I$, $A_i \ge Q_i$ for any $i \in I^c$, then $A$ is also an extremizer. 
        \item[(ii)] If $A=(A_j)_{1 \leq j \leq m}, A'=(A_j')_{1 \leq j \leq m}$ are Gaussian extremizers, then so is $(2(A_j^{-1}+(A_j')^{-1})^{-1})_{1 \leq j \leq m}.$ 
        \item[(iii)] If $A$ and $A'$ are Gaussian extremizers, then so is $\frac{A+A'}{2}.$ 
    \end{itemize} 
\end{corollary}

\begin{proof}
    (i) is clear by \cref{Closure_prop_extremizers} (iv); (ii) follows from \cref{Closure_prop_extremizers} (v), the computation of the convolution of two Gaussians and the fact that $$0<A_i,A_i' \leq Q_i \Longrightarrow A_i^{-1}+(A_i')^{-1} \geq 2Q_i^{-1} \Longrightarrow 0<2(A_i^{-1}+(A_i')^{-1})^{-1} \leq Q_i \quad \forall i \in I,$$ and similarly, $$A_i,A_i' \ge Q_i \Longrightarrow 0< A_i^{-1}+(A_i')^{-1} \le 2Q_i^{-1} \Longrightarrow (A_i^{-1}+(A_i')^{-1})^{-1} \ge Q_i \quad \forall i \in I^c.$$ Finally, (iii) follows from \cref{Closure_prop_extremizers} (vi).   
\end{proof}

\hfill \\




Now that we have shown that extremizable is equivalent to Gaussian extremizable, one would like to understand what are the necessary and sufficient conditions for Gaussian extremizers/extremizers to exist. 
We start by defining a notion of equivalence between generalized Brascamp-Lieb data. 

\begin{theorem}
    Two Brascamp-Lieb data $(\textbf{C},p,\textbf{Q},I)$ and $(\textbf{C'},p',\textbf{Q'},I')$ are equivalent if $I=I'$ and there exist invertible linear transformations $D: H' \rightarrow H$ and $D_j: H_j' \rightarrow H_j$ (called intertwining transformations) such that $C'_j=D_j^{-1}C_jD$ and $Q_j'=D_j^{*}Q_jD_j$ for all $j$ and $p=p'.$ 
\end{theorem}

The Brascamp-Lieb constants of equivalent data are closely related.

\begin{lemma} \label{Equiv_BL_Constants}
    Suppose that $(\textbf{C},p,\textbf{Q},I), (\textbf{C'},p',\textbf{Q'},I)$ are two equivalent BL data with intertwining maps $D: H' \rightarrow H, D_j: H_j' \rightarrow H_j.$ Then, $$\mathrm{BL}(\textbf{C'},p',\textbf{Q'},I)=\frac{\prod_{j=1}^m \det(D_j)^{p_j}}{\det(D)}\mathrm{BL}(\textbf{C},p,\textbf{Q},I).$$
\end{lemma}

\begin{proof}
    This lemma is a consequence of the following simple computation: $$\begin{array}{lll}
    \mathrm{BL}(\textbf{C'},p',\textbf{Q'}) &=& \ds \sup_{\substack{0<A_j \leq Q_j' \; \forall j \in I \\ A_j \ge Q_j' \; \forall j \in I^c}}\frac{\prod_{j=1}^m \det(A_j)^{p_j/2}}{\det \Bigl(\ds \sum_{j=1}^m p_j(C_j')^*A_jC_j'\Bigl)^{1/2}} \\
    \nm 
    &=& \ds \sup_{\substack{0<A_j \leq Q_j' \; \forall j \in I \\ A_j \ge Q_j' \; \forall j \in I^c}}\frac{\prod_{j=1}^m \det(A_j)^{p_j/2}}{\det \Bigl(D^* \Bigl[\ds \sum_{j=1}^m p_jC_j^*\bigl(D_j^{-*}A_jD_j^{-1} \bigl)C_j \Bigl] D\Bigl)^{1/2}} \\
    \nm 
    &=& \ds \frac{\prod_{j=1}^m \det(D_j)^{p_j}}{\det(D)} \sup_{\substack{0<A_j \leq Q_j' \; \forall j \in I \\ A_j \ge Q_j' \; \forall j \in I^c}} \frac{\prod_{j=1}^m \det(D_j^{-*}A_jD_j^{-1})^{p_j/2}}{\det \Bigl(\ds \sum_{j=1}^m p_jC_j^*\bigl(D_j^{-*}A_jD_j^{-1} \bigl)C_j\Bigl)^{1/2}}
    \\
    \nm 
    &=& \ds \frac{\prod_{j=1}^m \det(D_j)^{p_j}}{\det(D)} \sup_{\substack{0<A_j \leq (D_j)^{-*}Q_j'D_j^{-1} \; \forall j \in I \\ A_j \ge (D_j)^{-*}Q_j'D_j^{-1} \; \forall j \in I^c}} \frac{\prod_{j=1}^m \det(A_j)^{p_j/2}}{\det \Bigl(\ds \sum_{j=1}^m p_jC_j^*A_jC_j\Bigl)^{1/2}} \\
    \nm 
    &=& \ds \frac{\prod_{j=1}^m \det(D_j)^{p_j}}{\det(D)} \mathrm{BL}(\textbf{C},p,\textbf{Q},I).
    \end{array}$$
\end{proof}

\begin{definition}
    A Brascamp-Lieb datum $(\textbf{C},p,\textbf{Q},I)$ is a generalized geometric BL datum if $$C_jC_j^* \leq \mathrm{id}_{H_j}, Q_j \geq \mathrm{id}_{H_j}, (\mathrm{id}_{H_j}-C_jC_j^*)(Q_j-\mathrm{id}_{H_j})=0$$ for any $j \in I$, $$C_jC_j^* \geq \mathrm{id}_{H_j}, Q_j \leq \mathrm{id}_{H_j}, (\mathrm{id}_{H_j}-C_jC_j^*)(Q_j-\mathrm{id}_{H_j})=0$$ for any $j \in I^c$ and \begin{equation} \label{GGeometric_BL_datum} \sum_{j=1}^m p_jC_j^*C_j=\mathrm{id}_H.\end{equation}
\end{definition}



\begin{remark}
    Let $G>0$ be positive definite and let $u \in C^2$ be a positive function that is more log-convex than $g_G$ and $A \geq 0.$ Then, $$\mathrm{div}(A\nabla \log u)=\mathrm{Tr}(A^{1/2}D^2(\log u)A^{1/2}) \geq -2\pi \mathrm{Tr}(A^{1/2}GA^{1/2})=-2\pi \mathrm{Tr}(AG).$$ Similarly, if $u \in C^2$ is a positive function that is more log-concave than $g_G$ and $A \ge 0$, $$-\mathrm{div}(A\nabla \log u) \ge 2\pi \mathrm{Tr}(AG).$$
\end{remark}




\begin{theorem} \label{BL_Geometric_equals_1}
A generalized geometric Brascamp-Lieb datum has BL constant equal to $1.$
\end{theorem}

From the proof of \cref{Theorem_REG_BL}, one may see that $$1 \leq \mathrm{BL}(\textbf{C},p,\textbf{Q},[m])=\mathrm{BL}_g(\textbf{C},p,\textbf{Q},[m]) \leq \mathrm{BL}_{\text{type Q fcts}}(\textbf{C},p,\textbf{Q},[m])=1$$ where the last equality is shown in \cite[Prop 8.9]{BCCT} and the first inequality holds by testing $A_j=\mathrm{id}_{H_j}$ for any $1 \leq j \leq m.$ However, we repeat the proof to be able to treat the mixed log-concavity/log-convexity regularization (\text{i.e.} if $I \neq \{1,\cdots,m\}$) and as it will be a crucial starting point to study the equality cases. 

\hfill \\


\begin{proof}
Let $f_j$ be a Schwarz function that is more log-convex than $g_{Q_j}$ for all $j \in I$ and more log-concave than $g_{Q_j}$ for all $j \in I^c$. This implies that $\nabla \log \tilde f_j(x,t)$ grows at most polynomially in space, locally uniformly on $(1;+\infty)$ by \cite{Vald2} and the assumptions of \cite[Lemma 2.6]{BCCT} are then satisfied.
Let us define the following function:
$$\tilde Q(x,t)=t^{(\sum_{j=1}^m p_j\dim(H_j)-\dim(H))/2}\prod_{j=1}^m \tilde f_j(C_jx,t)^{p_j}$$ with $\tilde f_j: H_j \times [1;+\infty) \rightarrow \mathbb{R}_+$ solution to the heat equation  $$\tilde f_j(x_j,1)=f_j(x_j) \text{ and } \partial_t \tilde f_j(x_j,t)=\frac{1}{4\pi}\Delta \tilde f_j(x_j,t) \text{ for } t>1.$$
Then $\tilde Q(x,1)= \prod_{j=1}^m f_j(C_jx)^{p_j}$ and since $$\tilde f_j(x,t)=\int_{H_j}\frac{1}{(t-1)^{n_j/2}}\exp \Bigl(-\frac{\pi|x-y|^2}{t-1} \Bigl) f_j(y)\, dy=\int_{H_j}\exp(-\pi|z|^2)f_j(x+\sqrt{t-1}z)\, dz,$$ we have that $$\begin{array}{lll}
\ds \int_H \tilde Q(x,t)\, dx &=& \ds t^{\sum_{j=1}^m p_j\dim(H_j)/2}\int_H \prod_{j=1}^m \tilde f_j^{p_j}(\sqrt{t}C_jx,t)\, dx \\
\nm 
&=& \ds \frac{t^{\sum_{j=1}^m p_jn_j/2}}{(t-1)^{\sum_{j=1}^m p_jn_j/2}}\int_H \prod_{j=1}^m \Bigl( \int_{H_j} e^{-\pi \frac{|\sqrt{t} C_jx-y|^2}{t-1}} f_j(y)\, dy \Bigl)^{p_j}\, dx \\
\nm 
&=& \ds \frac{t^{\sum_{j=1}^m p_jn_j/2}}{(t-1)^{\sum_{j=1}^m p_jn_j/2}} \int_H \prod_{j=1}^m \Bigl( e^{-\pi \frac{t}{t-1}p_j|C_jx|^2} \Bigl(\int_{H_j}  e^{\frac{2\sqrt{t}}{t-1}\langle C_jx,y\rangle-\frac{|y|^2}{t-1}}f_j(y)\, dy\Bigl)^{p_j} \Bigl)\, dx\\
\nm 
&& \ds \rightarrow_{t \rightarrow +\infty} \int_H e^{-\pi \langle \sum_{j=1}^m p_jC_j^*C_jx,x\rangle}\prod_{j=1}^m \Bigl(\int_{H_j}f_j\Bigl)^{p_j}\, dx = \prod_{j=1}^m \Bigl( \int f_j \Bigl)^{p_j}.
\end{array}$$ In the last line, dominated convergence was used as well as the fact that $\sum_{j=1}^m p_jC_j^*C_j=\mathrm{id}_H$ since the BL datum is generalized geometric.
The goal is to show that $$\int_H \tilde Q(x,t)\, dx$$ is non-decreasing in $t$ (which suffices by density of Schwarz functions) and the proof employs \cite[Lemma 2.6]{BCCT} (multilinear monotonicity for transport equations).
First, let us show that $\tilde f_j(\cdot,t)$ is more log-convex than $g_{Q_j/t}$ for any $j \in I$ and more log-concave than $g_{Q_j/t}$ for any $j \in I^c.$ Indeed, setting $h_j(x)=e^{\pi \langle Q_jx,x\rangle}f_j(x)$, one has $$\begin{array}{lll}
\tilde f_j(x,t) &=& \ds \int_{H_j}e^{-\pi |z|^2}f_j(x+\sqrt{t-1}z) \, dz \\
\nm 
&=& \ds \int_{H_j}e^{-\pi |z|^2}e^{-\pi \langle Q_j(x+\sqrt{t-1}z),x+\sqrt{t-1}z\rangle}h_j(x+\sqrt{t-1}z)\, dz \\
\nm 
&=& \ds \int_{H_j}e^{-\pi \langle Q_jx,x\rangle}e^{-2\pi \sqrt{t-1}\langle Q_jx,z\rangle}e^{-\pi(t-1)\langle Q_jz,z\rangle}e^{-\pi |z|^2} h_j(x+\sqrt{t-1}z)\, dz.
\end{array}$$ Using the notation of  \cite[Theorem 4.3]{Brascamp_Lieb}, $$\begin{pmatrix}
    D_j^{(1)} \quad D_j^{(2)} \\
    (D_j^{(2)})^* \quad D_j^{(3)}
\end{pmatrix}=\begin{pmatrix}
    \pi Q_j \quad \pi \sqrt{t-1}Q_j \\
    \pi \sqrt{t-1}Q_j \quad \pi(t-1)Q_j+\pi \mathrm{id}_{H_j}
\end{pmatrix},$$ one has $$\begin{array}{lll} D_j &=& D_j^{(1)}-(D_j^{(2)})^*(D_j^{(3)})^{-1}D_j^{(2)} \\ 
&=& \pi Q_j-\pi^2(t-1)Q_j^2 \bigl[\pi(t-1)Q_j+\pi \mathrm{id}_{H_j} \bigl]^{-1} \\
\nm 
&=& \pi Q_j\Bigl( \mathrm{id}_{H_j}-(t-1)Q_j[(t-1)Q_j+\mathrm{id}_{H_j}]^{-1} \Bigl)=\pi Q_j[(t-1)Q_j+\mathrm{id}_{H_j}]^{-1} \\
\nm 
&=& \pi [(t-1)\mathrm{id}_{H_j}+Q_j^{-1}]^{-1} \leq \pi Q_j/t \quad \forall j \in I,
\end{array}$$ since $\mathrm{id}_{H_j} \geq Q_j^{-1} \Longrightarrow (t-1)\mathrm{id}_{H_j}+Q_j^{-1} \geq tQ_j^{-1}$ for any $j \in I.$ 
As $Q_j^{-1} \ge \mathrm{id}_{H_j} \Longrightarrow (t-1)\mathrm{id}_{H_j}+Q_j^{-1} \le tQ_j^{-1}$ for any $j \in I^c$, $$D_j \ge \pi Q_j/t \quad \forall j \in I^c.$$ 
Applying \cite[Theorem 4.3]{Brascamp_Lieb}, it follows that $\tilde f_j(\cdot,t)$ is more log-convex (resp. more log-concave) than $g_{D_j/\pi}$, in particular more log-convex (resp. more log-concave) than $g_{Q_j/t}$ for any $j \in I$ (resp. $j \in I^c$).  



Consider the functions $$u_j=\tilde f_j \circ C_j, \overrightarrow{v_j}=-\frac{1}{4\pi}D_j(\nabla \log \tilde f_j) \circ C_j$$ where $D_j$ is a right-inverse of $C_j$, $$\overrightarrow{v}=\sum_{j=1}^m p_jC_j^*C_jv_j, \alpha=\frac{1}{2}\Bigl(\dim(H)-\sum_{j=1}^m p_j \dim(H_j) \Bigl).$$ As $\partial_t \tilde f_j=\frac{1}{4\pi}\Delta \tilde f_j$, $$\partial_t u_j=\frac{1}{4\pi}(\Delta \tilde f_j)\circ C_j$$ and $$\begin{array}{lll}
    \mathrm{div}(\overrightarrow{v_j}u_j) &=& \ds -\frac{1}{4\pi}\mathrm{div}([D_j(\nabla \log \tilde f_j)\tilde f_j]\circ C_j)  \\ \\
    \nm
    &=& \ds -\frac{1}{4\pi}\mathrm{div}(C_jD_j\nabla \tilde f_j) \circ C_j \\ \\
    \nm 
    &=& \ds -\frac{1}{4\pi}(\Delta \tilde f_j)\circ C_j. \\
    \nm 
\end{array}$$ Thus, $\mathrm{div}(\overrightarrow{v_j}u_j)+\partial_t u_j=0$ for any $1 \leq j \leq m.$ Moreover, since $C_j(C_j^*C_j-\mathrm{id}_{H_j})D_j=C_jC_j^*-\mathrm{id}_{H_j}$, $$\begin{array}{lll}
\ds \mathrm{div}\Bigl(\overrightarrow{v}-\sum_{j=1}^m p_j\overrightarrow{v_j}\Bigl) &=& \ds \sum_{j=1}^m p_j \mathrm{div}((C_j^*C_j-\mathrm{id}_{H_j})\overrightarrow{v_j}) \\
\nm 
&=& \ds  -\frac{1}{4\pi}\sum_{j=1}^m p_j\mathrm{div}((C_j^*C_j-\mathrm{id}_{H_j})D_j(\nabla \log \tilde f_j)\circ C_j) \\
\nm 
&=& \ds -\frac{1}{4\pi}\sum_{j=1}^m p_j \mathrm{div}(C_j(C_j^*C_j-\mathrm{id}_{H_j})D_j(\nabla \log \tilde f_j))\circ C_j \\
\nm 
&=& \ds \frac{1}{4\pi}\sum_{j=1}^m p_j\mathrm{div}((\mathrm{id}_{H_j}-C_jC_j^*)\nabla \log \tilde f_j) \circ C_j.
\end{array}$$  By the previous remark, $$
\ds \mathrm{div}((\mathrm{id}_{H_j}-C_jC_j^*)\nabla \log \tilde f_j) \geq \ds -\frac{2\pi}{t}\mathrm{Tr}((\mathrm{id}_{H_j}-C_jC_j^*)Q_j)=-\frac{2\pi}{t}\mathrm{Tr}(\mathrm{id}_{H_j}-C_jC_j^*) 
$$ which implies by summing over $j=1,\cdots,m$, that $$\begin{array}{lll}
\ds \frac{1}{4\pi}\sum_{j=1}^m p_j\mathrm{div}((\mathrm{id}_{H_j}-C_jC_j^*)\nabla \log \tilde f_j) &\geq& \ds \frac{1}{2t}\sum_{j=1}^m p_j\mathrm{Tr}(C_j^*C_j)-\frac{1}{2t}\sum_{j=1}^m p_j\dim(H_j) \\
\nm
&=& \ds \frac{\alpha}{t}.
\end{array}$$ Finally, $$\nabla \log u_j=C_j^*(\nabla \log \tilde f_j)\circ C_j \Longrightarrow -C_j^*C_j\overrightarrow{v_j}=\frac{1}{4\pi}C_j^*C_jD_j(\nabla \log \tilde f_j)\circ C_j=\frac{1}{4\pi}\nabla \log u_j$$ which gives that $$\sum_{j=1}^m p_j\langle \overrightarrow{v}-\overrightarrow{v_j},\nabla \log u_j\rangle_H =4\pi \sum_{j=1}^m p_j \langle C_j^*C_j(\overrightarrow{v}-\overrightarrow{v_j}),-\overrightarrow{v_j}\rangle_H.$$  As $\sum_{j=1}^m p_jC_j^*C_j(\overrightarrow{v}-\overrightarrow{v_j})=\overrightarrow{v}-\sum_{j=1}^m p_jC_j^*C_j\overrightarrow{v_j}=0$, one has $$\begin{array}{lll}
\ds \sum_{j=1}^m p_j \langle \overrightarrow{v}-\overrightarrow{v_j},\nabla \log u_j\rangle_H &=& \ds 4\pi \sum_{j=1}^m p_j \langle C_j^*C_j(\overrightarrow{v}-\overrightarrow{v_j}),\overrightarrow{v}-\overrightarrow{v_j}) \geq 0 \\
\nm 
&=& \ds 4\pi \Bigl[ \sum_{j=1}^m p_j \langle C_j\overrightarrow{v_j},C_j\overrightarrow{v_j}\rangle-\langle \overrightarrow{v},\overrightarrow{v} \rangle \Bigl] \\
\nm 
&=& 4\pi \langle (I-P)A(x),A(x) \rangle \geq 0
\end{array}$$ where $$A(x)=\begin{pmatrix}
    p_1^{1/2}C_1\overrightarrow{v_1} \\
    \vdots \\
    p_m^{1/2}C_m\overrightarrow{v_m}
\end{pmatrix}, \; T=\begin{pmatrix}
    \cdots p_1^{1/2}C_1 \cdots \\
    \vdots \\
    \cdots p_m^{1/2}C_m \cdots
\end{pmatrix}$$ and $P=TT^*=T(T^*T)^{-1}T^*$ is  a projection (recall by assumption, $T^*T=\sum_{j=1}^m p_jC_j^*C_j=\mathrm{id}_{H}$). The conditions to apply Lemma $2.6$ are all fulfilled and hence we get the desired claim. 
\end{proof}


\begin{remark} \label{Important_rk_equality_cases}
There would be equality $$\int_H \prod_{j=1}^m (f_j \circ C_j)^{p_j} \, = \mathrm{BL}(\textbf{C},p,\textbf{Q})\prod_{j=1}^m \Bigl( \int_{H_j}f_j\Bigl)^{p_j}$$ if and only if $$\mathrm{Tr}((\mathrm{id}_{H_j}-C_jC_j^*)(D^2\log \tilde f_j+\frac{2\pi}{t}Q_j))=0 \quad \forall 1 \le j \le m, $$ as the trace of a product of two positive semi-definite (or two negative semi-definite) matrices and $$\langle (I-P)A(x),A(x)\rangle=0.$$ 
\end{remark}

\hfill \\


One can give a satisfactory algebraic characterization of Gaussian extremizable Brascamp-Lieb data using the previously introduced definitions. 

\begin{proposition} \label{Prop_Equiv_BL_constant}
    Let $(\textbf{C},p,\textbf{Q},I)$ be a Brascamp-Lieb datum with $p_j>0$ for all $j$, $B=(B_j)_{1 \leq j \leq m}$ such that $0<B_j \leq Q_j$ for any $j \in I$, $B_j \ge Q_j$ for any $j \in I^c.$ Denote by $M=\sum_{j=1}^m p_jC_j^* B_jC_j.$ Then the following assertions are equivalent: 
    \begin{enumerate}
        \item [(a)] $B$ is a global extremizer, \textit{i.e.} $$\mathrm{BL}(\textbf{C},p,\textbf{Q},I)=\frac{\prod_{j=1}^m \det(B_j)^{p_j/2}}{\det(M)^{1/2}},$$
        \item[(b)] $B$ is a local extremizer,
        \item[(c)] $M$ is invertible with $$B_j^{-1} \geq C_jM^{-1}C_j^* \quad \forall j \in I,\, B_j^{-1} \le C_jM^{-1}C_j^* \quad \forall j \in I^c$$ and $$ (B_j^{-1}-C_jM^{-1}C_j^*)(Q_j-B_j)=0 \quad \forall 1 \leq j \leq m,$$
        \item[(d)] $(\textbf{C},p,\textbf{Q},I)$ is equivalent to a generalized geometric BL datum $(\textbf{C'},p',\textbf{Q'},I)$ with intertwining operators $D=M^{-1/2}, D_j=B_j^{-1/2}.$
    \end{enumerate}
\end{proposition}


\begin{proof}
    The implication from $(a)$ to $(b)$ is trivial. Let us verify that $(b)$ implies $(c)$. 
    By assumption, $B$ is a local extremizer for the function $$f(A_1,\cdots,A_m)=\sum_{j=1}^m p_j \log \det A_j -\log \det \Bigl(\sum_{j=1}^m p_jC_j^*A_jC_j \Bigl)$$ under the constraint $0<A_j \leq Q_j$ for any $j \in I$, $A_j \ge Q_j$ for any $j \in I^c.$ Fix some $j \in I$ and denote by $V_j$ the kernel of the positive semi-definite operator $Q_j-B_j.$ We also consider $\textit{i}_j: V_j \rightarrow H_j$ the inclusion map and $G_j$ a self-adjoint operator that is negative semi-definite on $V_j$ (in other words,  $\textit{i}_j^*G_j\textit{i}_j \leq 0$). Hence, $0<B_j+\epsilon G_j \leq Q_j$ for all $\epsilon>0$ sufficiently small and as $B$ is a local extremizer, $$\limsup_{\epsilon \downarrow 0}\frac{f(B_1,\cdots,B_{j-1},B_j+\epsilon G_j,B_{j+1},\cdots,B_m)-f(B_1,\cdots,B_m)}{\epsilon} \leq 0$$ which gives that $$\mathrm{Tr}((B_j^{-1}-C_jM^{-1}C_j^*)G_j) \leq 0$$ whenever $G_j$ is negative semi-definite on $V_j.$ \hfill \\ In particular, as the previous inequality holds whenever $G_j$ is negative semi-definite on $H_j$, one obtains that $$B_j^{-1} \geq C_jM^{-1}C_j^*.$$ Furthermore, $$\mathrm{Tr}((B_j^{-1}-C_jM^{-1}C_j^*)G_j)=0 \text{ if } \textit{i}_j^*G_j\textit{i}_j=0, G_j^*=G_j.$$ \hfill \\ By positive semi-definiteness of $B_j^{-1}-C_jM^{-1}C_j^*$, one deduces that it must be equal to $\textit{i}_jN_j\textit{i}_j^*$ for a self-adjoint transformation $N_j: V_j \rightarrow V_j.$  But $(Q_j-B_j)\textit{i}_j=0$ by definition of $V_j$ which implies that $\textit{i}_j^*(Q_j-B_j)=0$ and so $$(B_j^{-1}-C_jM^{-1}C_j^*)(Q_j-B_j)=0.$$ 
    Similarly, we fix $j \in I^c$ and denote again by $V_j$ the kernel of the positive semi-definite operator $B_j-Q_j.$ Considering $G_j$ this time to be a self-adjoint operator that is positive semi-definite on $V_j$, we may observe that $B_j+\epsilon G_j \ge Q_j$ for all $\epsilon>0$ and as $B$ is a local extremizer, $$\limsup_{\epsilon \downarrow 0}\frac{f(B_1,\cdots,B_{j-1},B_j+\epsilon G_j,B_{j+1},\cdots,B_m)-f(B_1,\cdots,B_m)}{\epsilon} \leq 0$$ which gives that $$\mathrm{Tr}((B_j^{-1}-C_jM^{-1}C_j^*)G_j) \leq 0$$ whenever $G_j$ is positive semi-definite on $V_j.$ This implies in particular that $$B_j^{-1} \le C_jM^{-1}C_j^*$$ and by the same argument as previously, one also derives that $$(B_j^{-1}-C_jM^{-1}C_j^*)(B_j-Q_j)=0.$$
    
    \hfill \\ We prove that $(c)$ implies $(d)$. Let us define $D=M^{-1/2}, D_j=B_j^{-1/2}, Q_j'=B_j^{-1/2}Q_jB_j^{-1/2}$ and $C_j'=B_j^{1/2}C_jM^{-1/2}$ for any $1 \leq j \leq m.$ We claim that $(\textbf{C'},p,\textbf{Q'})$ is a generalized geometric BL datum. Indeed, as $B_j \leq Q_j,$ one has $Q_j' \geq \mathrm{id}_{H_j}$ and $$C_j'(C_j')^*=B_j^{1/2}C_jM^{-1}C_j^*B_j^{1/2} \leq B_j^{1/2}B_j^{-1}B_j^{1/2}=\mathrm{id}_{H_j}.$$ Moreover, $$\sum_{j=1}^m p_j(C_j')^*C_j'=M^{-1/2}\Bigl(\sum_{j=1}^m p_jC_j^*B_jC_j\Bigl)M^{-1/2}=\mathrm{id}_H$$ by definition of $M$ and $$(\mathrm{id}_{H_j}-C_j'(C_j')^*)(Q_j'-\mathrm{id}_{H_j})=B_j^{1/2}(B_j^{-1}-C_jM^{-1}C_j^*)B_j^{1/2}B_j^{-1/2}(Q_j-B_j)B_j^{-1/2}=0$$ by the second assumption of $(c).$ \\ Finally, $(d)$ implies $(a)$ as a consequence of \cref{Equiv_BL_Constants,BL_Geometric_equals_1}. 
\end{proof}

\hfill \\



The dual regularized Brascamp-Lieb inequality \cref{Reg_Dual_BL_Log_Convexity}, together with \cref{Prop_Equiv_BL_constant} for $I=\emptyset$, allows us to extend a result of \cite[Section 6.2]{Nakamura}:

\begin{corollary}
Let $p_1,p_2 \in (0;1)$ be such that $p_1+p_2<1.$ Then, for $\mathbb{R}^n \ni Q_1,Q_2>0$, if $f_j$ is more log-convex than $g_{Q_j^{-1}}$ for $j=1,2$, $$\Bigl(\int_{\mathbb{R}^n}f_1\Bigl)^{p_1}\Bigl( \int_{\mathbb{R}^n}f_2\Bigl)^{p_2} \leq  \Bigl( \frac{\det(Q_1)^{p_1}\det(Q_2)^{p_2}}{\det(p_1Q_1+p_2Q_2)}\Bigl)^{1/2}\int_{\mathbb{R}^n}\esssup_{x=p_1x_1+p_2x_2}f_1(x_1)^{p_1}f_2(x_2)^{p_2}\, dx$$ if and only if $$Q_j \geq p_1Q_1+p_2Q_2 \text{ for } j=1,2.$$ 
\end{corollary}

\bigskip

\begin{example}[Some examples on Gaussian extremizability]
\hfill \\

\begin{enumerate} 

    \item [(i)]  Consider $n=4, n_i=1, q_i=1$ for all $1 \leq i \leq 5.$ Moreover, take $C_1=(1,-1,0,0), C_2=(0,1,0,0), C_3=(1,0,0,0), C_4=(0,0,1,0), C_5=(0,0,0,1)$ and $p_1=p_4=p_5=2,p_2=p_3=\frac{1}{2}.$ 
    Then, one can check that $$\sum_{i=1}^5 p_i\dim(C_iV)=\sum_{V \not \subset \ker(C_i)}p_i \geq \dim(V) \text{ for any linear subspace } V \subset \mathbb{R}^4.$$ 
    Notice that 
    $$\sum_{i=1}^5 p_ib_iC_i^*C_i=\begin{pmatrix}
     2b_1+\frac{1}{2}b_3 & -2b_1 & 0 & 0 \\
    -2b_1 & 2b_1+\frac{1}{2}b_2 & 0 & 0 \\
        0 & 0 & 2b_4 & 0 \\
        0 & 0 & 0 & 2b_5 
    \end{pmatrix}$$ 
    which implies that $$f(b_1,b_2,b_3,b_4,b_5)=\frac{b_1^2b_2^{1/2}b_3^{1/2}b_4^2b_5^2}{(b_1b_2+b_1b_3+\frac{1}{4}b_2b_3)4b_4b_5}=\frac{b_1b_4b_5}{4}\frac{b_2^{1/2}b_3^{1/2}}{b_2+b_3+\frac{1}{4b_1}b_2b_3}.$$ To determine $$\sup_{0<b_i \leq 1}f(b_1,b_2,b_3,b_4,b_5),$$ it is clear that we need to take $b_1=b_4=b_5=1$ so we are left with the computation of $$\sup_{0<b_2,b_3 \leq 1}\frac{b_2^{1/2}b_3^{1/2}}{b_2+b_3+\frac{1}{4}b_2b_3}.$$ Since $b_2^{1/2}b_3^{1/2} \leq \frac{1}{2}(b_2+b_3)$, the supremum is less than $\frac{1}{2}.$ Moreover, if we take $b_2=b_3=\epsilon$, then $$\frac{b_2^{1/2}b_3^{1/2}}{b_2+b_3+\frac{1}{4}b_2b_3}=\frac{\epsilon}{2\epsilon+\frac{1}{4}\epsilon^2}=\frac{1}{2+\frac{1}{4}\epsilon} \rightarrow_{\epsilon \rightarrow 0} \frac{1}{2}.$$ So the supremum is $1/2$ but can never be attained. This shows that not every non-degenerate Brascamp-Lieb data is extremizable. 
    
    \hfill \\
    
    
    \item[(ii)] Let us instead consider a Gaussian extremizable Brascamp-Lieb datum. 
    Set $n=2=n_1,n_2=1, Q_1=I_2, Q_2=1$ and $C_1=\begin{pmatrix}
        1 \; 1 \\
        0 \; 1
    \end{pmatrix}, C_2=(1,2), p_1=1,p_2=2.$ This datum satisfies the desired conditions for finiteness with a critical subspace $V=\mathrm{span}\{(2,-1)\}=\ker(C_2)$. Writing $B_1$ as $\begin{pmatrix}
        b_1^{(1)} \; b_1^{(2)} \\
        b_1^{(2)} \; b_1^{(3)}
    \end{pmatrix}$, we find that $$p_1C_1^*B_1C_1+p_2C_2^*B_2C_2=\begin{pmatrix}
        b_1^{(1)}+2b_2 \quad b_1^{(1)}+b_1^{(2)}+4b_2 \\
        b_1^{(1)}+b_1^{(2)}+4b_2 \quad b_1^{(1)}+2b_1^{(2)}+b_1^{(3)}+8b_2
    \end{pmatrix}$$ and computing the determinant gives $$b_1^{(1)}b_1^{(3)}-(b_1^{(2)})^2+2b_2(b_1^{(1)}+b_1^{(3)}-2b_1^{(2)}).$$ 
    So, we need to maximize $$\frac{\det(B)b_2^2}{\det(B)+2b_2\langle B(1,-1)^T,(1,-1)^T\rangle}=\frac{b_2^2}{1+2b_2\langle \frac{B}{\det(B)}(1,-1)^T,(1,-1)^T\rangle}$$ or in other words, to minimize $\langle \frac{B}{\det(B)}(1,-1)^T,(1,-1)^T\rangle.$
Notice that any $0<B \leq I_2$ can be written as $$B=U\mathrm{diag}(\lambda_1,\lambda_2)U^T, 0<\lambda_1 \leq \lambda_2 \leq 1, U \in O(n)$$ and by Rayleigh's principle, $$\Bigl \langle \frac{B}{\det(B)}\frac{(1,-1)^T}{\sqrt 2},\frac{(1,-1)^T}{\sqrt 2}\Bigl \rangle \geq \frac{\lambda_1}{\lambda_1 \lambda_2}=\frac{1}{\lambda_2} \geq 1.$$
    
    Taking $$B_\epsilon=\begin{pmatrix}
        1 \; 1 \\
        -1 \; 1
    \end{pmatrix} \begin{pmatrix}
        \epsilon \; 0 \\
        0 \; 1
    \end{pmatrix} \begin{pmatrix}
        1 \; 1 \\
        -1 \; 1
    \end{pmatrix}^{-1}=\frac{1}{2}\begin{pmatrix}
        1+\epsilon \quad 1-\epsilon \\
        1-\epsilon \quad 1+\epsilon
    \end{pmatrix},$$ we see that $\det(B_\epsilon)=\epsilon, B_\epsilon(1,-1)^T=\epsilon(1,-1)^T$ which gives $$\frac{b_2^2}{1+4b_2}$$ maximized at $b_2=1$. So, the supremum is equal to $\frac{1}{5}$ and is achieved at the family $(B_\epsilon,1)$ for any $0<\epsilon<1$. 
\end{enumerate}
\end{example}

\hfill \\



\bigskip




\section{Optimizers of Regularized Brascamp-Lieb Inequalities}

\hfill \\

We would like to turn to the study of equality cases for Gaussian extremizable Brascamp-Lieb data. Indeed, investigating equality cases is only meaningful when the datum $(\textbf{C},p,\textbf{Q})$ has a Brascamp-Lieb constant $\mathrm{BL}(\textbf{C},p,\textbf{Q})$ that is extremizable. But, by \cref{Ext_implies_Gaussian_ext}, this implies that the Brascamp-Lieb constant must be Gaussian extremizable. Moreover, by \cref{Prop_Equiv_BL_constant}, one may first reduce to considering the general extremizers of a generalized geometric BL. 
For notational simplicity, we will focus on the case $I=\{1,\cdots,m\}$ but the next arguments allow to conclude that the extremizers for general $I \subset \{1,\cdots,m\}$ will have the same form as in the case where all functions are more log-convex than some fixed Gaussians.

\hfill \\

Suppose that the inputs $(f_j)_{1 \leq j \leq m}$ are Schwarz functions such that $f_j$ is more log-convex than $g_{Q_j}$ for any $1 \leq j \leq m$ and attain equality for the generalized geometric Brascamp-Lieb datum $(\textbf{C},p,\textbf{Q})$, \textit{i.e.} $$\int_H \prod_{j=1}^m (f_j \circ C_j)^{p_j} =\prod_{j=1}^m \Bigl (\int_{H_j}f_j\Bigl)^{p_j}.$$ 
\hfill \\
Fix some time $t>1.$ Let us first derive from the second condition obtained in \cref{Important_rk_equality_cases} an equality between the gradients of $\log \tilde Q$ and $\log \tilde f_j$, which will be fundamental for the rest of the argument. 
Since $\langle (I-P)A(x),A(x)\rangle=0$ with $P=TT^*$, one has $$A(x)=T\beta(x)$$ for almost every $x \in  H$ where $\beta: H \rightarrow H.$ This implies that $p_j^{1/2}C_j \overrightarrow{v_j}=p_j^{1/2}C_j\beta(x)$ for any $1 \leq j \leq m$ by equality of the rows. In other words, $$-\frac{1}{4\pi}(\nabla \log \tilde f_j)(C_jx)=C_j\beta(x) \quad \forall 1 \leq j \leq m.$$ This shows that $\beta$ is smooth and $A(x)=T\beta(x)$ must hold for all $x \in H.$ 


Then, as $$\log \tilde Q(x)=\frac{\sum_{j=1}^m p_j\dim(H_j)-\dim(H)}{2}\log(t)+\sum_{j=1}^m p_j\log (\tilde f_j(C_jx)),$$ one has, using $\sum_{j=1}^m p_jC_j^*C_j=\mathrm{id}_{H_j}$, $$\nabla \log \tilde Q(x)=\sum_{j=1}^m p_j C_j^*(\nabla \log \tilde f_j)(C_jx)=-4\pi \sum_{j=1}^m p_jC_j^*C_j\beta(x)=-4\pi \beta(x).$$ Thus, \begin{equation} \label{Fundamental_eq_equality_cases} C_j\nabla \log \tilde Q(x)=(\nabla \log \tilde f_j)(C_jx)\end{equation} for all $x \in H.$  
\hfill \\








\hfill \\


It is now time to extract as much information as possible from the first condition of \cref{Important_rk_equality_cases}. 
Setting $B_j= \bigl[ (t-1)\mathrm{id}_{H_j}+Q_j^{-1} \bigl]^{-1},$ note that $$\begin{array}{lll}
\ds \frac{Q_j}{t}-B_j &=& \ds \frac{Q_j}{t}- \bigl[ (t-1)\mathrm{id}_{H_j}+Q_j^{-1} \bigl]^{-1} \\ \\
\nm 
&=& \ds \Bigl(\frac{Q_j}{t} \bigl[ (t-1)\mathrm{id}_{H_j}+Q_j^{-1} \bigl]-\mathrm{id}_{H_j}\Bigl)\bigl[ (t-1)\mathrm{id}_{H_j}+Q_j^{-1} \bigl]^{-1} \\ \\
\nm 
&=& \ds \frac{t-1}{t}(Q_j-\mathrm{id}_{H_j})\bigl[ (t-1)\mathrm{id}_{H_j}+Q_j^{-1} \bigl]^{-1}.
\end{array}$$ By assumption, $$\mathrm{Tr}((\mathrm{id}_{H_j}-C_jC_j^*)(D^2\log \tilde f_j+\frac{2\pi }{t}Q_j))=0$$ and $(\mathrm{id}_{H_j}-C_jC_j^*)(Q_j-\mathrm{id}_{H_j})=0$, which shows that $$\mathrm{Tr}((\mathrm{id}_{H_j}-C_jC_j^*)(D^2\log \tilde f_j+2\pi B_j))=0.$$ Further, since $\mathrm{id}_{H_j}-C_jC_j^*$ and $ D^2\log \tilde f_j+2\pi B_j$ are both positive semi-definite, the previous trace equality actually means that $$(\mathrm{id}_{H_j}-C_jC_j^*)^{1/2}(D^2 \log \tilde f_j+2\pi B_j)(\mathrm{id}_{H_j}-C_jC_j^*)^{1/2}=0.$$ \hfill \\ For notational convenience, set $A_j=\mathrm{id}_{H_j}-C_jC_j^*<\mathrm{id}_{H_j}.$ As $A_j,Q_j$ commute, they are simultaneously diagonalizable or in other words, there exists an orthonormal basis of $H_j$, $\{v_1^{(j)},\cdots,v_{n_j}^{(j)}\}$, such that $$A_jv_i^{(j)}=a_i^{(j)}v_i^{(j)}, Q_jv_i^{(j)}=q_i^{(j)}v_i^{(j)} \quad \forall 1 \leq i \leq n_j$$ with $1>a_{1}^{(j)} \geq \cdots \geq a_{n_j}^{(j)} \geq 0,\;  q_i^{(j)} \geq 1$ for any $1 \leq i \leq n_j$ (we can only ask for one set of eigenvalues to be ordered).  Set $0 \leq k_j \leq n_j$ such that $a_i^{(j)}>0$ for any $1 \leq i \leq k_j, a_i^{(j)}=0$ for all $k_j+1 \leq i \leq n_j.$ Denoting $V_j \in O(n_j)$ the change of basis matrix, $$Q_j=V_j \,\mathrm{diag}(q_1^{(j)},\cdots,q_{n_j}^{(j)}) V_j^T \Longrightarrow B_j=V_j\, \mathrm{diag}\Bigl(\frac{1}{t-1+\frac{1}{q_1^{(j)}}},\cdots,\frac{1}{t-1+\frac{1}{q_{n_j}^{(j)}}}\Bigl)V_j^T.$$ Moreover, as $(\mathrm{id}_{H_j}-C_jC_j^*)(Q_j-\mathrm{id}_{H_j})=0$, $q_i^{(j)}=1$ for any $1 \leq i \leq k_j.$ In particular, the $k_j$-th first eigenvalues of $B_j$ are $t^{-1}.$ Applying the eigenvectors of $A_j$ that are not associated to a zero eigenvalue to the equality $A_j^{1/2}(D^2 \log \tilde f_j+2\pi B_j)A_j^{1/2}=0$, one obtains that $$A_j^{1/2}((D^2 \log \tilde f_j )v_i^{(j)}+2\pi B_jv_i^{(j)})=0 \quad \forall 1 \leq i \leq k_j$$ which implies that $$(D^2 \log \tilde f_j)v_i^{(j)}=-\frac{2\pi}{t}v_i^{(j)}+\sum_{l=k_{j}+1}^{n_j}\alpha_l v_l^{(j)}.$$ This means that decomposing $H_j$ as $\mathrm{Ran}(\mathrm{id}_{H_j}-C_jC_j^*)=\mathrm{span}\{v_1^{(j)},\cdots,v_{k_j}^{(j)}\}$ in direct sum with $\ker(\mathrm{id}_{H_j}-C_jC_j^*)=\mathrm{span}\{v_{k_j+1}^{(j)},\cdots,v_{n_j}^{(j)}\}$, one gets that $$ D^2(\log \tilde f_j)=V_j\left(\begin{array}{ c | c }
    \mathrm{diag}(-\frac{2\pi}{t},\cdots,-\frac{2\pi}{t}) & \cdot_j \\
    \hline
    \cdot_j & \bullet_j 
  \end{array}\right)V_j^T.$$ This uses completely the first condition of \cref{Important_rk_equality_cases} in the sense that if $D^2(\log \tilde f_j)$ is of this form, then it satisfies the trace equality. 
  Recall the two remaining conditions that are available to us: \begin{equation} \label{Remaining_conditions_eq} (\nabla \log \tilde f_j)(C_jx)=C_j\nabla \log \tilde Q(x) \text{ and } D^2\log \tilde f_j \geq -2\pi B_j \geq -2\pi \frac{Q_j}{t}.\end{equation}

\hfill \\

It will be useful to decompose $H$ as a direct sum containing $\ker(C_j)$. More precisely, one may write $$H=\ker(C_j) \oplus \mathrm{Ran}(C_j^*)=\ker(C_j) \oplus C_j^*\mathrm{Ran}(\mathrm{id}_{H_j}-C_jC_j^*)\oplus C_j^* \ker(\mathrm{id}_{H_j}-C_jC_j^*).$$ Indeed, $C_jx=\sum_{i=1}^{n_j}\beta_i v_i^{(j)}, C_j(C_j^*v_i^{(j)})=(1-a_i^{(j)})v_i^{(j)}$ which implies that $$C_j\Bigl(x-\sum_{i=1}^{n_j}\frac{\beta_i}{1-a_i^{(j)}}C_j^*v_i^{(j)}\Bigl)=0.$$

\hfill \\

Unless $Q_j>\mathrm{id}_{H_j}$ (in which case $C_jC_j^*=\mathrm{id}_{H_j}$ and one can reduce to the analysis of Valdimarsson in \cite{Vald2}), we have $$0 \leq D^2(\log \tilde f_j)+2\pi \frac{Q_j}{t}=V_j\left(\begin{array}{ c | c }
    \mathrm{diag}(0,\cdots,0) & \cdot_j \\
    \hline
    \cdot_j & \bullet_j+\mathrm{diag}(\frac{2\pi}{t}q_{k_j+1},\cdots,\frac{2\pi}{t}q_{n_j})
  \end{array}\right)V_j^{T},$$ which means that $\cdot_j=0, \bullet_j \geq -\mathrm{diag}(\frac{2\pi}{t}q_{k_j+1},\cdots,\frac{2\pi}{t}q_{n_j}).$ 


\hfill \\


By differentiating the first equality in \cref{Remaining_conditions_eq}, one gets $$D^2(\log \tilde f_j)(C_jx)C_j=C_jD^2\log \tilde Q(x)$$ and in particular, for $k_j+1 \leq i \leq n_j$, $$D^2\log \tilde f_j(C_jx)v_i^{(j)}=D^2\log \tilde f_j(C_jx)C_jC_j^*v_i^{(j)}=C_jD^2\log \tilde Q(x)C_j^*v_i^{(j)}$$ and so, $$D^2 (\log \tilde f_j)(C_jx)=V_j\left(\begin{array}{ c | c }
    \mathrm{diag}(-\frac{2\pi}{t},\cdots,-\frac{2\pi}{t}) & 0 \\
    \hline
    0 & C_jD^2(\log \tilde Q)(x)C_j^* \bigl|_{\ker(A_j)}
  \end{array}\right)V_j^{T}.$$ 



Thus, further investigation of $D^2 (\log \tilde Q)$ is needed.  

\hfill \\

Notice that for $c_j \in \mathrm{Ran}(C_j^*), c_j=C_j^*e_j$, $$\langle \nabla \log\tilde Q(x),c_j\rangle=-4\pi\langle \beta(x), C_j^*e_j\rangle=\langle \nabla \log \tilde f_j(C_jx),e_j\rangle.$$ Differentiating in a direction of $\ker(C_j)$, one obtains that for $a_j \in \ker(C_j)$, $$\langle D^2 (\log \tilde Q)(x)a_j,C_j^*e_j \rangle=0.$$ Moreover, for any $1 \leq i \leq k_j$, $$\begin{array}{lll}
C_jD^2(\log \tilde Q)(x)C_j^*v_i^{(j)} &=& D^2(\log \tilde f_j)(C_jx)C_jC_j^*v_i^{(j)} \\
\nm 
&=& \ds (1-a_i^{(j)})D^2(\log \tilde f_j)(C_jx)v_i^{(j)}=-\frac{2\pi}{t}(1-a_i^{(j)})v_i^{(j)}
\end{array}$$ which implies that $$D^2(\log \tilde Q)(x)C_j^*v_i^{(j)}=-\frac{2\pi}{t}C_j^*v_i^{(j)} \text{ since } D^2(\log \tilde Q)(x)C_j^*v_i^{(j)} \perp \ker(C_j).$$ Hence, $$D^2(\log \tilde Q)(x)=\left(\begin{array}{ c | c }
    \ker(C_j) & 0 \\
    \hline
    0 & \left. \begin{array}{ c | c }
    \mathrm{diag}(-\frac{2\pi}{t},\cdots,-\frac{2\pi}{t}) & 0  \\
    \hline
    0 & W_j
  \end{array} \right.
  \end{array}\right)$$
where $W_j=C_j^*\ker(\mathrm{id}_{H_j}-C_jC_j^*)$ and the second block is with respect to the basis $C_j^*\mathrm{Ran}(\mathrm{id}_{H_j}-C_jC_j^*) =: \tilde W_j.$ 

\bigskip


Set $$K=\bigcap_{j=1}^m \bigl( \ker(C_j) \oplus C_j^*\ker(\mathrm{id}_{H_j}-C_jC_j^*) \bigl).$$
As we already know the form of $D^2(\log \tilde Q)$ on $C_j^*\mathrm{Ran}(\mathrm{id}_{H_j}-C_jC_j^*)$ for any $1 \leq j \leq m,$ what is left to determine is the form of $D^2(\log \tilde Q)$ on $K.$ In fact, $K$ might not be critical so we will further decompose $D^2(\log \tilde Q)$ with respect to the largest critical subspace that is invariant under $C_iC_i^*$ (it will be seen why this invariance property is crucial for the argument to hold) and its orthogonal complement. 

\hfill \\

We introduce the following notation: if $A \subset B \subset H$, then the perpendicular of $A$ with respect to $B$ will be denoted by $$A^{\perp,B}=A^\perp \cap B.$$ 




\hfill \\

Let us remark that $$C_j^*\ker(\mathrm{id}_{H_j}-C_jC_j^*)=\ker(\mathrm{id}_H-C_j^*C_j) \text{ for any } 1 \leq j \leq m.$$ On the one hand, $x \in C_j^*\ker(\mathrm{id}_{H_j}-C_jC_j^*) \Longrightarrow x=C_j^*z $ where $z \in \ker(\mathrm{id}_{H_j}-C_jC_j^*)$ and $$(C_j^*C_j)x=C_j^*(C_jC_j^*z)=C_j^*z=x$$ so $x \in \ker(\mathrm{id}_H-C_jC_j^*).$  On the other hand, $x \in \ker(\mathrm{id}_H-C_j^*C_j)$ implies that $ x=C_j^*(C_jx)$. But $$C_jC_j^*(C_jx)=C_jx \Longrightarrow C_jx \in \ker(\mathrm{id}_{H_j}-C_jC_j^*).$$  





\hfill \\

To have a better grasp of the strategy of the proof, let us compute two examples where the behavior of the data with respect to $K$ is different. 

\begin{example}
First, we look at an example where $C_i^*C_iK \subset K$ for all $1 \leq i \leq m.$ 
    Let us take $n=2,n_1=2,n_2=1$ and $Q_1=\begin{pmatrix}
        1 \; 0 \\
        0 \; \lambda
    \end{pmatrix}, Q_2=1$ where $\lambda>1,$ $p_1=p_2=1.$ Finally, we set $$C_1=\begin{pmatrix}
        \frac{1}{\sqrt{2}} \; 0 \\
        0 \; 1 
    \end{pmatrix},C_2=\Bigl(\frac{1}{\sqrt 2},0 \Bigl).$$  Then, $$C_1C_1^* \leq I_2, (I_2-C_1C_1^*)(Q_1-I_2)=\begin{pmatrix}
        \frac{1}{2} \; 0\\
        0 \; 0
    \end{pmatrix} \begin{pmatrix}
        0 \; 0 \\
        0 \; \lambda-1
    \end{pmatrix}=0.$$ Moreover, $C_2C_2^*=\frac{1}{2}$ and $p_1C_1^*C_1+p_2C_2^*C_2=\begin{pmatrix}
        \frac{1}{2} \; 0 \\
        0 \; 1
    \end{pmatrix}+\begin{pmatrix}
        \frac{1}{2} \; 0 \\
        0 \; 0
    \end{pmatrix}=I_2.$ Hence, it is a generalized geometric BL and in particular, its BL constant is equal to $1.$ Let us first find the Gaussian extremizers: if $B_1=\begin{pmatrix}
        b_1^{(1)} \; b_1^{(2)} \\
        b_1^{(2)} \; b_1^{(3)}
    \end{pmatrix}$, then $$\det(p_1C_1^*B_1C_1+p_2b_2 C_2^*C_2)=\det \Bigl ( \begin{pmatrix}
        \frac{1}{2}b_1^{(1)} \; \frac{1}{\sqrt 2}b_1^{(2)} \\
        \frac{1}{\sqrt{2}}b_1^{(2)} \; b_1^{(3)}
    \end{pmatrix}+\begin{pmatrix}
        \frac{1}{2}b_2 \; 0 \\
        0 \; 0
    \end{pmatrix} \Bigl)=\frac{1}{2}\det(B_1)+\frac{1}{2}b_2b_1^{(3)}.$$ 
From the previous computation, it follows that $$\mathrm{BL}(\textbf{C},p,\textbf{Q})=\sup_{0<B_1 \leq Q_1, 0<b_2 \leq 1}\frac{\det(B_1)b_2}{\frac{1}{2}\det(B_1)+\frac{1}{2}b_2b_1^{(3)}}=\sup_{0<B_1 \leq Q_1}\frac{\det(B_1)}{\frac{1}{2}\det(B_1)+\frac{1}{2}b_1^{(3)}}$$ and $$\det(B_1)=b_1^{(1)}b_1^{(3)}-(b_1^{(2)})^2 \leq b_1^{(1)}b_1^{(3)} \leq b_1^{(3)}$$ since $B_1 \leq Q_1$ implies $b_1^{(1)} \leq 1.$ \hfill \\ \\  In conclusion, one has $$\frac{\det(B_1)}{\frac{1}{2}\det(B_1)+\frac{1}{2}b_1^{(3)}} \leq 1$$ with equality if $b_1^{(3)}=\det(B_1)$ or equivalently, $b_1^{(1)}=1,b_1^{(2)}=0.$ This shows that the supremum is indeed $1$ and the maximum is attained at $B_1=\begin{pmatrix}
    1 \; 0 \\
    0 \; \alpha
\end{pmatrix}$ for $\alpha \in (0;\lambda], b_2=1.$ \hfill \\ The Brascamp-Lieb inequality associated with the chosen data says that if $f_1$ is more log-convex than $g_{Q_1}$ and $f_2$ is more log-convex than $g_{Q_2}$, then $$\int_{\mathbb{R}^2}f_1 \Bigl(\frac{1}{\sqrt{2}}x_1,x_2 \Bigl)f_2 \Bigl(\frac{1}{\sqrt{2}}x_1 \Bigl)\, dx_1 dx_2 \leq \int_{\mathbb{R}^2}f_1 \int_{\mathbb{R}}f_2.$$ \hfill \\

Clearly, $\mathrm{Ran}(I_2-C_1C_1^*)=\mathrm{span}\{e_1\}, \ker(I_2-C_1C_1^*)=\mathrm{span} \{e_2\}, \ker(C_1)=\{0\}, W_2=\{0\}$ so that $K=W_1 \cap \ker(C_2)=\mathrm{span} \{e_2\}$ is critical.
From \cref{Fundamental_eq_equality_cases}, one may show that there is equality iff $$f_1(x_1,x_2)=c_1e^{-\pi x_1^2-ax_1}u(x_2),  f_2(x_1)=c_2e^{-\pi x_1^2-ax_1}$$ with $u$ being more log-convex than $e^{-\lambda \pi  x^2}.$ 
\end{example}

\hfill \\


\begin{example}

Let us now look at an example where the assumption $C_j^*C_jK \subset K$ for any $1 \leq j \leq m$ does not hold. 
    Set $n=2,n_1=2,n_2=n_3=n_4=1, p_1=p_2=p_3=p_4=\frac{1}{2},$ and $$C_1=\begin{pmatrix}
        1 \; 0 \\
        0 \; \frac{1}{\sqrt{2}}
    \end{pmatrix}, \, C_2=\Bigl(0,\frac{1}{\sqrt 2}\Bigl), \, C_3=\Bigl(\frac{1}{\sqrt 2},\frac{1}{\sqrt 2}\Bigl), C_4=\Bigl(\frac{1}{\sqrt 2},-\frac{1}{\sqrt 2}\Bigl).
   $$ Consider $$Q_1=\begin{pmatrix}
       2 \; 0 \\
       0 \; 1
   \end{pmatrix},Q_2=1,Q_3=Q_4=2.$$ Then, $$C_1^*C_1=\begin{pmatrix}
       1 \; 0 \\
       0 \; 1/2
   \end{pmatrix}, C_2^*C_2=\begin{pmatrix}
       0 \; 0 \\
       0 \; 1/2
   \end{pmatrix}, C_3^*C_3=\begin{pmatrix}
       1/2 \; 1/2 \\
       1/2 \; 1/2
   \end{pmatrix}, C_4^*C_4=\begin{pmatrix}
       1/2 \; -1/2 \\
       -1/2 \; 1/2
   \end{pmatrix}$$ which implies that $\sum_{i=1}^4 p_iC_i^*C_i=I_2$ and it is easy to see that the other conditions for $(\textbf{C},p,\textbf{Q})$ to be a generalized geometric Brascamp-Lieb datum are satisfied. \hfill \\ \\ Moreover, $$\begin{array}{lll}
   \ker(C_1)=\{0\}, C_1^*\ker(I_2-C_1C_1^*)=\mathrm{span}\{e_1\}; \\
   \nm 
   \ker(C_2)=\mathrm{span}\{e_1\}, C_2^*\ker(1-C_2C_2^*)=\{0\}; \\
   \nm 
   \ker(C_3)=\mathrm{span}\{(1,-1)\}, C_3^*\ker(1-C_3C_3^*)=\mathrm{span}\{(1,1)\}; \\
   \nm 
   \ker(C_4)=\mathrm{span}\{(1,1)\}, C_4^*\ker(1-C_4C_4^*)=\mathrm{span}\{(1,-1)\}.\end{array}$$ Hence, $K=\mathrm{span}\{e_1\}$ and one can notice that $C_3^*C_3K \not \subset K, C_4C_4^*K \not \subset K.$ Furthermore, $K$ is not critical as $\sum_{i=1}^4 p_i\dim(C_iK)=\frac{3}{2}>\dim(K)=1$ and in fact, the datum is simple. This means that the extremizers are expected to be Gaussian. Let us show that this expectation is founded. 
   \hfill \\ \\
   First, we write the Brascamp-Lieb inequality with such a datum: for $f_i$ more log-convex than $g_{Q_i}$ for any $1 \leq i \leq 4$, $$\int_{\mathbb{R}^2}f_1 \Bigl(x,\frac{1}{\sqrt 2}y \Bigl)^{1/2}f_2\Bigl(\frac{1}{\sqrt 2}y \Bigl)^{1/2}f_3 \Bigl(\frac{1}{\sqrt 2}x+\frac{1}{\sqrt 2}y\Bigl)^{1/2}f_4 \Bigl(\frac{1}{\sqrt 2}x-\frac{1}{\sqrt 2}y \Bigl)^{1/2}\, dx \, dy \leq \prod_{i=1}^4 \Bigl( \int_{\mathbb R^{n_i}}f_i\Bigl)^{1/2}.$$ 
   One can directly take $t=1$ as it does not change anything to the argument. Using the characterization of $D^2(\log \tilde Q)$ when in terms of special eigenspaces of $C_1$ or $C_2$, one gets $$D^2(\log \tilde Q)(x,y)=\begin{pmatrix}
       \partial_{11}(\log \tilde Q)(x,y) \; 0 \\
       0 \qquad \qquad -2\pi
   \end{pmatrix}$$ with $\partial_{11}(\log \tilde Q) \geq -4\pi.$ Then, $\partial_1 (\log \tilde Q)(x,y)$ is constant in $y$ and so $$\log \tilde Q(x,y)=\log Q(x)-\pi y^2-by \text{ and } \nabla (\log \tilde Q)(x,y)=((\log Q)'(x),-2\pi y-b).$$
   \hfill \\
   We  may forget about the constants since, by \cref{Closure_prop_extremizers}, the extremizers that will be derived can be multiplied by any positive constants. 
   Now, let us use \cref{Fundamental_eq_equality_cases} for $j=1,2,3,4.$ First, $$(\log f_3)'\Bigl(\frac{1}{\sqrt 2}x+\frac{1}{\sqrt 2}y \Bigl)=C_3\nabla \log \tilde Q(x,y)=\frac{1}{\sqrt 2}(\log Q)'(x)-\frac{2\pi}{\sqrt 2}y-\frac{b}{\sqrt 2},$$ $$(\log f_4)'\Bigl(\frac{1}{\sqrt 2}x-\frac{1}{\sqrt 2}y \Bigl)=C_4\nabla \log \tilde Q(x,y)=\frac{1}{\sqrt 2}(\log Q)'(x)+\frac{2\pi}{\sqrt 2}y+\frac{b}{\sqrt 2}$$ and taking $y=0$ shows that $$(\log f_3)'(\frac{1}{\sqrt 2}x)=\frac{1}{\sqrt 2}(\log Q)'(x)-\frac{b}{\sqrt 2},(\log f_4)'(\frac{1}{\sqrt 2}x)=\frac{1}{\sqrt 2}(\log Q)'(x)+\frac{b}{\sqrt 2}.$$ \hfill \\ This implies that $$\log f_3(x)=\frac{1}{2}(\log Q)(\sqrt 2x)-\frac{b}{\sqrt 2}x ,\log f_4(x)=\frac{1}{2}(\log Q)(\sqrt 2x)+\frac{b}{\sqrt 2}x.$$ For $j=1,2$, we have $$\nabla (\log f_1)\Bigl(x,\frac{1}{\sqrt 2}y \Bigl)=C_1\nabla \log \tilde Q(x,y)=\Bigl((\log Q)'(x),-\frac{2\pi}{\sqrt 2}y-\frac{b}{\sqrt 2}\Bigl)$$ and $$(\log f_2)'\Bigl(\frac{1}{\sqrt 2}y \Bigl)=C_2\nabla \log \tilde Q(x,y)=-\frac{2\pi}{\sqrt 2}y-\frac{b}{\sqrt 2}$$ which implies that $$\log f_1(x,y)=\log Q(x)-\pi y^2-\frac{b}{\sqrt 2}y, \log f_2(y)=-\pi y^2-\frac{b}{\sqrt 2}y.$$  To find the final form of $Q$, we use the equation of $(\log f_3)'$ to get $$(\log Q)'(y)=(\log Q)'(0)-2\pi y,$$ which means that $Q(x)=e^{-\pi x^2-ax}$ and consequently,
$$\begin{array}{lll} f_1(x,y)=c_1 e^{-\pi x^2-ax}e^{-\pi y^2-by}, f_2(y)=c_2e^{-\pi y^2-by}, \\ \\
f_3(x)=c_3Q(\sqrt{2}x)^{1/2}e^{-bx}=c_3e^{-\pi x^2-(\frac{a}{\sqrt 2}+b)x}, \\ \\
f_4(x)=c_4Q(\sqrt 2x)^{1/2}e^{bx}=c_4 e^{-\pi x^2-(\frac{a}{\sqrt 2}-b)x}\end{array}$$ for any $x,y \in \mathbb{R} \text{ and } a,b \in \mathbb{R}$ arbitrary. Let us show that if $f_1,f_2,f_3,f_4$ are of this form, then they are extremizers. Indeed, first notice that $f_i$ are more log-convex than $g_{Q_i}$ for any $1 \leq i \leq 4.$ Second, let us compute $$\begin{array}{lll}
\ds \int_{\mathbb{R}^2}\prod_{j=1}^4 f_j(C_j(x,y))^{1/2}\, dx \, dy &=& \ds \int_{\mathbb{R}^2}e^{-\pi x^2-ax}e^{-\pi y^2-\sqrt 2 by}\, dx \, dy \\
\nm 
&=& \ds e^{a^2/4\pi}e^{b^2/2\pi} 
\end{array}$$ and $$\begin{array}{lll}
    \ds \prod_{j=1}^4 \Bigl( \int f_j \Bigl)^{1/2} &=& \ds e^{a^2/8\pi}e^{b^2/8\pi} e^{b^2/8\pi}e^{(\frac{a}{\sqrt 2}+b)^2/8\pi}e^{(\frac{a}{\sqrt 2}-b)^2/8\pi}  \\
     &=& e^{a^2/4\pi}e^{b^2/2\pi},
\end{array}$$ since $\int_{\mathbb{R}}e^{-\pi x^2-cx}\, dx=e^{c^2/4\pi}$ for any $c \in \mathbb{R}.$  
\end{example}



\bigskip 

\hfill \\







Recall that $$D^2(\log \tilde Q)(x)=\left(\begin{array}{ c | c }
    \ker(C_j) & 0 \\
    \hline
    0 & \left. \begin{array}{ c | c }
    \mathrm{diag}(-\frac{2\pi}{t},\cdots,-\frac{2\pi}{t}) & 0 \\
    \hline
    0 & W_j
  \end{array} \right.
  \end{array}\right)$$
where $W_j=C_j^*\ker(\mathrm{id}_{H_j}-C_jC_j^*)$ and the second block is with respect to the basis $\{C_j^*v_i^{(j)}: 1 \leq i \leq k_j\}=C_j^*\mathrm{Ran}(\mathrm{id}_{H_j}-C_jC_j^*) =: \tilde W_j.$

We introduce the definition of generalized independent subspace that plays a crucial role in \cite{Vald2}. 

\begin{definition}
    A subspace $\tilde H$ of $K$ is generalized independent with respect to the generalized geometric datum if it is not $\{0\}$ and has the form $$\tilde H=\bigcap_{j=1}^m \tilde H_j^{a}$$ where for each $j, \tilde H_j^a$ is either $\ker(C_j)$ or $W_j=C_j^*\ker(\mathrm{id}_{H_j}-C_jC_j^*)=\ker(\mathrm{id}_H-C_j^*C_j).$ 
\end{definition}

Note that for $x \in \tilde H$ a generalized independent subspace of $K$, either $C_ix=0$ or $C_i^*C_ix=x$ for any $1 \leq i \leq m$. This implies that $C_i^*C_i \tilde H \subset \tilde H$ for all $1 \leq i \leq m.$ By \cref{Commuting_implies_critical}, any generalized independent subspace is critical. 
\hfill \\

Moreover, define $\tilde K_{\mathrm{dep}}$ as the orthogonal complement of $\bigoplus_{k=1}^{k_0}\tilde K_k$ as subspace of $K$ where $\{\tilde K_k : k=1,\cdots,k_0\}$ is an enumeration of the generalized independent subspaces of $K$. Then, arguing exactly as in \cite[Lemma 13]{Vald2}, one obtains $$D^2\log \tilde Q(x)=\left(\begin{array}{ c | c }
    \mathrm{diag}(-\frac{2\pi}{t},\cdots,-\frac{2\pi}{t}) & 0 \\
    \hline
    0 & \left. \begin{array}{ c | c } 
    \nm
    \tilde K_1   & 0   \\ 
    \hline
    0 &  \left. \begin{array}{c | c} \nm \ddots & 0 \\
    \hline
    0 & \left. \begin{array}{c | c} \nm \tilde K_{k_0} & 0 \\
    \nm
    \hline
     0 & \tilde K_{\mathrm{dep}}  \end{array} \right.  \end{array} \right. 
    \end{array} \right.
  
  \end{array}\right) .$$

\hfill \\



Intuitively, we will decompose further $\tilde K_{\mathrm{dep}}$ into critical subspaces and a non-critical part, \textit{i.e.} $$\tilde K_{\mathrm{dep}}=\bigoplus_{k=k_0+1}^{k_1}\tilde K_k \oplus \text{Remaining non-critical part}.$$ On $\bigoplus_{k=1}^{k_1}\tilde K_k$, the Brascamp-Lieb data will be reduced to a geometric Brascamp-Lieb data with $\bigoplus_{k=k_0+1}^{k_1}\tilde K_k$ being the dependent part of the independent decomposition using the terminology of \cite{Vald2}. Moreover, on $\tilde K_j$ for $1 \leq j \leq k_0$, we have an arbitrary Jacobian as in the situation of Hölder's inequality and on $\tilde K_{\mathrm{dep}}$, we have a constant Jacobian, which means that the functions are Gaussians as in the case of Young's inequality. 

\hfill \\ Thus, for any $1 \leq j \leq m$, we have $$\log \tilde Q(x)=-\frac{\pi}{t} \langle P_{\tilde W_j}x,P_{\tilde W_j}(x+b_j)\rangle+u_j^{\parallel}(P_{\ker(C_j)}x)+u_j^\perp(P_{W_j}x)$$ 
and
$$\log \tilde Q(x)=-\frac{\pi}{t} \langle P_{K^\perp}x,P_{K^\perp}(x+b_\perp)\rangle+\sum_{k=1}^{k_0}u_{\tilde K_k}(P_{\tilde K_k}x)+u_{\tilde K_{\mathrm{dep}}}({P_{\tilde K_{\mathrm{dep}}}}x).$$ 





\begin{proposition}
    Suppose that $K$ has no generalized independent subspace. Then any optimizer must be Gaussian. 
\end{proposition}

\begin{proof}
Recall that
    $$\nabla \log \tilde Q(x)=-\frac{\pi}{t} P_{\tilde W_j}^*P_{\tilde W_j}(2x+b_j)+P_{\ker(C_j)}^*\nabla u_j^{||}(P_{\ker(C_j)}x)+P_{W_j}^*\nabla u_j^\perp(P_{W_j}x).$$ As it is well-known that the Fourier transform of a distribution supported at the origin is a polynomial (and reversely, the Fourier transform of a polynomial is a distribution supported at the origin), the Fourier transform of $\nabla \log \tilde Q(x)$ is supported on $\ker(C_j) \cup W_j$ for any $1 \leq j \leq m.$ As this holds for any $1 \leq j \leq m$, it is in fact supported on $$\bigcap_{j=1}^m (\ker(C_j) \cup W_j)=\{0\},$$ since $K$ has no generalized independent subspace. Hence, using growth estimates, $\nabla \log \tilde Q(x)$ is a linear polynomial. More precisely, as we are assuming that $f_j$ is a Schwartz function for any $1 \leq j \leq m$, one may show that $\nabla \log(\tilde f_j)(C_jx)$ has linear growth in $x$ (see \cite[Lemma 5]{Vald2}) which shows the previous claim as $$\nabla \log \tilde Q(x)=\sum_{j=1}^m p_jC_j(\nabla \log \tilde f_j)(C_jx).$$ This means that all $\tilde f_j$ are Gaussians. 
\end{proof} 


We can consider the largest subspace of $\tilde K_{\mathrm{dep}}$ such that it is invariant under $C_i^*C_i$ for all $1 \leq i \leq m$ and denote its orthogonal complement in $\tilde K_{\mathrm{dep}}$ by $H_{\mathrm{dep}}^{\perp,K}.$ Moreover, by the next lemma, this subspace is critical. 

\begin{lemma} \label{Commuting_implies_critical}
    If $V \subset K$ satisfies $C_i^*C_iV \subset V$ for all $1 \leq i \leq m$ where $(\textbf{C},p,\textbf{Q})$ is a generalized geometric BL, then $V$ is critical.
\end{lemma}

\begin{proof}
     By assumption, $C_i^*C_iV \subset V$ for any $1 \leq i \leq m.$ Thus for any $x \in C_iV, y \in V,$ $$\langle C_i^*x-(C_i|_V)^*x,y\rangle=0 \Longrightarrow (C_i|_V)^*x=C_i^*x \quad \forall 1 \leq i \leq m \text{ since } C_i^*x \in V.$$ This shows that $(C_i|_V)^*=C_i^*|_{C_iV}$ for any $1 \leq i \leq m.$     
     Moreover, for $x \in C_iV$, there exists $z \in V \subset K$ such that $x=C_iz.$ As $z=C_i^*z_i+\tilde z_i$ where $z_i \in \ker(\mathrm{id}_{H_i}-C_iC_i^*), \tilde z_i \in \ker(C_i),$ $$ x=C_iC_i^* z_i= z_i$$ so $C_i|_V(C_i|_V)^*x=C_iC_i^*x=x.$ In other words, one has $$C_i|_V(C_i|_V)^*=\mathrm{id}_{C_iV} \quad \forall 1 \leq i \leq m.$$ Moreover, $$\sum_{i=1}^m p_i(C_i|_V)^*C_i|_V=\mathrm{id}_V$$ by restricting the equality $\ds \sum_{i=1}^m p_iC_i^*C_i=\mathrm{id}_H$ on $V$, which holds since $(\textbf{C},p,\textbf{Q})$ is a generalized geometric BL.  Taking traces, the claim follows.  
\end{proof}




This shows that there exists a generalized maximum critical decomposition of $\tilde K_{\mathrm{dep}}$ $$\tilde K_{\mathrm{dep}}=\bigoplus_{k=k_0+1}^{k_1}\tilde K_k \oplus H_{\mathrm{dep}}^{\perp,K},$$ which implies that $$H_{\mathrm{dep}}=\bigoplus_{k=1}^{k_1}\tilde K_k,$$ such that the purely quadratic term in this Gaussian is the tensor product of multiples of the identity operator on each $\tilde K_k$ for $k_0+1 \leq k \leq k_1$, $H_\mathrm{dep}^{\perp,K}$ appearing in this decomposition. Furthermore, by definition of $\bigoplus_{k=k_0+1}^{k_1}\tilde K_k$ and as $\tilde K_k$ for any $1 \leq k \leq k_0$ is invariant under all $C_i^*C_i$, this is also the case for $H_{\mathrm{dep}}$, which means that $H_{\mathrm{dep}}$ is critical by the previous lemma. Hence, we deduce that $$\begin{array}{lll}
\log \tilde Q(x)&=& \ds -\frac{\pi}{t} \langle P_{K^\perp}x,P_{K^\perp}(x+b_\perp)\rangle+\sum_{k=1}^{k_0}u_{\tilde K_k}(P_{\tilde K_k}x) \\
&& \ds -\sum_{k=k_0+1}^{k_1}d_k\langle P_{\tilde K_k}x,P_{\tilde K_k}(x+b_k)\rangle-c\langle P_{H_{\mathrm{dep}}^{\perp,K}}x,P_{H_{\mathrm{dep}}^{\perp,K}}(x+\tilde b_\perp)\rangle.\end{array}$$ where $b_\perp \in K^\perp, b_k \in \tilde K_k$ and $\tilde b_\perp \in H_{\mathrm{dep}}^{\perp,K}.$ 

\hfill \\

The final step to be able to determine the extremizers is to understand the special form of the Gaussian constant $c$ on the block $H_{\mathrm{dep}}^{\perp,K}.$  

\hfill \\

Note that $(C_jH_{\mathrm{dep}})^{\perp,H_j}=C_j(H_{\mathrm{dep}}^\perp)=C_j(K^\perp \oplus H_{\mathrm{dep}}^{\perp,K}).$ Indeed, if $x \in H_{\mathrm{dep}}^\perp$, then $$\langle C_jx,y\rangle=\langle x,C_j^*y\rangle=0 \quad \forall y \in C_jH_{\mathrm{dep}}$$ since $C_j^*C_jH_{\mathrm{dep}} \subset H_{\mathrm{dep}}$ for any $1 \leq j \leq m.$ This implies that $C_j(H_{\mathrm{dep}}^\perp) \subset (C_jH_{\mathrm{dep}})^{\perp,H_j}$. Furthermore, $H_j=C_jH_{\mathrm{dep}} +C_j(H_{\mathrm{dep}}^\perp)$ since $C_j$ is surjective and consequently, $$\dim((C_jH_{\mathrm{dep}})^{\perp,H_j})=\dim(H_j)-\dim(C_jH_{\mathrm{dep}}) \leq \dim(C_jH_{\mathrm{dep}}^\perp).$$ This proves the desired claim. Using $C_j(H_{\mathrm{dep}}^\perp)=(C_jH_{\mathrm{dep}})^{\perp,H_j}$, we have \begin{equation} \label{Proj_commute_under_C_j}  C_jP_{H_{\mathrm{dep}}^\perp}^*P_{H_{\mathrm{dep}}^\perp}=P_{(C_jH_{\mathrm{dep}})^{\perp,H_j}}^*P_{(C_jH_{\mathrm{dep}})^{\perp,H_j}}C_j.\end{equation} 

\hfill \\



Let us now prove that $c=\frac{\pi}{t}$.
From \cref{Fundamental_eq_equality_cases}, we see that for any $x \in H$, $$\begin{array}{lll}
    P_{(C_jH_{\mathrm {dep}})^{\perp,H_j}}^*P_{(C_jH_{\mathrm{dep}})^{\perp,H_j}}\nabla \log \tilde f_j(C_jx) &=& P_{(C_jH_{\mathrm {dep}})^{\perp,H_j}}^*P_{(C_jH_{\mathrm{dep}})^{\perp,H_j}}C_j\nabla \log \tilde Q(x)
    \\
    \nm 
    &=& C_jP_{H_{\mathrm{dep}}^\perp}^*P_{H_{\mathrm{dep}}^\perp}\nabla \log \tilde Q(x) \\
    \nm 
    &=& \ds -\frac{\pi}{t} C_jP_{K^\perp}^*P_{K^\perp}(2x+b_\perp) \\
    \nm 
    && \ds -c C_jP_{H_\mathrm{dep}^{\perp,K}}^*P_{H_{\mathrm{dep}}^{\perp,K}}(2x+\tilde b_\perp).\end{array}$$ For $x \in \ker(C_j)$, denote by $y=P_{K^\perp}^*P_{K^\perp}x, z=P_{H_{\mathrm{dep}}^{\perp,K}}^*P_{H_{\mathrm{dep}}^{\perp,K}}x.$ Then by \cref{Proj_commute_under_C_j}, $$C_j(P_{H_\mathrm{dep}^{\perp}}^*P_{H_\mathrm{dep}^{\perp}}x)=0 \Longrightarrow C_jy=-C_jz$$ and hence the previous equation for $x \in \ker(C_j)$ reads $$\begin{array}{lll}P_{(C_jH_{\mathrm {dep}})^{\perp,H_j}}^*P_{(C_jH_{\mathrm{dep}})^{\perp,H_j}}\nabla \log \tilde f_j(0) &=& \ds 2\Bigl(-\frac{\pi}{t}+c \Bigl)C_jP_{K^\perp}^*P_{K^\perp}x-\frac{\pi}{t} C_jP_{K^\perp}^*P_{K^\perp}b_\perp \\
    \nm 
    && \ds -c \, C_jP_{H_\mathrm{dep}^{\perp,K}}^*P_{H_{\mathrm{dep}}^{\perp,K}}\tilde b_\perp.\end{array}$$ Subtracting this equation from the one obtained when considering $x=0$, we get $$\Bigl(-\frac{\pi}{t}+c \Bigl)C_jP_{K^\perp}^*P_{K^\perp}x=0 \quad \forall x \in \ker(C_j).$$
    
    Unless $P_{K^\perp}^*P_{K^\perp}\ker(C_j) \subset \ker(C_j)$ for any $1 \leq j \leq m$, we can conclude that $c=\frac{\pi}{t}.$ \hfill \\ \\ Else, $P_{K^\perp}^*P_{K^\perp}\ker(C_j) \subset \ker(C_j)$ or equivalently, $$P_K^*P_K(\ker(C_j)) \subset \ker(C_j) \text{ for any } 1 \leq j \leq m.$$ 
    This means that\footnote{As $P_K^*P_K\ker(C_j) \subset \ker(C_j)$ for any $1 \leq j \leq m, P_K^*P_K\ker(C_j)^\perp \subset \ker(C_j)^\perp$, which means that for any $x \in K,$ $$ x=P_K^*P_KP_{\ker(C_j)}^*P_{\ker(C_j)}x+P_K^*P_KP_{\ker(C_j)^\perp}^*P_{\ker(C_j)^\perp}x.$$} $$K=(K \cap \ker(C_j))\oplus (K \cap \ker(C_j)^\perp)=(K \cap \ker(C_j)) \oplus (K \cap \ker(\mathrm{id}_H-C_j^*C_j)).$$ But then, for $x \in K,$ $$ x=y_j+z_j \text{ for } y_j \in K \cap \ker(C_j), z_j \in K \cap \ker(\mathrm{id}_H-C_j^*C_j),$$ which implies that $C_j^*C_jx=z_j \in K.$ In conclusion, $$C_j^*C_jK \subset K \quad \forall 1 \leq j \leq m.$$ This proves that $H_{\mathrm{dep}}=K, H_{\mathrm{dep}}^{\perp,K}=\{0\}$ and in this case, there is nothing to show.   \hfill \\
    

Having verified that $c=\frac{\pi}{t},$ we can conclude as follows: 








By the previous form of $\log \tilde Q$, we obtain  $$\begin{array}{lll}
    \nabla \log \tilde f_j(C_jx) &=& \ds C_j\nabla \log \tilde Q(x)=-\frac{\pi}{t} C_jP_{K^\perp}^*P_{K^\perp}(2x+b_\perp)-\frac{\pi}{t} C_jP_{H_{\mathrm{dep}}^{\perp,K}}^*P_{H_{\mathrm{dep}}^{\perp,K}}(2x+\tilde b_\perp)  \\
    \nm 
    &&  \ds +\sum_{k=1}^{k_0}C_jP_{\tilde K_k}^*\nabla u_{\tilde K_k}(P_{\tilde K_k}x)-\sum_{k=k_0+1}^{k_1}d_kC_jP_{\tilde K_k}^*P_{\tilde K_k}(2x+b_k) \\
    \nm 
    &=& \ds -\frac{\pi}{t} C_jP_{H_{\mathrm{dep}}^\perp}^*P_{H_{\mathrm{dep}}^\perp}(2x+b_\perp+\tilde b_\perp)+\sum_{k=1}^{k_0}C_jP_{\tilde K_k}^*\nabla u_{\tilde K_k}(P_{\tilde K_k}x) \\
    \nm 
    && \ds -\sum_{k=k_0+1}^{k_1}d_kC_jP_{\tilde K_k}^*P_{\tilde K_k}(2x+b_k) 
\end{array}$$ and recall \cref{Proj_commute_under_C_j}. 
Moreover, each term in the first sum is zero unless $\tilde K_k \subset C_jH_{\mathrm{dep}}$ in which case $C_jP_{\tilde K_k}=P_{\tilde K_k}C_j^*C_j.$ For the second sum, we have $$C_jP_{\tilde K_k}^*P_{\tilde K_k}=C_jP_{\tilde K_k}^*C_j^*C_j.$$ 



Letting $t \rightarrow 1$ and by the next remark, we get the following statement. 

\begin{theorem} \label{Main_Thm_Eq_Geometric_BL}
Let $(\textbf{C},p,\textbf{Q})$ be a generalized geometric Brascamp-Lieb datum and suppose that $(f_j)$ is an extremizer for this datum. 
Recall that $H_{\mathrm{dep}}$ is the largest subspace of $K$ that contains $\bigoplus_{k=1}^{k_0}\tilde K_k$ (where $\{\tilde K_k : k=1,\cdots,k_0\}$ is an enumeration of the generalized independent subspaces of $K$) and is invariant under $C_i^*C_i$ for any $1 \leq i \leq m.$ Then, there exist integrable functions $u_k: H_k \rightarrow \mathbb{R}, k=1,\cdots,k_0$, a critical decomposition $\tilde K_1 \oplus \cdots \oplus \tilde K_{k_0} \oplus \tilde K_{k_0+1} \oplus \cdots \oplus \tilde K_{k_1}$ of $H_{\mathrm{dep}}$, positive constants $c_j$ for $j=1,\cdots,m$ and $d_k$ for $k_0+1 \leq k \leq k_1$ and two elements $b \in H_{\mathrm{dep}}^\perp, \tilde b \in \bigoplus_{k=k_0+1}^{k_1}\tilde K_k$ such that 
$$\begin{array}{lll}f_j(x) &=&
    
    \ds c_je^{-\pi \langle P_{(C_jH_{\mathrm{dep}})^{\perp,H_j}}x,P_{(C_jH_{\mathrm{dep}})^{\perp,H_j}}(x+C_jb)\rangle} \\
    \nm 
    && \ds \prod_{k=1}^{k_0}u_k(\tilde P_{j,k}C_j^*x)\prod_{k=k_0+1}^{k_1}e^{-d_k \langle \tilde P_{j,k}C_j^*x,\tilde P_{j,k}(C_j^*x+\tilde b)\rangle} \end{array}$$ where $\tilde P_{j,k}$ is the orthogonal projection from $H$ to $C_jH_{\mathrm{dep}} \cap \tilde K_k.$ 
    Moreover, $u_k$ and $d_k $ are such that $f_j$ is more log-convex than $g_{Q_j}$ for any $1 \leq j \leq m.$
\end{theorem}


\begin{remark}
    To remove the restriction that $(f_j)_{1 \leq j \leq m}$ is a tuple of Schwartz functions, we use \cref{Closure_prop_extremizers}. Indeed, when $(f_j)_{1 \leq j \leq m}$ is an optimizer where each $f_j$ is only assumed to be an $L^1$ function on $H_j$, denoting $$g_j(x)=e^{-\pi \|x\|^2_{H_j}},$$ $(g_j)_{1 \leq j \leq m}$ is an optimizer since the datum is generalized geometric and $((f_j * g_j)g_j(\frac{\cdot}{\sqrt 2}))_{1 \leq j \leq m}$ is also an optimizer. Each function in the  tuple we just constructed is a positive Schwartz function and by what was proven, this new optimizer has the desired form. But if $(f * g)g=k$ where $k(x)=k_P(Px)k_{P^\perp}(P^\perp x), P$ and $P^\perp$ are projections onto orthogonal subspaces and $g(x)=g_P(Px)g_{P^\perp}(P^\perp x)$ is strictly positive, then $$f(x)=f_P(Px)f_{P^\perp}(P^\perp x).$$ Moreover, if $g$ and $k$ are Gaussians, $f$ will also be Gaussian. Hence, the previous lemmas will also hold for the $L^1$ tuple $(f_j)_{1 \leq j \leq m}$ and any tuple must have the desired form. 
\end{remark}








\bigskip







\hfill \\ Let us now drop the condition that $(\textbf{C},p,\textbf{Q})$ is a generalized geometric Brascamp-Lieb datum. From \cref{Prop_Equiv_BL_constant}, the equations $$\begin{array}{ccc}
    M = \ds \sum_{j=1}^m p_jC_j^*B_jC_j  \\
    \nm 
     B_j^{-1} \geq C_jM^{-1}C_j, (B_j^{-1}-C_jM^{-1}C_j^*)(Q_j-B_j)=0
\end{array}$$ have a solution $M$ and $B_j$ with symmetric positive definite linear transformations and denoting $C_j'=B_j^{1/2}C_jM^{-1/2}, Q_j'=B_j^{-1/2}Q_jB_j^{-1/2}$, $(\textbf{C'},p,\textbf{Q'})$ is a generalized geometric datum. Moreover, if $(f_j)_{1 \leq j \leq m}$ is an optimizer for $(\textbf{C},p,\textbf{Q})$, then $(f_j \circ B_j^{-1/2})_{1 \leq j \leq m}$ is an optimizer for $(\textbf{C'},p,\textbf{Q'}).$ Conversely, if $(f_j')_{1 \leq j \leq m}$ is an optimizer for $(\textbf{C'},p,\textbf{Q'})$, then $(f_j' \circ B_j^{1/2})_{1 \leq j \leq m}$ is an optimizer for $(\textbf{C},p,\textbf{Q}).$ Hence, we get the following theorem: 

\begin{theorem} \label{Main_Thm_Eq_BL}
    Let $(\textbf{C},p,\textbf{Q})$ be an extremizable Brascamp-Lieb datum. Assume that $(f_j)_{1 \leq j \leq m}$ is an extremizer for this datum. Set $(\textbf{C'},p,\textbf{Q'})$ to be the generalized geometric datum equivalent to $(\textbf{C},p,\textbf{Q}).$ In addition, let $H_\mathrm{dep}$ corresponding to the datum $(\textbf{C'},p,\textbf{Q'}).$ Then, there exist integrable functions $u_k: H_k \rightarrow \mathbb{R}, k=1,\cdots,k_0$, a critical decomposition $H_{\mathrm{dep}}=\bigoplus_{k=1}^{k_0}\tilde K_k \oplus \tilde K_{k_0+1}\oplus \cdots \oplus \tilde K_{k_1}$, positive constants $c_j$ for $j=1,\cdots,m$ and $d_k$ for $k_0+1 \leq k \leq k_1$ and two elements $b \in H_{\mathrm{dep}}^\perp, \tilde b \in \bigoplus_{k=k_0+1}^{k_1}\tilde K_k$ such that $$\begin{array}{lll}
    f_j(x) &=&
    
    \ds c_je^{-\pi \langle P_{(C'_jH_{\mathrm{dep}})^{\perp,H_j}}B_j^{1/2}x,P_{(C'_jH_{\mathrm{dep}})^{\perp,H_j}}(B_j^{1/2}x+C_j'b)\rangle} \\
    \nm 
    && \ds \prod_{k=1}^{k_0}u_k(\tilde P_{j,k}C_j'^*B_j^{1/2}x)\prod_{k=k_0+1}^{k_1}e^{-d_k \langle \tilde P_{j,k}C_j'^*B_j^{1/2}x,\tilde P_{j,k}(C_j'^*B_j^{1/2}x+\tilde b)\rangle} \end{array}$$ where $\tilde P_{j,k}$ is the orthogonal projection from $H$ to $C'_jH_{\mathrm{dep}} \cap \tilde K_k.$ Moreover, $u_k$ and $d_k$ are such that $f_j$ is more log-convex than $g_{Q_j}$ for any $1 \leq j \leq m.$ 
    \hfill \\ Conversely, all functions of this form are optimizers for this problem. 
\end{theorem}

\hfill \\

\begin{remark}
Recall that we have a dual inequality \cref{Reg_Dual_BL_Log_Concavity} to the generalized geometric BL that we will call the generalized geometric Barthe's inequality: Assume that $\sum_{j=1}^m p_jC_j^*C_j=\mathrm{id}_H, C_jC_j^* \leq \mathrm{id}_{H_j}$ and $Q_j \leq \mathrm{id}_{H_j}, (\mathrm{id}_{H_j}-C_jC_j^*)(\mathrm{id}_{H_j}-Q_j)=0.$ Then, for $g_j$ more log-concave than $g_{Q_j}$ for any $1 \leq j \leq m$, $$\int_{H}\sup_{x=\sum_{j=1}^m p_jC_j^*y_j, y_j \in H_j}\prod_{j=1}^m g_j(y_j)^{p_j}\, dx \geq \prod_{j=1}^m \Bigl(\int_{H_j}g_j \Bigl)^{p_j}.$$ One can determine the equality of cases of the previous inequality by the same method as in \cite{Reverse_BL_equality} (which consequently generalize the main result of \cite{Reverse_BL_equality}).
\end{remark}


\begin{proposition}
    If $g_j$ is an extremizer for the generalized geometric Barthe's inequality, $$\begin{array}{lll}
    \ds g_j(x) &=& \ds c_j e^{-\pi\langle P_{(C_jH_{\mathrm{dep}})^{\perp,H_j}}x, P_{(C_jH_{\mathrm{dep}})^{\perp,H_j}}(x+C_jb)\rangle} \\
    \nm 
    && \ds \prod_{k=1}^{k_0}h_k(\tilde P_{j,k}C_j^*(x-w_j))\prod_{k=k_0+1}^{k_1}e^{-d_k\langle \tilde P_{j,k}C_j^*x,\tilde P_{j,k}(C_j^*x-b_j)\rangle}
    \end{array}$$ where \begin{itemize}
        \item [(i)] $\tilde P_{j,k}$ is the orthogonal projection from $H$ to $C_jH_{\mathrm{dep}} \cap \tilde K_k$, $c_j>0, b_j \in C_jH_\mathrm{dep} \cap \Bigl( \bigoplus_{k=k_0+1}^{k_1}\tilde K_k\Bigl)$ and $w_j \in C_jH_{\mathrm{dep}}$ for $j=1,\cdots,k$, 
        \item[(ii)] $h_k \in L^1(\tilde K_k)$ is non-negative for $k=1,\cdots,l$ and in addition, $h_k$ is log-concave if there exists $\alpha \neq \beta$ with $\tilde K_k \subset C_\alpha H_{\mathrm{dep}} \cap C_\beta H_{\mathrm{dep}},$ 
        \item[(iii)] $d_k>0$ for $k_0+1 \leq k \leq k_1$ and $h_k$ for $1 \leq k \leq k_0$ are such that $g_j$ is more log-concave than $g_{Q_j}$ for any $1 \leq j \leq m.$
    \end{itemize}
\end{proposition}

\hfill \\

\section{Some Applications}

The generalized geometric Brascamp-Lieb inequality reformulated for Gaussian measures gives the following, which generalizes the results of Chen-Dafnis-Paouris in \cite{Chen_Dafnis_Paouris}: 

\begin{theorem}
    Assume that $n \leq m, n_1,\cdots,n_m \leq n$ are positive integers. For any $1 \leq i \leq m$, consider the $n_i \times n_i$ matrices $C_i$ with $$C_iC_i^* \leq I_{n_i}, Q_i \ge I_{n_i} \; \forall i \in I; \quad C_iC_i^* \ge I_{n_i}, Q_i \le I_{n_i} \; \forall i \in I^c, $$ and $(I_{n_i}-C_iC_i^*)(Q_i-I_{n_i})=0$ and $p_i>0$ such that $$\sum_{i=1}^m p_iC_i^*C_i=I_n.$$ Then, for functions $f_i$ that are more log-convex (resp. more log-concave) than $g_{\frac{Q_i-I_{n_i}}{2\pi}}$ for any $i \in I$ (resp. for any $i \in I^c$), $$\int_{\mathbb{R}^n}\prod_{i=1}^m f_i(C_ix)^{p_i}\, d\gamma_n(x) \leq \prod_{i=1}^m \Bigl(\int_{\mathbb{R}^{n_i}}f_i \, d\gamma_{n_i}\Bigl)^{p_i}.$$ 
\end{theorem}


This allows us to retrieve \cite[Theorem 4]{Chen_Dafnis_Paouris} by taking $I=\{1,\cdots,m\}, Q_i=\Lambda I_{n_i}$ for $\Lambda>1$, which implies that $C_iC_i^*=I_{n_i}$ and letting $\Lambda \rightarrow +\infty$. 


\begin{proof}
    Since $(\textbf{C},p,\textbf{Q},I)$ is a generalized geometric BL by assumption, $$\mathrm{BL}(\textbf{C},p,\textbf{Q},I)=1.$$ Now, from $$\mathrm{BL}\Bigl(\textbf{C},p,\frac{\textbf{Q}}{2\pi},I\Bigl)=(2\pi)^{(n-\sum_{i=1}^m p_in_i)/2}\mathrm{BL}(\textbf{C},p,\textbf{Q},I)=(2\pi)^{(n-\sum_{i=1}^m p_in_i)/2},$$
    we obtain
    $$\begin{array}{lll}
        \ds \int_{\mathbb{R}^n} \prod_{i=1}^m f_i(C_ix)^{p_i} \, d\gamma_n(x)  &=& \ds \frac{1}{(2\pi)^{n/2}}\int_{\mathbb{R}^n}\prod_{i=1}^m f_i(C_ix)^{p_i} e^{-\frac{1}{2}\sum_{i=1}^m p_i |C_ix|^2}\, dx  \\
        \nm
        &=& \ds  \frac{1}{(2\pi)^{n/2}}\int_{\mathbb{R}^n}\prod_{i=1}^m \Bigl(f_i(C_ix)e^{-\frac{|C_ix|^2}{2}} \Bigl)^{p_i} \, dx  \\
        \nm 
        &\leq& \ds \frac{(2\pi)^{(n-\sum_{i=1}^m p_in_i)/2}}{(2\pi)^{n/2}}\prod_{i=1}^m \Bigl( \int_{\mathbb{R}^{n_i}}f_i e^{-\frac{|\cdot|^2}{2}} \Bigl)^{p_i} \\
        \nm 
        &=& \ds \prod_{i=1}^m \Bigl( \int f_i \, d\gamma_{n_i}\Bigl)^{p_i}.
    \end{array}$$ Here, we used the fact that $f_i$ being more log-convex (resp. more log-concave) than $g_{\frac{Q_i-I_{n_i}}{2\pi}}$ implies that $f_ie^{-|\cdot|^2/2}$ is more log-convex (resp. more log-concave) than $g_{\frac{Q_i}{2\pi}}.$ 
\end{proof}




\begin{theorem} \label{Gaussian_Version_BL}
    Let $n \leq m$ and $n_1,\cdots,n_m \leq n$ be positive integers and let $N=n_1+\cdots+n_m$. For every $i=1,\cdots,m$, consider the $n_i \times n$ matrices $C_i$ and $n_i \times n_i$ matrices $$C_iC_i^* \leq I_{n_i}, Q_i \geq I_{n_i} \; \forall i \in I; \quad C_iC_i^* \geq I_{n_i}, Q_i \le I_{n_i} \; \forall i \in I^c$$  such that $$(I_{n_i}-C_iC_i^*)(Q_i-I_{n_i})=0 \quad \forall 1 \le i \le m.$$  Let the $N \times N$ diagonal matrix $P$ be given by  $$P=\mathrm{diag}(p_1I_{n_1},\cdots,p_mI_{n_m}).$$ For functions $f_i$ that are more log-convex (resp. more log-concave) than $g_{\frac{Q_i-I_{n_i}}{2\pi}}$ for any $i \in I$ (resp. for any $i \in I^c$), we have that if $CC^* \leq P^{-1}$, then $$\int_{\mathbb{R}^n}\prod_{i=1}^m f_i(C_ix)^{p_i}\, d\gamma_n(x) \leq \prod_{i=1}^n \Bigl( \int_{\mathbb{R}^{n_i}}f_i\, d\gamma_{n_i} \Bigl)^{p_i}.$$
\end{theorem}

\hfill \\

As previously, this theorem recovers \cite[Theorem 3 (i)]{Chen_Dafnis_Paouris}. 


\begin{proof}
    Assume $CC^* \leq P^{-1}.$ This means that $\|A\|_{\mathrm{op}} \leq 1$ where $A=C^*PC$ since $$A^2=C^*P(CC^*)PC \leq C^*P^*P^{-1}PC=A.$$ Let $\lambda_1,\cdots,\lambda_n \geq 0$ be the eigenvalues of $A$ in non-increasing order and $\theta_1,\cdots,\theta_n$ the corresponding orthonormal eigenvectors. Take $k$ to be the largest integer such that $\lambda_1=\cdots=\lambda_k.$ Then, $$\sum_{i=1}^m \frac{p_i}{\lambda_1}C_i^*C_i+\sum_{i=k+1}^n \Bigl(1-\frac{\lambda_i}{\lambda_1} \Bigl)\theta_i\theta_i^*=\frac{A}{\lambda_1}+\sum_{i=k+1}^n \Bigl(1-\frac{\lambda_i}{\lambda_1} \Bigl)\theta_i\theta_i^*=I_n$$ as $A=\sum_{i=1}^n \lambda_i \theta_i \theta_i^*.$ If $\lambda_1<1$, one may write the previous equation as $$\sum_{i=1}^m p_iC_i^*C_i+\sum_{i=k+1}^n \Bigl(1-\frac{\lambda_i}{\lambda_1}\Bigl)\theta_i\theta_i^*+\sum_{i=1}^m p_i\Bigl(\frac{1}{\lambda_1}-1\Bigl)C_i^*C_i=I_n$$ with the coefficient terms being always positive as $1 \geq \lambda_1 \geq \lambda_i$ for any $1 \leq i \leq n.$  To conclude, there exists some $\nu \in \mathbb{N} \cup \{0\}$ such that there are $k_j \times n$ matrices $B_j$ with $B_jB_j^* \leq I_{k_j}$ or $B_jB_j^* \geq I_{k_j}$ (either $B_j=\theta_j^*$ then $B_jB_j^*=|\theta_j|^2=1$ or $B_j=C_{i_j}$) and $b_j>0$ for $1 \leq j \leq \nu$ satisfying $$\sum_{i=1}^m p_iC_i^*C_i+\sum_{j=1}^\nu b_jB_j^*B_j=I_n.$$ Taking $f_i$ given functions which are more log-convex than $g_{\frac{Q_i-I_{n_i}}{2\pi}}$ for $1 \leq i \leq m$, we set $g_1=f_1,\cdots,g_m=f_m$ and $g_{m+1}=\cdots=g_{m+\nu}=1$,  $Q'_j=I_{k_j}$ for $1 \leq j \leq \nu$ and $I'=\{m+j: j \in \{1,\cdots,\nu\} \text{ s.t. }  B_jB_j^* \le I_{k_j}\}.$ Since the conditions to apply the previous theorem are satisfied with $$((C_i)_{1 \leq i \leq m},(B_j)_{1 \leq j \leq \nu}), ((Q_i)_{1 \leq i \leq m},(Q_j')_{1 \leq j \leq \nu},I \cup I')$$ and $\int g_{i}\, d\gamma_{k_i}=1$ for all $m+1 \leq i \leq m+\nu$, the desired claim is proven. 
\end{proof}

\begin{theorem}
    Let $m,n_1,\cdots,n_m$ positive integers, $N=n_1+\cdots+n_m.$ Suppose $X_i$ is an $n_i$-dimensional random vector for $1 \leq i \leq m$ such that their joint law $$X=(X_1,\cdots,X_m)$$ forms a centered jointly $N$-dimensional Gaussian random vector with covariance matrix $T_{ii} \leq I_{n_i}, (I_{n_i}-T_{ii})(Q_i-I_{n_i})=0$ for matrices $Q_i \geq I_{n_i}$ for any $i \in I$ and $T_{ii} \geq I_{n_i}, (I_{n_i}-T_{ii})(Q_i-I_{n_i})=0$ for matrices $Q_i \leq I_{n_i}$ for any $i \in I^c.$ Let $P$ be the block diagonal matrix $$P=\mathrm{diag}(p_1I_{n_1},\cdots,p_mI_{n_m}).$$ Then for functions $f_i$ that are more log-convex (resp. more log-concave) than $g_{\frac{Q_i-I_{n_i}}{2\pi}}$ for any $i \in I$ (resp. for any $i \in I^c$), if $T \leq P^{-1}$, $$\mathbb{E}\prod_{i=1}^m f_i(X_i)^{p_i} \leq \prod_{i=1}^m \Bigl(\int_{\mathbb{R}^{n_i}} f_i \, d\gamma_{n_i} \Bigl)^{p_i}.$$ 
\end{theorem}

\begin{proof}
    Set $n=\mathrm{rk}(T).$ According to \cite[Lemma 2]{Chen_Dafnis_Paouris}, there exists a $N \times n$ matrix $U$ such that $T=UU^*.$ We write $U_1,\cdots,U_m$ for the block rows of $U$ and denote by $u_1^i,\cdots,u_{n_i}^i$ the rows of $U_i.$ Then, $T=UU^*=(U_iU_j^*)_{1 \leq i,j \leq m}$ which implies that $T_{ii}=U_iU_i^*$ for any $1 \leq i \leq m.$ Moreover, $$X=(X_1,\cdots,X_m)^T\overset{d}{=}(U_1Z,\cdots,U_mZ)^T=UZ$$ where $Z$ is an $n$-dimensional standard Gaussian random vector. Then, we can make use of the previous theorem since $$\mathbb{E}\prod_{i=1}^m f_i(X_i)^{p_i}=\int_{\mathbb{R}^n}\prod_{i=1}^m f_i(U_ix)^{p_i}\, d\gamma_n(x)$$ and $U_i$ satisfies the conditions of \cref{Gaussian_Version_BL}. 
\end{proof}


Thanks to our results, one may improve the forward regularized hypercontractivity inequalities of \cite[Corollary 6.6]{Nakamura_Tsuji}. Set $$\mathcal{Q}=\frac{1}{2\pi(1-e^{-2s})}\begin{pmatrix}
    \Bigl( 1-\frac{1-e^{-2s}}{p}\Bigl) \mathrm{id}_{\mathbb{R}^n} & -e^{-s}\mathrm{id}_{\mathbb{R}^n} \\
    \nm 
    -e^{-s}\mathrm{id}_{\mathbb{R}^n} & \Bigl( 1-\frac{1-e^{-2s}}{q'}\Bigl)\mathrm{id}_{\mathbb{R}^n}
\end{pmatrix}>0, p_1=\frac{1}{p}, p_2=\frac{1}{q'}$$ such that $$\frac{q-1}{p-1} \le e^{2s}, 1<p,q<+\infty$$ and let $$C_1=\begin{pmatrix}\mathrm{id}_{\mathbb{R}^n} & 0\end{pmatrix}, C_2=\begin{pmatrix} 0 & \mathrm{id}_{\mathbb{R}^n}\end{pmatrix}.$$  

The hypercontractivity inequality on $\mathbb{R}^n$ can be considered as a special case of the Brascamp-Lieb inequality $$\int_{\mathbb{R}^n}e^{-\pi \langle \mathcal{Q}x,x\rangle}\prod_{j=1}^2 f_j(C_jx)^{p_j}\ dx \leq \mathrm{H}(p_1,p_2)\prod_{j=1}^2 \Bigl( \int_{\mathbb{R}^n} f_j\Bigl)^{p_j}.$$



For simplicity, let us consider $n=1.$ 
Then, $$\begin{array}{lll}M &=& p_1aC_1^*C_1+p_2bC_2^*C_2+\mathcal{Q}=\begin{pmatrix}
    \frac{1}{p}a & 0 \\
    \nm 
    0 & \frac{1}{q'}b
\end{pmatrix}+\mathcal{Q} \\
\nm 
&=& \ds \frac{1}{2\pi(1-e^{-2s})}\begin{pmatrix}
    1+\frac{1-e^{-2s}}{p}(2\pi a-1) & e^{-s} \\
    \nm 
    e^{-s} & 1+\frac{1-e^{-2s}}{q'}(2\pi b-1)
\end{pmatrix}.
\end{array}$$
In particular, $$M^{-1}=\frac{2\pi}{ 1+\frac{2\pi a-1}{p}+\frac{2\pi b-1}{q'}+\frac{1-e^{-2s}}{pq'}(2\pi a-1)(2\pi b-1)}\begin{pmatrix}
    1+\frac{1-e^{-2s}}{q'}(2\pi b-1) & e^{-s} \\
    \nm 
    e^{-s} & 1+\frac{1-e^{-2s}}{p}(2\pi a-1)
\end{pmatrix}.$$

By \cref{Prop_Equiv_BL_constant},  \begin{equation} \label{Sup_Det_Hypercontractivity} \sup_{\substack{0<a_1 \leq \frac{1}{2\pi \beta} \\ 0<a_2 \leq \frac{1}{2\pi \alpha}}}\frac{a_1^{p_1/2}a_2^{p_2/2}}{\det(p_1a_1C_1^*C_1+p_2a_2C_2^*C_2+\mathcal{Q})^{1/2}}=\frac{a^{p_1/2}b^{p_2/2}}{\det(M)^{1/2}}\end{equation} if and only if $$2\pi \beta \leq \frac{1}{a}= C_1M^{-1}C_1^*=(M^{-1})_{11} \text{ or } \frac{1}{a}=2\pi \beta \geq (M^{-1})_{11} $$ and $$2\pi \alpha \leq \frac{1}{b} =C_2M^{-1}C_2^*=(M^{-1})_{22} \text{ or } \frac{1}{b}=2\pi \alpha \geq (M^{-1})_{22}.$$ Notice that if the supremum is taken over $0<a_1 \leq \frac{1}{2\pi \beta}$ and $0<a_2$, then the condition for $a_2$ restricts to $\frac{1}{b}=(M^{-1})_{22}.$ \hfill \\

Let us rewrite the equations $\frac{1}{a}=(M^{-1})_{11}, \frac{1}{b}=(M^{-1})_{22}.$ Indeed, we have $$\begin{array}{lll}
\ds \frac{1}{b}=\frac{2\pi}{1+\frac{2\pi a-1}{p}+\frac{2\pi b-1}{q'}+\frac{1-e^{-2s}}{pq'}(2\pi a-1)(2\pi b-1)}\Bigl( 1+\frac{1-e^{-2s}}{p}(2\pi a-1)\Bigl) \\
\nm 
\Longrightarrow \ds 1-\frac{1}{q'}+\frac{2\pi b}{q'}+\frac{2\pi a-1}{p}+\frac{1-e^{-2s}}{pq'}(2\pi a-1)(2\pi b-1)=2\pi b+\frac{1-e^{-2s}}{p}2\pi b(2\pi a-1) \\
\nm 
\Longrightarrow \ds -\bigl(1-\frac{1}{q'}\bigl)(2\pi b-1)+\frac{1-e^{-2s}}{p}(2\pi a-1)\Bigl( -\bigl(1-\frac{1}{q'}\bigl)(2\pi b-1)-1\Bigl) +\frac{2\pi a-1}{p}=0 \\
\nm 
\Longrightarrow \ds \Bigl( \frac{1-e^{-2s}}{p}(2\pi a-1)+1 \Bigl) \frac{1}{q}(2\pi b-1)=\frac{e^{-2s}}{p}(2\pi a-1) \\
\nm 
\Longrightarrow \ds b=\frac{1}{2\pi}\Bigl(1 +q\frac{\frac{e^{-2s}}{p}(2\pi a-1)}{\frac{1-e^{-2s}}{p}(2\pi a-1)+1}\Bigl)
\end{array}$$ and analogously, $$a=\frac{1}{2\pi} \Bigl( 1+p'\frac{\frac{e^{-2s}}{q}(2\pi b-1)}{\frac{1-e^{-2s}}{q}(2\pi b-1)+1}\Bigl).$$



\begin{theorem}
Assume that $\beta>1.$ If $$\alpha>\frac{1+\frac{1-\beta}{p\beta}(1-e^{-2s})}{1+\frac{1-\beta}{p\beta}(1+(q-1)e^{-2s})},$$ then the following inequality holds for $f_1$ that is more log-convex than $g_{1/(2\pi\beta)}, \text{ formally } (\log f_1)'' \geq -\frac{1}{\beta}$ and $f_2$ that is more log-convex than $g_{1/(2\pi\alpha)}, \text{ formally } (\log f_2)'' \geq -\frac{1}{\alpha}$: $$\frac{\ds \int_{\mathbb{R}^2}e^{-\pi \langle \mathcal{Q}x,x\rangle}\prod_{j=1}^2 f_j(C_jx)^{p_j}\, dx}{\ds \prod_{j=1,2}\Bigl(\int_{\mathbb{R}} f_j\Bigl)^{p_j}} \leq (2\pi)^{1-\frac{1}{2}(p_1+p_2)}(1-e^{-2s})^{1/2}\frac{\beta^{-1/2p}\alpha^{-1/2q'}}{\Bigl(1+\frac{1}{p}\frac{1-\beta}{\beta}+\frac{1}{q'}\frac{1-\alpha}{\alpha}+\frac{1-e^{-2s}}{pq'}\frac{1-\beta}{\beta}\frac{1-\alpha}{\alpha}\Bigl)^{1/2}}.$$ 
\end{theorem}

\begin{remark}
    If $\alpha$ is smaller than this exponent, one cannot expect to improve the inequality and \cite[Corollary 6.6]{Nakamura_Tsuji} gives already the best possible bound. This is coherent with the fact that at $\alpha=\frac{1+\frac{1-\beta}{p\beta}(1-e^{-2s})}{1+\frac{1-\beta}{p\beta}(1+(q-1)e^{-2s})}$, the constant in the previous theorem is $$\begin{array}{lll}
    \ds (2\pi)^{1-\frac12 (p_1+p_2)}(1-e^{-2s})^{1/2}\beta^{-1/2p}\Bigl(1+\frac{1-\beta}{p\beta}(1-e^{-2s})\Bigl)^{-1/2q'}\Bigl( 1+\frac{1-p}{p\beta}(1+(q-1)e^{-2s}) \Bigl)^{-1/2q} \\
    \ds = (2\pi)^{1-\frac12 (p_1+p_2)}(1-e^{-2s})^{1/2}\Bigl\|P_s[(\frac{\gamma_\beta}{\gamma})^{\frac 1p}]\Bigl\|_{L^q(\gamma)}
    \end{array}$$ which is precisely the constant of Corollary 6.6.(i) in \cite{Nakamura_Tsuji}. 
\end{remark}

\begin{proof}
    From what was proven before, it suffices to show that for $\alpha, \beta$ as in the assumptions of the theorem, the map $$\begin{array}{lll}\ds F: \Bigl[0;\frac{1}{2\pi \beta}\Bigl] \times  \Bigl[ 0; \frac{1}{2\pi \alpha}\Bigl] \rightarrow \mathbb{R}^2, \\ \ds F(x,y)=\Bigl(\frac{1}{2\pi} \Bigl( 1+p'\frac{\frac{e^{-2s}}{q}(2\pi y-1)}{\frac{1-e^{-2s}}{q}(2\pi y-1)+1}\Bigl) ,\frac{1}{2\pi}\Bigl(1 +q\frac{\frac{e^{-2s}}{p}(2\pi x-1)}{\frac{1-e^{-2s}}{p}(2\pi x-1)+1}\Bigl)\Bigl) \end{array}$$ has no fixed points. The previous claim, whose proof follows by straightforward computations, proves that the supremum in \cref{Sup_Det_Hypercontractivity} must be attained at $$a=\frac{1}{2\pi \beta}, b=\frac{1}{2\pi \alpha}$$ and one may conclude by applying \cref{Theorem_REG_BL} with the data defined above and some rearranging. 
\end{proof}






\bigskip

\section*{Acknowledgements}

The author is grateful to his supervisor Alessio Figalli, for his support and useful discussions during the elaboration of this work. He also would like to thank Shohei Nakamura for many helpful comments. Finally, he wishes to thank Guido de Philippis and Yair Shenfeld for a conversation about \cref{De_Phil_Shenfeld_contraction}, and for drawing the author's attention to reference \cite{Gozlan}. 

\hfill \\



\appendix 

\section{A general Caffarelli's Contraction Theorem}

\hfill \\

A celebrated result of Caffarelli in \cite{Caff} concerns the Lipschitz property of the optimal transportation map under log-convexity (resp. log-concavity) assumptions on the source measure (resp. the target measure). The following statement was shown in \cite{Caff} and subsequently extended in \cite{Kolesnikov}:  if $d\mu=e^{-V}\, dx$ and $d\nu=e^{-W}\, dx$ are probability measures in $\mathbb{R}^n$, with $\mu$ supported in all of $\mathbb{R}^n$, such that there exist $\lambda,\Lambda>0$ with $$\nabla^2 V \leq \lambda I_n \text{ and } \nabla^2 W \geq \Lambda I_n$$ and $\nabla \Phi$ is the Brenier map transporting $\mu$ to $\nu$, then $$\nabla^2\Phi \leq \sqrt{\frac{\lambda}{\Lambda}}I_n.$$
For our purposes, we would like to prove a similar statement with slightly more general assumptions, replacing $\lambda I_n, \Lambda I_n$ by positive definite matrices. 


\begin{theorem} \label{General_Caff_Contraction}
    Let $d\mu=e^{-V}\, dx$ and $d\nu=e^{-W}\, dx$ be probability measures on $\mathbb{R}^n$, with $\mu$ supported on all of $\mathbb{R}^n$, such that there exist $\mathbb{R}^{n \times n} \ni A,B>0$ with \begin{equation} \label{Caff_condition}\nabla^2 V \leq A \text{ and } \nabla^2 W \geq B.\end{equation} Let $\nabla \Phi$ be the Brenier map transporting $\mu$ to $\nu.$ Then, $$\nabla^2\Phi \leq A^{1/2}(A^{1/2}BA^{1/2})^{-1/2}A^{1/2}.$$
\end{theorem}

\hfill \\

This result can actually be deduced from the work of Chewi and Pooladian (see \cite{Chewi_Pooladian}) who proved the theorem for two commuting matrices $A$ and $B$ under mild regularity assumptions on $V$ and $W.$ It was also proved directly by Gozlan and Sylvestre in \cite{Gozlan}. 
However, the proof we would like to present does not use any advanced tools such as the entropic regularization, which the aforementioned authors brought into play to deduce an extension of Caffarelli's theorem.

\begin{remark}
Recall that the Hessian of the Brenier potential between the Gaussian distributions $N(0,A^{-1})$ and $N(0,B^{-1})$, which trivially satisfy the constraint \cref{Caff_condition}, is the matrix $$B^{-1/2}(B^{-1/2}A^{-1}B^{-1/2})^{-1/2}B^{-1/2}=B^{-1/2}(B^{1/2}AB^{1/2})^{1/2}B^{-1/2}$$ by \cite[Lemma 2.3]{Takatsu}. But setting $C=B^{1/2}A^{1/2}(A^{1/2}BA^{1/2})^{-1/2}A^{1/2}B^{1/2}$, one has $$\begin{array}{lll}
C^2 &=& B^{1/2}A^{1/2}(A^{1/2}BA^{1/2})^{-1/2}A^{1/2}BA^{1/2}(A^{1/2}BA^{1/2})^{-1/2}A^{1/2}B^{1/2} \\
\nm 
&=& B^{1/2}A^{1/2}A^{1/2}B^{1/2}=B^{1/2}AB^{1/2}.
\end{array}$$ This proves that $$A^{1/2}(A^{1/2}BA^{1/2})^{-1/2}A^{1/2}=B^{-1/2}(B^{1/2}AB^{1/2})^{1/2}B^{-1/2}.$$

This observation shows that the previous theorem is sharp for every pair of matrices and implies that the Brenier potential between Gaussians achieves the largest possible Hessian among all source and target measures obeying the constraint \cref{Caff_condition}. 
\end{remark}

\hfill \\

The proof strategy is based on the maximum principle and is very similar to \cite{Vald1}. 

\begin{proof}
    Set $$\tilde V=V(A^{-1/2}\cdot)+\frac{1}{2}\log(\det A), \tilde W=W(A^{1/2}\cdot)-\frac{1}{2}\log(\det A)$$ and $\tilde \mu=e^{-\tilde V}\, dx=(A^{1/2})_{\#}\mu, \tilde \nu=e^{-\tilde W}\, dx=(A^{-1/2})_{\#}\nu$.  Then, defining $\tilde \Phi(x)=\Phi(A^{-1/2}x)$ for all $x \in \mathbb{R}^n$ with $\nabla \Phi$ being the Brenier map that transports $\mu$ to $\nu$, $\nabla \tilde \Phi$ is the Brenier map that transports $\tilde \mu$ to $\tilde \nu$. Moreover, $$\nabla^2 \tilde V= A^{-1/2}\nabla^2 V(A^{-1/2}\cdot)A^{-1/2} \leq I_n, \nabla^2 \tilde W \geq A^{1/2}BA^{1/2}.$$ Since $$\nabla^2 \Phi=A^{1/2}\nabla \tilde \Phi(A^{1/2}\cdot)A^{1/2},$$ we see that it suffices to show that $$\nabla^2 \tilde \Phi \leq (A^{1/2}BA^{1/2})^{-1/2}.$$ In other words, we reduced the claim to showing that: $$\text{ If } D^2V \leq I_n \text{ and } D^2W \geq A>0, \text{ then } \nabla^2 \Phi \leq A^{-1/2}.$$ 
    \hfill \\ First of all, we replace $e^{-W}$ by $e^{-W_r}$ with $W_r: \overline{B}_r(0) \rightarrow \mathbb{R}$ given by $W_r(x)=W(x)+C'$ where $C'$ is a normalizing constant chosen so that $\int e^{-W_r}\, dx=1.$ Let us write the Brenier map transporting $e^{-V}\, dx$ to $e^{-W_r}\, dx$  as $\nabla \Phi^r$ with $\Phi^r$ convex.
    \hfill \\ \\ We will be interested in the function $$K(x,\alpha)=\langle A^{-1/2}\alpha, \alpha \rangle-\frac{\Phi^r(x+h\alpha)+\Phi^r(x-h\alpha)-2\Phi^r(x)}{h^2} \quad \forall x \in \mathbb{R}^n, \forall \alpha \in S^{n-1}$$ for some fixed $h>0.$ One may assume by approximation that $V,W \in C^{1,1}_{\mathrm{loc}}(\mathbb{R}^n).$ Hence, by Caffarelli's regularity theory, $K$ is $C^2$ on $\mathbb{R}^n \times S^{n-1}$ and the goal is to show that $K$ is non-negative. 
    By the convexity of $\Phi^r$, $$K(x,\alpha) \leq \langle A^{-1/2}\alpha,\alpha\rangle \quad \forall \alpha \in S^{n-1}.$$ As in the limit as $x$ tends to infinity, this inequality becomes an equality (by \cite[Lemma 2.1]{Vald1} or by \cite[Lemma 3.1]{Alessio}). This implies that $K$ has a global minimum that we assume to be attained at $(x_0,
    \alpha_0)$. \\ Then, $$\partial_{x_i}K(x_0,\alpha_0)=0 \Longrightarrow \Phi_i^r(x_0+h\alpha_0)+\Phi_i^r(x_0-h\alpha_0)-2\Phi_i^r(x_0)=0 \quad \forall 1 \leq i \leq n$$ which gives $$\nabla \Phi^r(x_0+h\alpha_0)+\nabla \Phi^r(x_0-h\alpha_0)-2\nabla \Phi^r(x_0)=0.$$ In addition, for any vector $v \in \mathrm{span}\{\alpha\}^{\perp}$, $$\begin{array}{lll}\partial_{v}K(x_0,\alpha_0) &=& \ds \lim_{t \rightarrow 0}\frac{K \bigl(x_0,\frac{\alpha_0+tv}{\sqrt{1+t^2}} \bigl)-K(x_0,\alpha_0)}{t} \\ \\
    \nm 
    &=& \ds 2\langle A^{-1/2}\alpha,v\rangle-\frac{\langle \nabla \Phi^r(x_0+h\alpha_0), hv\rangle-\langle \nabla \Phi^r(x_0-h\alpha_0),hv\rangle}{h^2}=0 \end{array}$$ since $(x_0,\alpha_0)$ is a local minimum. Hence, there exists $\lambda \in \mathbb{R}$ such that $$A^{-1/2}\alpha_0=\frac{\nabla \Phi^r(x_0+h\alpha_0)-\nabla \Phi^r(x_0-h\alpha_0)}{2h}-\lambda \alpha_0$$ and so $$\nabla \Phi^r(x_0 \pm h\alpha_0)=\nabla \Phi^r(x_0) \pm h(A^{-1/2}\alpha_0+\lambda \alpha_0).$$ Taking the logarithm of the Monge-Ampère equation, one has \begin{equation}
    \label{Monge_Amp}  \begin{array}{lll}
    \log(\det \nabla^2 \Phi^r(x_0+h\alpha_0))+\log(\det \nabla^2 \Phi^r(x_0-h\alpha_0))-2\log(\det \nabla^2 \Phi^r(x_0)) \\ \\
    \nm 
    
    = W_r(\nabla \Phi^r(x_0+h\alpha_0))+W_r(\Phi^r(x_0-h\alpha_0))-2W_r(\Phi^r(x_0)) \\ \\
    \nm 
     -[V(x_0+h\alpha_0)+V(x_0-h\alpha_0)-2V(x_0)].
    \end{array}\end{equation} Since $A \mapsto \det(A)$ is a log-concave function, the left-hand side is less than $$D(\log \det)(\nabla^2 \Phi^r(x_0))E=\mathrm{Tr}(\nabla^2 \Phi^r(x_0)^{-1}E)=\mathrm{Tr}(\nabla^2 \Phi^r(x_0)^{-1/2} \, E \; \nabla^2 \Phi^r(x_0)^{-1/2})$$ where $E=\nabla^2 \Phi^r(x_0+h\alpha_0)+\nabla^2 \Phi^r(x_0-h\alpha_0)-2\nabla^2 \Phi^r(x_0).$ As $E$ is the Hessian of the function $$x \mapsto \Phi^r(x+h\alpha_0)+\Phi^r(x-h\alpha_0)-2\Phi^r(x)$$ which attains a local maximum at $x_0$, $E$ is negative semi-definite. Hence, the left-hand side of \cref{Monge_Amp} is non-positive. \\ On the right-hand side of \cref{Monge_Amp}, we notice that $\nabla^2 V \leq I_n$ gives $$V(x_0+h\alpha_0)+V(x_0-h\alpha_0)-2V(x_0) \leq h^2|\alpha_0|^2.$$ Moreover, denoting $v=\nabla \Phi^r(x_0+h\alpha_0)-\nabla \Phi^r(x_0)=\nabla \Phi^r(x_0)-\nabla \Phi^r(x_0-h\alpha_0)$, one has $$ \begin{array}{lll} W_r(\nabla \Phi^r(x_0+h\alpha_0))+W_r(\nabla \Phi^r(x_0-h\alpha_0))-2W_r(\nabla \Phi^r(x_0)) \\ \\
    \nm 
    = \ds W_r(\nabla \Phi_r(x_0)+v)-W_r(\nabla \Phi_r(x_0))+W_r(\nabla \Phi_r(x_0)-v)-W_r(\nabla \Phi^r(x_0)) \\ \\
    \nm 
    =\ds \int_0^1 \langle v, \nabla W_r(\nabla \Phi^r(x_0)+sv)-\nabla W_r(\nabla \Phi^r(x_0)-sv)\rangle \, ds \\ \\
    \nm 
    = \ds \int_0^1 \int_{-s}^s \langle v, \nabla^2 W_r(\nabla \Phi^r(x_0)+tv)v\rangle \, dt \, ds.
    \end{array}$$ 
 As $\nabla^2 W_r \geq A$, we conclude that
    $$\begin{array}{lll}W_r(\nabla \Phi^r(x_0+h\alpha_0))+W_r(\nabla \Phi^r(x_0+h\alpha_0))-2W_r(\nabla \Phi^r(x_0)) \\
    \nm 
    \geq h^2 \langle A(A^{-1/2}\alpha_0+\lambda \alpha_0),A^{-1/2}\alpha+\lambda \alpha_0\rangle.
    \end{array}$$ since $v=h(A^{-1/2}\alpha_0+\lambda \alpha_0).$ Hence, $$2\lambda \langle A^{1/2}\alpha_0,\alpha_0\rangle+\lambda^2 \langle A\alpha_0,\alpha_0\rangle=-\langle \alpha_0,\alpha_0\rangle+\langle A(A^{-1/2}\alpha_0+\lambda \alpha_0),A^{-1/2}\alpha_0+\lambda \alpha_0\rangle \leq 0$$
    and so $\lambda \leq 0.$ Thus, $$\langle A^{-1/2}\alpha_0,\alpha_0\rangle \geq \frac{\langle \nabla \Phi^r(x_0+h\alpha_0),\alpha_0\rangle-\langle \nabla \Phi^r(x_0-h\alpha_0),\alpha_0\rangle}{2h}.$$ This implies by the convexity of $\Phi^r$ that $$\begin{array}{lll}
    \ds \frac{\Phi^r(x_0+h\alpha_0)+\Phi^r(x_0-h\alpha_0)-2\Phi^r(x_0)}{h^2} &\leq& \ds \frac{\langle \nabla \Phi^r(x_0+h\alpha_0), \alpha_0\rangle-\langle \nabla \Phi^r(x_0-h\alpha_0),\alpha_0\rangle}{h} \\
    \nm 
    &\leq& 2\langle A^{-1/2}\alpha_0,\alpha_0\rangle\end{array}$$ which allows us to conclude that $$K(x,\alpha) \geq K(x_0,\alpha_0) \geq -\langle A^{-1/2}\alpha_0,\alpha_0\rangle \quad \forall x \in \mathbb{R}^n, \alpha \in S^{n-1}$$ since $(x_0,\alpha_0)$ is a global minimum. This is equivalent to $$\frac{\Phi^r(x+h\alpha)+\Phi^r(x-h\alpha)-2\Phi^r(x)}{h^2} \leq \langle A^{-1/2}\alpha, \alpha \rangle+\langle A^{-1/2}\alpha_0,\alpha_0\rangle.$$ Taking $M$ to be the ratio of the largest and smallest eigenvalues of $A^{-1/2}$, one obtains $$\langle A^{-1/2}\alpha_0,\alpha_0\rangle \leq M\langle A^{-1/2}\alpha,\alpha\rangle$$ by Rayleigh's principle, \textit{i.e.} for $\mathbb{R}^n \ni B>0$, $$\langle Bx,x\rangle \geq \lambda_{\min}(B) \text{ and } \langle Bx,x\rangle \leq \lambda_{\max}(B) \quad \forall x \in S^{n-1}.$$ Therefore, $$\frac{\Phi^r(x_0+h\alpha)+\Phi^r(x_0-h\alpha)-2\Phi^r(x)}{h^2} \leq (M+1)\langle A^{-1/2}\alpha,\alpha\rangle.$$ 
    %Using the fact that the estimate is uniform in $h$, we apply the same strategy of bootstrapping as in \cite{Kolesnikov,Vald1} to get that $$\frac{\Phi^r(x_0+h\alpha)+\Phi^r(x_0-h\alpha)-2\Phi^r(x)}{h^2} \leq \langle A^{-1/2}\alpha,\alpha\rangle$$ which asserts that $K(x_0,\alpha_0) \geq 0$

    
    We now may apply the same strategy of bootstrapping as in \cite{Kolesnikov,Vald1} to get the desired estimate; more precisely, assuming that we have an estimate of the type $$\frac{\Phi^r(x_0+h\alpha)+\Phi^r(x_0-h\alpha)-2\Phi^r(x)}{h^2} \le a \langle A^{-1/2}\alpha,\alpha\rangle,$$ one gets letting $h \rightarrow 0$ that $$\langle \nabla^2 \Phi^r(x)\alpha, \alpha \rangle \le a\langle A^{-1/2}\alpha,\alpha \rangle.$$ Using this additional information, one has $$\langle \nabla \Phi^r(x_0+s\alpha_0)-\nabla \Phi^r(x_0-s\alpha),\alpha_0\rangle=\int_{-s}^s \langle D^2\Phi^r(x_0+t\alpha_0)\alpha_0,\alpha_0\rangle \, dt \le 2\langle A^{-1/2}\alpha_0,\alpha_0\rangle as$$ and recall that we have already shown the following bound: $$\begin{array}{lll}
    \ds \langle \nabla \Phi^r(x_0+s\alpha_0)-\nabla \Phi^r(x_0-s\alpha_0),\alpha_0\rangle &\le& \ds \langle \nabla \Phi^r(x_0+h\alpha_0)-\nabla \Phi^r(x_0-h\alpha_0),\alpha_0\rangle \\ \\
     &\le& \ds 2h\langle A^{-1/2}\alpha_0,\alpha_0\rangle
    \end{array}$$ for any $0 \le s \le h.$ 
    
    This implies $$\begin{array}{lll}
    \ds \frac{\Phi^r(x_0+h\alpha_0)+\Phi^r(x_0-h\alpha_0)-2\Phi^r(x)}{h^2} &\le& \ds \frac1{h^2}\int_0^h \langle \nabla \Phi^r(x_0+s\alpha_0)-\nabla \Phi^r(x_0-s\alpha_0),\alpha_0\rangle \, ds \\ \\
    &\le& \ds \frac1{h^2}\int_0^h 2\langle A^{-1/2}\alpha_0,\alpha_0\rangle \min(a s,h)\, ds \\ \\
    &=& \ds \frac{2a-1}{a}\langle A^{-1/2}\alpha_0,\alpha_0\rangle.
    \end{array}$$
    As $a_1=M+1, a_{n+1}=\frac{2a_n-1}{a_n}$ defines a decreasing sequence converging to $1$, we get by letting $n \rightarrow +\infty$ the following: $$\frac{\Phi^r(x_0+h\alpha_0)+\Phi^r(x_0-h\alpha_0)-2\Phi^r(x_0)}{h^2} \le \langle A^{-1/2}\alpha_0,\alpha_0\rangle.$$
    In other words, this proves that $$K(x,\alpha) \ge K(x_0,\alpha_0) \ge  0 \quad \forall x \in \mathbb{R}^n \text{ and } \alpha \in S^{n-1}.$$ Letting $h\rightarrow 0$, this shows that $$\langle \nabla^2 \Phi^r(x)\alpha,\alpha\rangle \leq \langle A^{-1/2}\alpha,\alpha\rangle$$ which translates to $$\frac{\langle \nabla \Phi^r(x+h\alpha)-\nabla \Phi^r(x-h\alpha),\alpha\rangle}{2h} \leq \langle A^{-1/2}\alpha,\alpha \rangle$$ by integration. Since $e^{-W_r}\, dx$ tends to $e^{-W}\, dx$ weakly, the stability of the transport map gives that $$\frac{\langle \nabla \Phi(x+h\alpha)-\nabla \Phi(x-h\alpha),\alpha \rangle}{2h} \leq \langle A^{-1/2}\alpha,\alpha \rangle.$$ Taking the limit as $h \rightarrow 0$, we get $$\langle \nabla^2 \Phi(x)\alpha,\alpha\rangle \leq \langle A^{-1/2}\alpha,\alpha \rangle \quad \forall \alpha \in S^{n-1}$$ which concludes the proof.
\end{proof}

\bigskip

Furthermore, we would like to add that it is possible to retrieve by PDE arguments a generalization of \cite{Guido}, which was proven in \cite{Gozlan}. Let us first give an anisotropic extension of \cite[Thm 1.4]{Guido}:

\begin{theorem} \label{De_Phil_Shenfeld_contraction}
    Let $d\mu=e^{-V}\, dx$ and $d\nu=e^{-W}\, dx$ be probability measures on $\mathbb{R}^n$ with $\mu$ supported on all of $\mathbb{R}^n$, such that there exist $\alpha>0$ and $A,G \in \mathbb{R}^{n \times n} \text{ satisfying } A \geq 0 , G>0,$ $$\mathrm{div}(A \nabla V) \leq \alpha \text{ and } \nabla^2 W \geq G^{-1}.$$ Let $\nabla \Phi : \mathbb{R}^n \rightarrow \mathbb{R}^n$ be the Brenier map transporting $\mu$ to $\nu.$ Then, $$\| \mathrm{div}(A\nabla \Phi)\|_{L^\infty} \leq \sqrt{\alpha \cdot \mathrm{Tr}(AG)}.$$ 
\end{theorem}


\hfill \\ First, let us establish the following simple trace inequality: 

\begin{lemma} \label{Simple_trace_ineq}
For $ \mathbb{R}^{n \times n} \ni B>0, \mathbb{R}^{n \times n} \ni C,D \geq 0, C \neq 0$, $$\mathrm{Tr}(B^{-1}CDC) \geq \mathrm{Tr}(DB)^{-1}\mathrm{Tr}(DC)^2.$$
\end{lemma}

\begin{proof}
By Cauchy-Schwarz, $$\begin{array}{lll}
    \mathrm{Tr}(DC)^2 &=& \ds \mathrm{Tr}((D^{1/2}CB^{-1/2})(D^{1/2}B^{1/2})^T)^2 \\
    \nm 
    &\leq& \ds \mathrm{Tr}(B^{-1/2}CDCB^{-1/2})\mathrm{Tr}(D^{1/2}BD^{1/2}) \\
    \nm 
    &\leq& \ds \mathrm{Tr}(B^{-1}CDC)\mathrm{Tr}(DB).
    \end{array}$$
\end{proof}

\hfill \\ \begin{proof} We only sketch a formal proof.
Let us first assume that $A=\mathrm{Id}_{\mathbb{R}^n}.$ 
For a given unit vector $e \in \mathbb{R}^n$ and a function $\xi: \mathbb{R}^n \rightarrow \mathbb{R}$, we denote by $\xi_e$ (resp. $\xi_{ee}$) the first (resp. second) directional derivative of $\xi$ in the direction $e.$ Taking the logarithm in the Monge-Ampère equation and differentiating twice in the direction $e$ gives $$V_{ee}=\langle \nabla^2 W(\nabla \Phi)\nabla^2 \Phi e, \nabla^2 \Phi e \rangle+\langle \nabla W, \nabla \Phi_{ee}\rangle-\mathrm{Tr}((\nabla^2 \Phi)^{-1}\nabla^2 \Phi_{ee})+\mathrm{Tr}\bigl(( (\nabla^2 \Phi)^{-1}\nabla^2 \Phi_e)^2\bigl).$$ Using $\nabla^2 W \geq G^{-1}$, the previous inequality leads to $$V_{ee} \geq \langle G^{-1}\nabla^2 \Phi e, \nabla^2 \Phi e\rangle+\langle \nabla W, \nabla \Phi_{ee}\rangle-\mathrm{Tr}((\nabla^2 \Phi)^{-1}\nabla^2 \Phi_{ee})+\mathrm{Tr}\bigl(( (\nabla^2 \Phi)^{-1}\nabla^2 \Phi_e)^2\bigl).$$ Summing on both sides over the standard basis $\{e_i\}$ of $\mathbb{R}^n$ yields $$\Delta V \geq \sum_{i=1}^n \langle \nabla^2 \Phi \, G^{-1}\, \nabla^2 \Phi e_i,e_i\rangle+\langle \nabla W, \nabla \Delta \Phi \rangle-\mathrm{Tr}((\nabla^2 \Phi)^{-1}\nabla^2 \Delta \Phi)+\sum_{i=1}^n \mathrm{Tr}\bigl(( (\nabla^2 \Phi)^{-1}\nabla^2 \Phi_{e_i})^2\bigl).$$ Assume that $\Delta \Phi$ attains its maximum at $x_0.$ On the one hand, $\nabla \Delta \Phi(x_0)=0$ and $\nabla^2 \Delta \Phi(x_0) \leq 0$ and so $$\langle \nabla W(x_0),\nabla \Delta \Phi(x_0)\rangle-\mathrm{Tr}[(\nabla^2 \Phi(x_0))^{-1}\nabla^2 \Delta \Phi(x_0)] \geq 0.$$ On the other hand, $$\mathrm{Tr}\bigl(( (\nabla^2 \Phi)^{-1}\nabla^2 \Phi_e)^2\bigl)=\mathrm{Tr}(A^2) \geq 0 \text{ with } A=(\nabla^2 \Phi)^{-1/2}\nabla^2 \Phi_e (\nabla^2 \Phi)^{-1/2}.$$ By \cref{Simple_trace_ineq}, $$\alpha \geq \Delta V(x_0) \geq \mathrm{Tr}(G^{-1}(\nabla^2 \Phi(x_0))^2) \geq \mathrm{Tr}(G)^{-1}(\Delta \Phi(x_0))^2$$ and therefore, $$\Delta \Phi(x_0) \leq \sqrt{\alpha \cdot \mathrm{Tr}(G)}$$ using that $\Delta \Phi(x_0) \geq 0$ as $\Phi$ is convex. Since $x_0$ is a point where $\Delta \Phi$ attains its maximum, we obtain the desired claim.  \\
Let us now prove the case with $A>0$. We consider $\tilde V=V(A^{1/2}\cdot)-\frac{1}{2}\log(\det A)$ and $\tilde \mu=e^{-\tilde V}\, dx.$ Notice that for any Borel set $B \subset \mathbb{R}^n$, $$(A^{-1/2})_{\#}\mu(B)=\int_{A^{1/2}B}e^{-V(x)}\, dx=\int_B e^{-V(A^{1/2}y)}\det(A)^{1/2}\, dy$$ \textit{i.e.} $(A^{-1/2})_{\#}\mu=\tilde \mu.$ Consider $\tilde \nu=(A^{1/2})_{\#}\nu=e^{-\tilde W}\, dx$ with $\tilde W=W(A^{-1/2}\cdot)+\frac{1}{2}\log(\det A).$ Set $\tilde \Phi(x)=\Phi(A^{1/2}x)$ for any $x \in \mathbb{R}^n$, where $\nabla \Phi$ is the Brenier map pushing $\mu$ forward to $\nu.$ Then, $\tilde \Phi$ is also convex and $\nabla \tilde \Phi(x)=A^{1/2} \nabla \Phi(A^{1/2}x)$ for any $x \in \mathbb{R}^n.$ This implies that $$\begin{array}{lll}
(\nabla \tilde \Phi)_{\#}\tilde \mu &=& (A^{1/2} \circ \nabla \Phi \circ A^{1/2})_{\#}\tilde \mu=(A^{1/2} \circ \nabla \Phi)_{\#}(A^{1/2})_{\#}((A^{-1/2})_{\#}\mu) \\
\nm 
&=& (A^{1/2})_{\#}((\nabla \Phi)_{\#}\mu)=(A^{1/2})_{\#}\nu \\
\nm 
&=& \tilde \nu.
\end{array}$$ Hence, by the uniqueness part of Brenier's theorem, $\nabla \tilde \Phi$ is the optimal transport map from $\tilde \mu$ to $\tilde \nu.$ Moreover, $$\alpha \geq \mathrm{div}(A\nabla V)(A^{1/2}x)=\mathrm{Tr}(A^{1/2}D^2V(A^{1/2}x)A^{1/2})=\mathrm{Tr}(D^2 \tilde V(x))=\Delta \tilde V(x) \text{ for a.e. } x \in \mathbb{R}^n$$ and similarly, $$D^2 \tilde W(x)=A^{-1/2}D^2W(A^{-1/2}x)A^{-1/2} \geq A^{-1/2}G^{-1}A^{-1/2}=(A^{1/2}GA^{1/2})^{-1} \text{ for a.e. } x \in \mathbb{R}^n.$$ Thus, $$\|\Delta \tilde \Phi\|_{L^\infty} \leq \sqrt{\alpha \mathrm{Tr}(A^{1/2}GA^{1/2})}=\sqrt{\alpha \mathrm{Tr}(AG)}.$$ The claim is proven if we notice that, as $\tilde \Phi(x)=\Phi(A^{1/2}x)$, $$\Delta \tilde \Phi(x)=\mathrm{Tr}(D^2\tilde \Phi(x))=\mathrm{Tr}(A^{1/2}D^2\Phi(A^{1/2}x)A^{1/2})=\mathrm{div}(A\nabla \Phi(A^{1/2}x))$$ which shows that $$\|\mathrm{div}(A\nabla \Phi)\|_{L^\infty} \leq \sqrt{\alpha \mathrm{Tr}(AG)}.$$  


\hfill \\ To treat the case where $A \geq 0$ with $A \neq 0$, it suffices by conjugating to show that this holds for $$A=\mathrm{diag}(\lambda_1,\cdots,\lambda_m,0,\cdots,0) \text{ for } m \leq n, \text{ where } \lambda_i>0 \quad \forall 1 \leq i \leq m.$$ As in the previous proof while considering the Laplacian, one obtains $$\begin{array}{lll} \alpha &\geq& \ds \sum_{i=1}^m \lambda_i V_{e_i e_i} \geq \sum_{i=1}^m \lambda_i \langle \nabla^2 \Phi \, G^{-1}\, \nabla^2 \Phi e_i,e_i\rangle \\ 
\nm 
&& \ds +\langle \nabla W, \nabla \sum_{i=1}^m \lambda_i \Phi_{e_ie_i} \rangle-\mathrm{Tr}((\nabla^2 \Phi)^{-1}\nabla^2 \sum_{i=1}^m \lambda_i \Phi_{e_ie_i})+\sum_{i=1}^m \lambda_i \mathrm{Tr}\bigl(( (\nabla^2 \Phi)^{-1}\nabla^2 \Phi_{e_i})^2\bigl).
\end{array}$$ Assume that $\sum_{i=1}^m \lambda_i \Phi_{e_ie_i}=\mathrm{div}(A \nabla \Phi)=\mathrm{Tr}(A \nabla^2 \Phi)$ attains its maximum at $x_0.$ As before, the contribution of the last three terms is positive and hence $$\alpha \geq \sum_{i=1}^n \langle \nabla^2 \Phi(x_0) G^{-1}\nabla^2\Phi(x_0) A^{1/2}e_i,A^{1/2}e_i\rangle=\mathrm{Tr}(A^{1/2}\nabla^2 \Phi(x_0) \, G^{-1}\nabla^2 \Phi(x_0) A^{1/2}).$$ 
By \cref{Simple_trace_ineq}, $$\begin{array}{lll}
\alpha &\geq& \mathrm{Tr}(G^{-1}\nabla^2 \Phi(x_0)A\nabla^2\Phi(x_0)) \\
\nm 
&\geq& \mathrm{Tr}(AG)^{-1}\mathrm{Tr}(A\nabla^2 \Phi(x_0))^2
\end{array}$$ and the desired claim is proven. 
\end{proof}

\hfill \\

The result of \cite[Thm 6.3]{Gozlan} for the Kantorovitch potential may be re-derived using \cref{Simple_trace_ineq}. Indeed, this lemma implies, in particular, that for any $A,B>0$, $$\mathrm{Tr}(AB^2) \geq \mathrm{Tr}(A^{-1})^{-1}\mathrm{Tr}(B)^2.$$ If $V$ is $\alpha_V$-superharmonic and $W^*$ is $\frac{1}{\beta_W}$-superharmonic, one gets $$\begin{array}{lll}
\ds n\alpha_V \geq \Delta V(x_0) &\geq& \ds \mathrm{Tr}\Bigl(\nabla^2 W(\nabla \Phi(x_0))(\nabla^2 \Phi(x_0))^2\Bigl) \\
\nm 
&\geq& \ds \mathrm{Tr}\Bigl( \bigl(\nabla^2 W(\nabla \Phi (x_0))\bigl)^{-1}\Bigl)^{-1}(\Delta \Phi(x_0))^2 \geq \frac{\beta_W}{n}(\Delta \Phi(x_0))^2
\end{array}$$ where $x_0$ is the maximum point of $\Delta \Phi.$ This implies that $$n\sqrt{\frac{\alpha_V}{\beta_W}} \geq \Delta \Phi(x_0) =\|\Delta \Phi\|_\infty.$$ 

\hfill \\ 

By allowing directional dependencies, one is able to generalize \cite[Thm 1.18]{Guido}.

\hfill \\

Indeed, consider matrices $A_1,\cdots,A_N \in \mathbb{R}^{2 \times 2}$ and set $A=\mathrm{diag}(A_1,\cdots,A_N).$
Let $\mu_{\beta,N}$ be the probability measure on $\mathbb{C}^N$ defined by $$d\mu_{\beta,N,Q}(z)=\frac{e^{-Q_{\beta,N}(z)-G_{\beta,N}(z)}}{\mathcal{Z}_{\beta,N}} \text{ with } \mathcal{Z}_{\beta,N}=\int_{\mathbb{C}^N}e^{-Q_{\beta,N}(z)-G_{\beta,N}(z)}\, dz$$ where $$Q_{\beta,N}(z)=\beta N\sum_{j=1}^N Q_j(z_j) \text{ with } Q_j: \mathbb{C} \rightarrow \mathbb{R} \text{ for any } 1 \leq j \leq N,$$ and $$G_{\beta,N}(z)=-\beta \sum_{i<j=1}^N \log |A_i^{-1/2}z_i-A_j^{-1/2}z_j|$$ for any $z \in \mathbb{C}^N.$ Since $\mathrm{div}(A\, \nabla G_{\beta,N}) \leq 0$, one can apply the previous results to the two-dimensional (one component) Coulomb gas $\mu_{\beta,N,Q}$ and the pure background model $\nu_{\beta,N,Q}$ given by $$d\nu_{\beta,N,Q}(z)=\frac{e^{-Q_{\beta,N}(z)}}{\int e^{-Q_{\beta,N}(z)}\, dz}\, dz.$$ 

\begin{theorem}
    Suppose that there exists $\mathbb{R}^{2 \times 2} \ni G_i,A_i>0,$ and $\alpha_i>0$ such that $$\mathrm{div}(A_i\nabla Q_i(x)) \leq \alpha_i, \nabla^2 Q_i(x) \geq G_i^{-1} \text{ for all } 1 \leq i \leq N$$ and for a.e. $x \in \mathbb{R}^2$, where $\nabla^2 Q_i$ is the distributional derivative on $\mathbb{R}^2.$ Let $\nabla \Phi_{\beta,N,Q}$ be the optimal transport map between $\mu_{\beta,N,Q}$ and $\nu_{\beta,N,Q}.$ Then, $$\|\mathrm{div}(A\nabla \Phi_{\beta,N,Q})\|_{L^\infty(dx)} \leq \sqrt{\alpha \sum_{i=1}^N \mathrm{Tr}(A_iG_i)}$$ where $$\alpha=\sum_{i=1}^N \alpha_i.$$
\end{theorem}

\begin{proof}
    Denote $d\mu_{\beta,N,Q}=e^{-V_{\beta,N,Q}}\, dz.$ The assumption $\nabla^2 Q_i \geq G_i^{-1}$ implies that $\nabla^2 Q_{\beta,N} \geq \beta N \mathrm{diag}(G_1^{-1},\cdots,G_N^{-1}).$ Since $\mathrm{div}(A\nabla G_{\beta,N}) \leq 0$, it follows that $$\mathrm{div}(A\nabla V_{\beta,N,Q})=\mathrm{div}(A\nabla Q_{\beta,N})+\mathrm{div}(A\nabla G_{\beta,N}) \leq \beta N \sum_{i=1}^N\mathrm{div}(A_i\nabla Q_i) \leq \beta N\alpha$$ where $\alpha=\sum_{i=1}^N \alpha_i.$ Hence, we can apply \cref{De_Phil_Shenfeld_contraction} with $$n=2N, \mu=\mu_{\beta,N,Q}, \nu=\nu_{\beta,N,Q}, A=\mathrm{diag}(A_1,\cdots,A_N), G=(\beta N)^{-1}\mathrm{diag}(G_1,\cdots,G_N).$$ This shows that $$\begin{array}{lll}
    \ds \|\mathrm{div}(A\nabla \Phi_{\beta,N,Q})\|_{L^\infty(dx)} &\leq& \ds \sqrt{\beta N\alpha\, \mathrm{Tr}(AG)} \\
    \nm 
    &\leq& \ds \sqrt{\alpha \sum_{i=1}^N \mathrm{Tr}(A_iG_i)}.
    \end{array}$$
\end{proof}


\bigskip

\printbibliography

\end{document}
